% ============================================================================
%  Zero-shot object-goal navigation by LLM-guided room search
%  Single self-contained file: preamble, body, appendix and bibliography.
%  Compile with:  pdflatex main && pdflatex main   (no BibTeX run needed)
%
%  Placeholders: \TODO{...} prints a red note; \tbd prints a red "TODO" in a
%  table cell. Search the source for "TODO" to find every open item.
% ============================================================================
\documentclass[10pt]{article}

\usepackage[utf8]{inputenc}
\usepackage[T1]{fontenc}
\usepackage{lmodern}
\usepackage{microtype}
\usepackage[margin=1in]{geometry}
\usepackage{amsmath,amssymb,mathtools}
\usepackage{graphicx}
\usepackage{booktabs}
\usepackage{multirow}
\usepackage{array}
\usepackage{xcolor}
\usepackage{xspace}
\usepackage{enumitem}
\usepackage{caption}
\usepackage{subcaption}
\usepackage[numbers,sort&compress]{natbib}
\usepackage[ruled,vlined,linesnumbered]{algorithm2e}
\usepackage[colorlinks=true,linkcolor=blue!60!black,citecolor=green!40!black,urlcolor=blue!60!black]{hyperref}
\usepackage[capitalise,nameinlink]{cleveref}

% --- Placeholders -------------------------------------------------------------
\newcommand{\TODO}[1]{\textcolor{red}{\textbf{[TODO:} #1\textbf{]}}}
\newcommand{\tbd}{\textcolor{red}{TODO}}

% --- Method name ----------------------------------------------------------------
% TODO: choose the final method name. It is NOT "ZSON": that is the name of
% Majumdar et al. (NeurIPS 2022), which this paper cites and compares against.
% "LGOS" (LLM-Guided Object Search) is the working name taken from the caption
% of the method's own pseudocode.
\newcommand{\ours}{LGOS\xspace}

% --- Notation -------------------------------------------------------------------
\newcommand{\G}{\mathcal{G}}
\newcommand{\M}{\mathcal{M}}
\newcommand{\Ob}{\mathcal{O}}
\newcommand{\D}{\mathcal{D}}
\newcommand{\R}{\mathcal{R}}
\newcommand{\N}{\mathcal{N}}
\newcommand{\cstar}{c^{*}}
\newcommand{\SR}{\textsc{SR}\xspace}
\newcommand{\SPL}{\textsc{SPL}\xspace}
\newcommand{\SoftSPL}{\textsc{SoftSPL}\xspace}
\newcommand{\DTG}{\textsc{DTG}\xspace}
\newcommand{\cmark}{\checkmark}
\newcommand{\xmark}{--}

% --- algorithm2e keywords (shared by every algorithm in the paper) ----------------
\SetKwInOut{Input}{Input}
\SetKwInOut{Output}{Output}
\SetKwInOut{State}{State}
\SetKwComment{Comment}{$\triangleright$\ }{}
\SetCommentSty{textit}
\SetAlFnt{\small}
\SetFuncSty{textsc}
\SetKw{KwOr}{or}
\SetKw{KwAnd}{and}
\SetKw{KwContinue}{continue}
\SetKw{Stop}{Stop}
\SetKw{Found}{Found}
\SetKw{Fail}{Fail}
\SetKwFunction{WarmUp}{WarmUp360}
\SetKwFunction{InitSceneGraph}{InitSceneGraph}
\SetKwFunction{UpdateSceneGraph}{UpdateSceneGraph}
\SetKwFunction{ClassifyRooms}{ClassifyRooms}
\SetKwFunction{CandidateNodes}{CandidateNodes}
\SetKwFunction{Dijkstra}{Dijkstra}
\SetKwFunction{LLMProbabilities}{LLMProbabilities}
\SetKwFunction{LLM}{LLM}
\SetKwFunction{RPT}{RPT\ensuremath{^*}}
\SetKwFunction{AStar}{A\ensuremath{^*}}
\SetKwFunction{Navigate}{Navigate}
\SetKwFunction{TargetClosing}{TargetClosing}
\SetKwFunction{ExploreLocal}{ExploreLocal}
\SetKwFunction{ExploreFrontier}{ExploreFrontier}
\SetKwFunction{IntegrateDepth}{IntegrateDepth}
\SetKwFunction{SegmentRooms}{SegmentRooms}

% --- Macros requested by the sections -------------------------------------
\SetKwFunction{Release}{Release}

\begin{document}

% ============================================================================
% 01_front
% ============================================================================
% ============================================================================
%  01_front.tex -- title block, abstract and keywords
% ============================================================================

\title{Zero-Shot Object-Goal Navigation via LLM Room Probabilities and Minimum-Expected-Time Ordering on an Online Scene Graph}
\author{\TODO{authors and affiliations}}
\date{}
\maketitle
\noindent\TODO{choose the final method name; \ours is a working name and must not be ZSON}

\begin{abstract}
Zero-shot object-goal navigation (ObjectNav) asks an agent to find an instance of a named object category in an unseen indoor environment from egocentric RGB-D observations and pose, without navigation-specific training. Most training-free methods score frontiers with a foundation model and pursue the best one, weighing the chance of success against travel cost one step at a time rather than over the remaining search. We present \ours, a training-free modular agent that casts each high-level search decision as minimising the expected time-to-find over an online scene graph of rooms, doorways and object landmarks. At each decision point, a large language model (LLM) that is never shown distances estimates, for each unfinished room and each opening onto unseen space not covered by one, the probability that searching it finds the target. Together with geodesic travel costs measured on the occupancy grid, these probabilities, after rule-based discounts, define a Hamiltonian path problem with probabilistic terminals, which RPT$^*$ solves exactly for the current instance. The agent executes only the first visit of the resulting order before replanning, and a confident detection hands control to a verify--approach--inspect routine that issues STOP only on confirmed, close-range evidence. We evaluate on four datasets and five validation splits, Gibson, HM3D v1 and v2, MP3D and RoboTHOR, and disclose every privileged input; \ours \TODO{summarise SR and SPL relative to the strongest training-free baselines on each benchmark}.
\end{abstract}

\noindent\textbf{Keywords:} object-goal navigation, zero-shot, large language models, scene graphs, probabilistic search.

% ============================================================================
% 02_intro
% ============================================================================
% ============================================================================
%  02_intro.tex -- Introduction
% ============================================================================
\section{Introduction}\label{sec:intro}

Finding an object in a home it has never seen is one of the most basic competences asked of a household robot.
Object-goal navigation (ObjectNav) formalises it: an agent placed in an unknown building receives the name of an
object category, such as ``toilet'', and must reach an instance of that category from egocentric RGB-D
observations within a fixed action budget, signalling arrival with an explicit STOP action in the standard
protocol~\citep{batra2020objectnav}. Photorealistic simulators, building-scale scans of real homes and simulated
apartments with physical
counterparts~\citep{savva2019habitat,xia2018gibson,chang2017matterport3d,ramakrishnan2021habitat,deitke2020robothor}
have turned the task into a standard benchmark, and many strong systems learn it: a modular semantic map with a
learned exploration policy~\citep{chaplot2020object}, potential functions predicted from top-down semantic
maps~\citep{ramakrishnan2022poni}, or end-to-end policies trained by imitation of about 80k human
demonstrations~\citep{ramrakhya2022habitat} and then fine-tuned with reinforcement
learning~\citep{ramrakhya2023pirlnav}. Such agents are bound to the scenes,
embodiment and category vocabulary of their training data, and each new setting calls for new data and new
training.

The zero-shot setting removes this dependence. The agent receives no ObjectNav training at all; it composes
pretrained foundation models for the two competences the task needs, recognising the target and knowing where to
look for it. Open-vocabulary detectors built on vision--language
pre-training~\citep{radford2021learning,cheng2024yolo} recognise categories named only at test time, and large
language models (LLMs) carry commonsense about how homes are laid out. Zero-shot agents assembled this
way~\citep{gadre2023cows,zhou2023esc,yu2023l3mvn,yokoyama2024vlfm,yin2024sg,zhang2025apexnav} are now standard
baselines on the established benchmarks, and the modular, map-based design most of them share is also the one
that a large real-world study found to transfer from simulation to real homes far better than end-to-end
policies~\citep{gervet2023navigating}.

Most of these agents, however, share one decision rule. They maintain frontiers, the boundaries between observed
free space and unknown space~\citep{yamauchi1997frontier}, or a dense value map, attach a semantic score to each
candidate, and drive to the best one. ESC turns commonsense about objects and rooms into soft constraints on the
choice of frontier~\citep{zhou2023esc}; L3MVN asks a language model to score each frontier from the objects near
it~\citep{yu2023l3mvn}; LFG uses an LLM's ``semantic guesswork'' as a frontier heuristic~\citep{shah2023navigation};
VLFM sends the agent to the frontier of highest image--text value~\citep{yokoyama2024vlfm}; SG-Nav prompts an LLM over
an online scene graph to score frontiers~\citep{yin2024sg}; and ApexNav switches between semantic and geometric
frontier selection~\citep{zhang2025apexnav}. Distance enters such rules, if at all, as a discount on the next goal,
or, in the semantic mode of ApexNav, as the length of a shortest tour over the frontiers whose score passes a
threshold, a tour that no longer weighs how promising each of them is. Neither rule expresses the trade-off that
actually determines how long a search takes. Whether to first check a nearby room of moderate promise or to head
straight for a distant room of high promise depends on what each visit costs, on how likely it is to end the search, and on what remains to be searched if it does not; these
quantities only make sense over the whole remaining search order. Always taking the most promising candidate walks
past cheap checks and returns for them later, and always taking the nearest discards commonsense altogether.

Two observations shape our approach. The first concerns the unit of decision. Frontiers are numerous, fragmented
and transient, and the semantic context of a frontier is whatever objects happen to lie near it. A person searching
an unfamiliar home reasons about rooms instead: a toilet is in a bathroom, a bed in a bedroom, and an unexplored
doorway off a hallway may lead to either. Commonsense is phrased at this level; classical semantic object search
exploited it through object--room and object--object priors~\citep{kollar2009utilizing,aydemir2013active}, and
room-segmented scene graphs~\citep{hughes2022hydra,werby2024hierarchical,honerkamp2024language} make rooms
available to a robot as first-class entities. The second observation concerns the objective. A searcher that wants
to find the target soon should minimise the expected time until it is found, the classical objective of search
theory~\citep{stone1975theory}. Its deterministic relative is the minimum-latency, or traveling-repairman, problem,
which orders visits to minimise the sum of arrival times~\citep{blum1994minimum}. When each location instead ends
the search independently with some probability, the problem becomes a Hamiltonian path problem with probabilistic
terminals, which RPT$^*$~\citep{lyu2026rpt} solves exactly by best-first search with an admissible heuristic and
dominance pruning. Recent zero-shot agents have begun to adopt objectives of this family, over frontier viewpoints
weighted by vision--language similarity~\citep{liu2026global} or over a voxel belief map that fuses LLM priors with
observations~\citep{zhou2025beliefmapnav}. We bring the objective to the level at which the commonsense lives,
rooms and the openings that lead from them into unseen space, and solve each decision's instance of it exactly.

We present \ours, a training-free ObjectNav agent that searches a home room by room in the order that minimises the
expected time to find the target. From RGB-D and pose, \ours maintains an online scene graph $\G=(\M,\Ob,\D,\R)$: a
2-D occupancy grid $\M$; open-vocabulary object landmarks $\Ob$, detected by YOLO-World~\citep{cheng2024yolo} and
confirmed in at least two frames; confirmed doorways $\D$; and rooms $\R$, segmented geometrically by a clearance
watershed over the observed free space and cut at the doorways. A small LLM, Qwen2.5-3B-Instruct~\citep{yang2024qwen25},
labels each room from the objects confirmed inside it. At every decision point \ours forms candidate nodes, namely
unfinished rooms, openings onto unseen space and sighted but still unverified target landmarks, and asks a second
LLM, Qwen2.5-14B-Instruct, for an independent probability that searching each room or opening finds the target;
sighted landmarks receive a fixed probability of 0.85 without a model call.
Fixed geometric rules then discount the probabilities of rooms the agent has already entered and left and of doorless
openings that would reveal little unseen floor.
The prompt describes the house and each node but withholds distance, so that travel is charged exactly once, by the
planner: travel costs come from Dijkstra's algorithm~\citep{dijkstra1959note} on the occupancy grid, a
service cost per node, which for a room grows with its number of live frontiers, charges each visit, and RPT$^*$ returns the visit order of minimum expected time-to-find for that instance. Only the first node is
executed, in a receding horizon. \ours follows an A$^*$~\citep{hart1968formal} route to it, performs a bounded local
visit (a walk to the room's frontiers in order of the unseen floor each would reveal per metre walked, until none
is left, or a three- or five-heading look from the threshold of an opening), and replans on the updated graph. Whenever the target is detected with sufficient confidence, a
verify--approach--inspect routine takes over and normally decides to stop only on a fresh, confident detection measured
within reach, falling back on recent evidence in a few bounded cases, then takes at most two further forward steps
toward the target before issuing STOP (\cref{sec:execution}). \Cref{fig:teaser}
illustrates one decision.

\begin{figure}[t]
  \centering
  \fbox{\parbox[c][4.5cm][c]{0.95\linewidth}{\centering\TODO{Teaser figure from one real episode. Left: the
  online scene graph seen from above, with geometric rooms shaded by the LLM's probability of finding the target,
  confirmed doorways, openings onto unseen space and confirmed object landmarks. Middle: the RPT$^*$ visit order
  with its expected time-to-find, beside the order a greedy most-probable-first rule would follow. Right: the
  executed first leg, the replanned order after that visit, and the verified STOP at the target.}}}
  \caption{\ours searches an unknown home room by room. An LLM estimates, for each candidate room and opening
  onto unseen space, the probability that searching it finds the target, without being told how far away it is.
  Geodesic travel costs measured on the observed map, together with a search cost per node, turn these
  estimates into the visit order that minimises the expected time to find the target, solved exactly with RPT$^*$ when the instance meets the solver's precondition. The agent executes only the first visit and
  then replans.}
  \label{fig:teaser}
\end{figure}

The contribution is this coupling of commonsense probabilities over geometrically grounded rooms, travel costs
measured on the observed map, an ordering that is certified optimal for each decision's instance, and closed-loop
replanning that revises all three as evidence accrues; it is not a new detector, language model or solver. We are
also explicit about what the guarantee covers. It holds for the routing instance built from the current map and
probabilities, which contains at most twelve candidates (the nearest and the eleven most probable), not for the
partially observed search problem as a whole, and the probabilities are model estimates, not calibrated
quantities. The certificate is also lost in two cases (\cref{sec:planning}): if the solver's time budget of 2\,s
expires, it returns the best complete order found so far, or a nearest-neighbour order; and when an instance violates the triangle inequality that the
exact solver requires, the implementation instead draws the next node at random in proportion to its probability
(\cref{tab:runtime}). In summary, we make the following contributions.
\begin{itemize}[leftmargin=*,itemsep=2pt]
  \item \ours, an integrated, training-free ObjectNav agent built from open-weight models that runs perception,
  geometric room segmentation, LLM room labelling, search and target closing in one closed loop
  (\cref{sec:perception,sec:planning,sec:execution}).
  \item A probabilistic room-search formulation in which independent LLM search-success probabilities for rooms
  and openings, discounted by fixed geometric rules, and a fixed probability for sighted target landmarks, are combined with geodesic travel costs and
  per-node service costs into a Hamiltonian path problem with probabilistic terminals; RPT$^*$ solves exactly
  each instance that meets its metric precondition within its time budget, for the visit order of minimum expected
  time-to-find, with a disclosed fallback otherwise, and the order is executed in a receding horizon with explicit
  anti-oscillation rules (\cref{sec:planning}).
  \item A verified target-closing routine that confirms a detection over consecutive frames, approaches a standoff
  pose, decides to stop on a fresh, centred, confident and close detection or, in a few bounded cases, on recent
  evidence of one, and issues STOP after at most two further forward steps, releasing false candidates back to the search (\cref{sec:execution}).
  \item An evaluation on four datasets and five splits, Gibson, HM3D v1 and v2, MP3D and RoboTHOR (1{,}000, 2{,}000,
  1{,}000, 2{,}195 and 1{,}800 validation episodes), under each benchmark's own success rule, with every privileged
  input disclosed (\cref{tab:privileges}) and with ablations of the main design choices
  (\cref{sec:experiments,sec:ablations}). \TODO{final list of ablations} \TODO{headline results once the runs
  are final: \SR and \SPL of \ours on each split, compared with the strongest zero-shot baseline on that split,
  with the comparability caveats of \cref{sec:experiments}}
\end{itemize}

We review related work in \cref{sec:related} and formalise the task and the search objective in
\cref{sec:problem}. \Cref{sec:perception} describes the online scene graph, \cref{sec:planning} the probabilistic
room search, and \cref{sec:execution} execution, local exploration and target closing. The experimental setup and
results follow in \cref{sec:experiments,sec:results}, and limitations and conclusions in
\cref{sec:limitations,sec:conclusion}. The appendices give the sub-routines and hyper-parameters
(\cref{app:algorithms,app:hyperparams}), the LLM prompts (\cref{app:prompts}), the benchmark protocols
(\cref{app:protocols}), and additional results (\cref{app:additional}).

% ============================================================================
% 03_related
% ============================================================================
% ============================================================================
%  03 -- Related Work
% ============================================================================
\section{Related Work}
\label{sec:related}

\ours draws on five lines of work, reviewed below together with how \ours departs from each.

\paragraph{Object-goal navigation and learned agents.}
ObjectNav asks an embodied agent to reach an instance of a named object category in an unseen
environment and to declare success with a STOP action~\citep{batra2020objectnav}; agents are scored by
success rate and SPL~\citep{anderson2018evaluation}, and often also by SoftSPL~\citep{datta2021integrating}.
The standard testbeds are the Habitat simulator~\citep{savva2019habitat,szot2021habitat} with scenes from
Gibson~\citep{xia2018gibson}, Matterport3D~\citep{chang2017matterport3d} and
HM3D~\citep{ramakrishnan2021habitat,yadav2023habitat}, and the AI2-THOR family~\citep{kolve2017ai2} with
RoboTHOR~\citep{deitke2020robothor} and the procedurally generated ProcTHOR houses~\citep{deitke2022procthor}.
Learned approaches fall into two families. End-to-end policies map observations directly to actions:
Habitat-Web~\citep{ramrakhya2022habitat} imitates crowd-sourced human demonstrations,
PIRLNav~\citep{ramrakhya2023pirlnav} follows behaviour cloning with reinforcement-learning fine-tuning,
EmbCLIP~\citep{khandelwal2022simple} trains agents on frozen CLIP visual backbones, and
OVRL-v2~\citep{yadav2023ovrl} uses a recurrent ViT policy with no map or detector at all. Modular
approaches keep an explicit semantic map and learn only where to go next: SemExp~\citep{chaplot2020object}
learns a goal-oriented exploration policy over an episodic semantic map; PONI~\citep{ramakrishnan2022poni}
learns, by supervision on top-down semantic maps and without any interaction with the environment, to
predict area and object potential functions; 3D-aware ObjectNav~\citep{zhang20233d} trains separate
exploration and target-identification policies on a 3-D point-based scene representation;
SGM~\citep{zhang2024imagine} learns a self-supervised generative map; and
T-Diff~\citep{yu2024trajectory} guides navigation with trajectory diffusion. Stubborn~\citep{luo2022stubborn}
showed that a semantics-agnostic explorer with careful stuck-escape and multi-frame aggregation of
detections can beat semantic policies, a reminder that low-level execution and detection filtering matter
alongside the choice of long-term goal. In a large real-world study, \citet{gervet2023navigating} found
that modular, map-based pipelines transfer from simulation to reality far better than end-to-end policies.
\ours keeps the modular decomposition into a metric map, a long-term goal and a local planner, but
involves no training. The learned long-term-goal policy is replaced by language-model probabilities over
room-level nodes and an exact expected-time ordering; no navigation policy is trained, whether on ObjectNav
data or on another task in simulation; and the pretrained detector and language models are used without
fine-tuning. PONI is the closest learned analogue, since it too estimates where the target is likely to
be found; our estimate comes from an LLM reading the observed scene graph rather than from a network
trained on semantic maps of the benchmark's training scenes.

\paragraph{Zero-shot ObjectNav with foundation models.}
Zero-shot methods use no ObjectNav training. A first group still trains a navigation policy, but on
another task: ZSON~\citep{majumdar2022zson} trains an image-goal agent on CLIP~\citep{radford2021learning}
goal embeddings and substitutes text embeddings at test time, PSL~\citep{sun2024prioritized} trains on
image-goal navigation with prioritized semantic learning, and PixNav~\citep{cai2024bridging} drives a
learned pixel-goal skill with an LLM planner. A second group trains no navigation policy for the task and
instead composes pretrained perception and language models with classical frontier
exploration~\citep{yamauchi1997frontier}, although the boundary between the groups is blurred: some
members still drive to their frontiers with a point-goal policy trained in simulation, as
VLFM~\citep{yokoyama2024vlfm} does, or rely on a segmentation network fine-tuned in domain, as
L3MVN~\citep{yu2023l3mvn} and TriHelper~\citep{zhang2024trihelper} do. CoW~\citep{gadre2023cows} pairs
frontier exploration with CLIP- or OWL-ViT-based target localisation, and most later work keeps a
frontier, or a location on a top-down map, as the unit of decision while changing how it is scored: by
soft commonsense constraints over objects and rooms in ESC~\citep{zhou2023esc}; by a language model reading the objects near each frontier
in L3MVN~\citep{yu2023l3mvn}; by LLM ``semantic guesswork'' used as a planning heuristic in
LFG~\citep{shah2023navigation}; by an LLM choosing a direction from detected objects in
LGX~\citep{dorbala2024can}; by a semantic frontier utility over a training-free scene graph in
SemUtil~\citep{chen2023how}; by a BLIP-2~\citep{li2023blip} image--text value map in
VLFM~\citep{yokoyama2024vlfm}; by multi-scale geometric-part and affordance maps in
GAMap~\citep{yuan2024gamap}; by a vision-language semantic score map in OpenFMNav~\citep{kuang2024openfmnav};
by multi-sourced value maps in an LLM-driven chain of navigation in InstructNav~\citep{long2024instructnav};
by a vision-language world model with a curiosity value map in WMNav~\citep{nie2025wmnav}; or by
multimodal spatial reasoning on a top-view map in TopV-Nav~\citep{zhong2024topvnav}.
TriHelper~\citep{zhang2024trihelper} adds collision, exploration and detection helpers to such a pipeline,
ApexNav~\citep{zhang2025apexnav} switches between semantic and geometric frontier exploration and fuses
detections around the target, and VoroNav~\citep{wu2024voronav} lets an LLM choose among nodes of a reduced
Voronoi graph described in text. Closest to our representation, SG-Nav~\citep{yin2024sg} builds an online
3-D scene graph, prompts an LLM to score frontiers from it, and re-perceives the target to reject false
detections; UniGoal~\citep{yin2025unigoal} extends such a graph to several goal types;
CogNav~\citep{cao2025cognav} lets an LLM drive a finite-state cognitive process over a heterogeneous map;
and OSG-Nav~\citep{loo2025open} supplies a schema-defined open scene graph as memory for foundation-model
reasoning. A few methods also optimise beyond the next step: in its semantic mode ApexNav orders a subset
of high-scoring frontiers by a traveling-salesman tour, which minimises path length but ignores that the
search may end early; BeliefMapNav~\citep{zhou2025beliefmapnav} fuses LLM priors with observations in a
3-D voxel belief map and plans by expected distance-to-detect; and WTRP-Searcher~\citep{liu2026global},
closest to \ours in spirit, casts the search as a weighted traveling repairman problem over VLM-scored
frontier viewpoints.

\ours differs from this body of work in three respects. First, its unit of decision is a place in an
online scene graph: an unfinished room or an opening onto unseen space, both derived from the geometry of
the observed map, or a sighted but unverified target. The language model therefore reasons about a few
semantically meaningful places rather than about individual frontier points, and it reasons over a text
description of the graph, whose objects come from an open-vocabulary detector~\citep{cheng2024yolo},
rather than over images. Second, semantics and geometry are kept apart: the LLM scores each room and
opening without seeing distances, and travel enters only through a geodesic cost matrix. Third, the
candidate nodes are ordered jointly to minimise the expected time to find the target, and the routing
instance posed at a decision point is solved exactly rather than greedily, with a fallback, described
in~\cref{sec:planning}, for instances that violate the solver's assumptions; most methods above instead
commit to the single best-scoring frontier.

\paragraph{Scene graphs and room-level semantic maps.}
3-D scene graphs organise a scene into a layered graph of building, rooms, objects and
cameras~\citep{armeni20193d}. 3-D dynamic scene graphs~\citep{rosinol20203d}, in the Kimera line of work,
build such hierarchies, with places and rooms, automatically from visual-inertial data, and
Hydra~\citep{hughes2022hydra} does so
incrementally and in real time, segmenting rooms from a layer of places. Open-vocabulary variants fuse 2-D
foundation models into object-centric graphs that an LLM can query for planning, as in
ConceptGraphs~\citep{gu2024conceptgraphs}, or into floor, room and object hierarchies with bird's-eye-view
room segmentation for language-grounded navigation, as in HOV-SG~\citep{werby2024hierarchical}. LLM
planners then exploit the graph: SayPlan~\citep{rana2023sayplan} searches a collapsed scene graph and
replans iteratively, SayNav~\citep{rajvanshi2024saynav} grounds step-wise LLM search plans in an
incrementally built graph, MoMa-LLM~\citep{honerkamp2024language} performs interactive object search over
dynamically built, room-segmented scene graphs, and MORE~\citep{mohammadi2025more} combines room
segmentation with active filtering of task-relevant objects for rearrangement. Room segmentation of 2-D
grid maps has a long history: \citet{bormann2016room} benchmark morphological, distance-transform, Voronoi
and feature-based methods, and the watershed transform~\citep{vincent1991watersheds} is a standard tool of
the distance-transform family. Inferring a room's category from the objects it contains, with doors as
properties of places, goes back at least to the probabilistic semantic mapping of
\citet{pronobis2012large}. \ours borrows the layered object--room structure but builds a lightweight 2-D
graph during the episode from RGB-D and pose alone: rooms come from a clearance watershed over the
observed free space, cut at confirmed doorways and at narrow passages the agent has walked through, and
re-segmented as the map grows; objects enter the graph only after they have been detected in more than
one frame; and an LLM labels each room from its confirmed objects, marking the label firm or tentative
depending on the evidence. The graph is not only a memory to be
queried: its rooms and openings are, directly, the terminals of the routing problem solved at each
decision point.

\paragraph{Language-model commonsense for planning.}
Several systems ground the commonsense of large language models in robot decisions.
SayCan~\citep{ichter2023do} multiplies an LLM's estimate of a skill's usefulness by a learned value
function's estimate that the skill can succeed; LM-Nav~\citep{shah2022lm} composes pretrained
language, vision and navigation models over a topological graph without fine-tuning;
VLMaps~\citep{huang2023visual} fuses visual-language features into a map to reach open-vocabulary and
spatial goals; and NavGPT~\citep{zhou2024navgpt} reasons explicitly with an LLM for vision-and-language
navigation. Prompting a model to write intermediate reasoning before its answer improves multi-step
reasoning~\citep{wei2022chain}. \ours follows SayCan's principle of separating what the language model is
good at, semantic plausibility, from what it is not, geometry and feasibility, but couples the two
through an expected-cost objective over a whole visit order rather than a per-step product. The LLM never
sees distances, so it cannot double-count the travel that the cost matrix already contains, and its reply
follows a reasoning-first schema, in the spirit of chain-of-thought prompting, before it states any
probability. Both the room labeller and the probability oracle are open-weight, instruction-tuned Qwen2.5
models~\citep{yang2024qwen25} served locally.

\paragraph{Probabilistic object search and routing.}
Planning a search with priors over object locations predates foundation models.
\citet{kollar2009utilizing} learn object--object and object--scene co-occurrence priors from Flickr
captions and use them to plan searches for novel objects, and \citet{aydemir2013active} search unknown
environments at the level of rooms, using probabilistic room and object priors and a planner over
uncertain semantics. Classical search theory~\citep{stone1975theory} allocates search effort by detection
probability to maximise the chance of detection or to minimise the expected time to find. When the search
visits discrete locations, minimising the expected time to find becomes a routing problem of the
minimum-latency, or traveling repairman, type~\citep{blum1994minimum}, which minimises the sum of arrival
times rather than the length of the tour. A weighted repairman objective, as in
WTRP-Searcher~\citep{liu2026global}, charges each vertex a fixed weight times its arrival time, which
equals the expected time to find when the weights are the probabilities of mutually exclusive target
locations. RPT$^*$~\citep{lyu2026rpt} instead formalises target search as a Hamiltonian path problem with
probabilistic terminals: each vertex ends the search with its own probability, each arc is weighted by the
probability that the search is still running when it is traversed, and the problem is solved optimally by
an A$^*$-style best-first search with an admissible heuristic and dominance pruning. Its authors embed the
solver in a hierarchical search system that regenerates and re-solves the problem as the robot moves; in
unknown environments its terminals are clustered frontiers whose probabilities combine the density of
nearby unknown space, the fraction of unknown space visible from the frontier, and a supplied
Gaussian-mixture prior over object locations.
BeliefMapNav~\citep{zhou2025beliefmapnav} pursues a related expected-distance criterion, but over a dense
voxel belief map rather than over a discrete set of places. \ours adopts RPT$^*$, re-implemented from its
published description, as its solver and keeps its receding-horizon use, executing only the first node of
each optimal order and replanning at visit boundaries. It feeds the solver a different instance, however:
the terminals are room-level nodes of an online scene graph; the probabilities of rooms and openings come
from an LLM's commonsense judgement of the observed house, scaled down by deterministic rules for openings
onto little unseen floor and for rooms entered and left, and are, consistently with independent
terminals, not normalised across nodes; and the arc costs combine geodesic travel with a node-specific
service time. The resulting order is optimal only for the instance posed at each decision point, not for the underlying
partially observable search problem. Relative to WTRP-Searcher, the closest zero-shot system, \ours uses
room-level nodes instead of frontier viewpoints, LLM success probabilities instead of VLM similarity scores,
and the probabilistic-terminal objective, solved to optimality for each instance that meets the solver's
precondition, instead of a weighted-latency objective \TODO{state whether WTRP-Searcher solves its routing
instances exactly or heuristically}. Our contribution is the closed-loop coupling summarised in
\cref{sec:intro}, not a new solver.

% ============================================================================
% 04_problem
% ============================================================================
% ============================================================================
%  Section 4: Problem Formulation and Overview
% ============================================================================
\section{Problem Formulation and Overview}
\label{sec:problem}

\subsection{Task and Setting}
\label{sec:problem-task}

\paragraph{Episodes.}
In each episode of object-goal navigation (ObjectNav)~\citep{batra2020objectnav}, an agent is placed at a start pose
$x_0$ in a previously unseen indoor environment and given a target category $\cstar \in \mathcal{C}$ from a fixed,
benchmark-specific set of 6 to 21 categories (\cref{tab:benchmarks}). At every step $t$ it receives an RGB image $I_t$,
a metric depth image $D_t$ and its pose $x_t$ (position, heading and, where the camera tilts, pitch); in our experiments
the pose comes from the simulator, a privilege listed in \cref{tab:privileges}. It answers with one discrete action
$a_t$: STOP, MOVE\_FORWARD (a $0.25$\,m step), TURN\_LEFT or TURN\_RIGHT (a $30^\circ$ turn); HM3D, MP3D and
RoboTHOR add LOOK\_UP and LOOK\_DOWN, which tilt the camera by $30^\circ$, and we add the same pair on Gibson, a
deviation from its published action set (\cref{sec:experiments-privileges}). The episode ends at step
$t_{\mathrm{end}}$, when the agent calls STOP or when the budget of
$T = 500$ actions is spent.

\paragraph{Success.}
In a general form, the success indicator of an episode is
\begin{equation}
  S \;=\; \mathbb{I}\left[\, \operatorname{dist}\left(x_{t_{\mathrm{end}}}, \mathcal{X}_{\cstar}\right) \le \delta
  \;\wedge\; \nu_{\cstar}\left(x_{t_{\mathrm{end}}}\right) \,\right] \cdot
  \mathbb{I}\left[\, a_{t_{\mathrm{end}}} = \text{STOP} \,\right]^{\rho} \;\in\; \{0, 1\},
  \label{eq:success}
\end{equation}
and the episode succeeds when $S = 1$. Here $\mathcal{X}_{\cstar}$ is the goal set (the instances of $\cstar$, view
points around them, or a region dilated around them), $\operatorname{dist}$ a geodesic distance, $\delta$ the success
threshold (a strict inequality on HM3D and MP3D), $\nu_{\cstar}$ a visibility test that is identically true where the
protocol has none, and $\rho \in \{0, 1\}$ indicates whether STOP is required. Gibson does not require
STOP ($\rho = 0$), so an episode is also scored at its final pose when the budget runs out; HM3D, MP3D and
RoboTHOR require it ($\rho = 1$), and on RoboTHOR the criterion reduces to the visibility test, which carries a 1\,m
visibility distance. \cref{sec:experiments} gives each instantiation. \ours decides STOP from its own
detections and depth (\cref{sec:execution}) and never reads the evaluator's distances, goal positions or success
signal.

\paragraph{Zero-shot, training-free setting.}
We use \emph{zero-shot} in the training-free sense: \ours performs no ObjectNav training of any kind, neither
reinforcement nor imitation learning nor fine-tuning. Its learned components, the open-vocabulary detector
YOLO-World~\citep{cheng2024yolo} and two instruction-tuned language models, Qwen2.5-3B-Instruct and
Qwen2.5-14B-Instruct~\citep{yang2024qwen25}, are used frozen, and the target enters only as a category name. The agent
builds its own metric map and receives no prior map, goal coordinates or distance field; the privileged inputs it does
use (ground-truth pose, noise-free simulator depth with enclosed holes filled, stair geometry revealed only once a
staircase is visible, and the projection of its approach goal onto the navigable surface) are listed in
\cref{tab:privileges}. \cref{sec:experiments} also discloses the development protocol, including the validation episodes inspected
during development and the revisions of the search, of perception and mapping, of target closing and of Gibson's
action set that were informed by validation failures and, in part, selected by validation scores.

\subsection{The Search Objective}
\label{sec:problem-objective}

Searching an unseen building is a sequential decision problem under partial observability; rather than solving it as
a whole, \ours poses, at each decision point, a finite combinatorial problem built from what it has observed. Let
$\N = \{1, \dots, m\}$ be the current candidate nodes: rooms not yet finished, openings onto unseen space, and target
landmarks sighted but not yet verified (finished rooms do not re-enter, and at most twelve nodes are kept;
\cref{sec:planning}). Each node $i$ has an entry point $e_i$ and an estimated probability $p_i$ that searching node $i$
finds $\cstar$. A depot vertex $0$, the agent's pose $x_t$, is the fixed start of every order and has $p_0 = 0$. For
vertices $i \neq j$, $C_{ij}$ is the estimated cost, in actions (the unit of $T$), of travelling from vertex $i$ (the
agent's position for $i = 0$, otherwise $e_i$) to $e_j$ and searching node $j$ there, and $C_{ii} = 0$
(\cref{eq:arccost}). A visit order is a Hamiltonian path $\pi = (\pi_0 = 0, \pi_1, \dots, \pi_m)$ from the depot
through every node, and its expected time-to-find is
\begin{equation}
  J(\pi) \;=\; \sum_{k=0}^{m-1} q_k \, C_{\pi_k \pi_{k+1}},
  \qquad
  q_k \;=\; \prod_{j=0}^{k} \left(1 - p_{\pi_j}\right),
  \label{eq:objective}
\end{equation}
where $q_k$ is the probability that the target has not been found once the first $k$ nodes of the order have been
searched ($q_0 = 1$ because $p_0 = 0$); under the independent-terminal model below it is the probability that the arc
from $\pi_k$ to $\pi_{k+1}$ is travelled at all. Late arcs are thus discounted by the chance that the search has already
ended, and minimising $J$ weighs each node's probability against the cost of reaching and searching it, a trade-off
that neither always visiting the most probable node nor always the nearest one captures. This is the Hamiltonian path
problem with probabilistic terminals of \citet{lyu2026rpt}. We use paths rather than tours because the search ends
where the target is found; a tour would charge a return leg to the start that is never travelled. Related formulations
are the minimum-latency problem~\citep{blum1994minimum}, applied to frontier viewpoints by \citet{liu2026global}, and the
expected time to detection of classical search theory~\citep{stone1975theory}.

\paragraph{Independent terminals.}
The probabilities are deliberately not normalised: each $p_i$ is a separate estimate of the chance that searching node
$i$ finds the target, and $\sum_i p_i$ may differ from one. A normalised distribution would assert that the target lies
in exactly one current candidate, which is false whenever part of the building is not yet represented by any node, or
when the category has several instances, as chairs and plants often do. \Cref{eq:objective} then treats the outcomes of
the individual searches as independent events. In general they are not (a room holding the only instance makes every
other search fail), and this independent-terminal assumption is the approximation under which the objective is
defined. Because $p_i$ refers to \emph{searching} node $i$, it also absorbs whether the bounded local visit would detect
the object. Every $p_i$ is kept strictly below one, which keeps every $q_k$ positive and the solver's heuristic well
defined.

\paragraph{Receding horizon.}
The instance is a snapshot: as the agent moves, rooms are discovered, split or merged, openings appear and vanish, and
new objects change the probabilities. \ours therefore rebuilds the instance at each decision point, computes an order
minimising \cref{eq:objective} with the exact solver RPT$^*$~\citep{lyu2026rpt}, executes only $\pi_1$, discards the
rest of the order and plans again at the next decision point. The order is optimal for the instance it was computed
from, not for the underlying partially observed problem. RPT$^*$ certifies this optimality only when the costs satisfy
the triangle inequality and the search completes within its budget of 2\,s; because our costs mix distances measured at
two clearances, an instance can violate the inequality, and the implementation then draws the next node at random with
probability proportional to $p_i$ (\cref{sec:planning}).

\subsection{Method Overview}
\label{sec:problem-overview}

\ours couples four components (\cref{fig:overview}). \emph{Perception} (\cref{sec:perception}) maintains an online
scene graph $\G = (\M, \Ob, \D, \R)$: a 2-D occupancy grid $\M$, confirmed open-vocabulary object landmarks $\Ob$,
confirmed doorways $\D$, and rooms $\R$ obtained by a clearance watershed over observed free space and cut at the
doorways, which the smaller language model labels from the objects confirmed in them. \emph{Planning}
(\cref{sec:planning}) forms the candidate nodes, asks the larger language model for a probability of each room
and opening from a textual description of the house and of the node, discounts it by fixed rules for exits that would
reveal little and for rooms already entered to obtain $p_i$ (a sighted landmark receives a fixed
$p_i = 0.85$), computes $C$ by Dijkstra's algorithm~\citep{dijkstra1959note} on the occupancy grid, and solves
\cref{eq:objective}. Distances are withheld from the language model so that travel is charged once, through $C$,
rather than a second time by the model. \emph{Execution} (\cref{sec:execution}) follows a committed A$^*$
route~\citep{hart1968formal} over a clearance-weighted grid to the entry point of $\pi_1$ and performs a bounded local
visit there: in a room, a walk over its frontiers, best first by the unknown floor a look from each would reveal per
metre of travel, until none would reveal at least 1\,m$^2$ (at most 60 actions); at an opening, or at a sighted
landmark from a 1.5\,m standoff, a look along three headings, toward the unseen side or the object and $30^\circ$ to
either side; where the agent stopped short of an opening's passable threshold, it steps onto the threshold and widens
the look to $60^\circ$ to either side (at most 30 actions). \emph{Target closing}
(\cref{sec:execution}) preempts all of this on a confident detection of the target: it verifies the sighting,
approaches the target and inspects it, and decides to stop on a fresh, confident detection within reach or, in a
few bounded cases detailed in \cref{sec:execution}, on recent evidence of one, then calls STOP after at most two
further steps toward the target; otherwise it releases the candidate and
the interrupted search resumes.

\begin{figure}[t]
  \centering
  \fbox{\parbox[c][4.5cm][c]{0.95\linewidth}{\centering\TODO{overview figure: an RGB-D frame feeding the scene
  graph (occupancy grid, object landmarks, doorways, rooms with their LLM labels); the candidate nodes (rooms,
  openings, an unverified target sighting) annotated with LLM probabilities and Dijkstra costs; the RPT$^*$ visit
  order with only its first node executed; the committed A$^*$ route and the local visit (room frontier walk, opening look);
  a dashed preemption arrow from the detector to target closing (verify, approach, inspect, then Stop or
  release and resume)}}}
  \caption{Overview of \ours. Perception updates the scene graph $\G = (\M, \Ob, \D, \R)$ after every action. At
  each decision point the agent forms candidate nodes, obtains for each an independent probability $p_i$ (from a
  language model, discounted by fixed rules, except for sighted target landmarks) and a travel-plus-service cost matrix from the occupancy grid,
  and computes the visit order that minimises the expected time-to-find (\cref{eq:objective}). Only the first node is
  visited before the agent plans again. A confident target detection preempts the search and hands control to target
  closing.}
  \label{fig:overview}
\end{figure}

\paragraph{Main loop.}
\cref{alg:main} summarises the procedure; the sub-routines are given in \cref{app:algorithms}. After a warm-up of
ten $30^\circ$ turns in place, which records the spawn spot as a completed scan, the agent reaches a decision point
whenever no node is in force: when a visit ends (node finished, visit budget spent, node unreachable, transit blocked
or timed out) and after each fallback exploration step. At a decision point the agent queries the oracle on the
places it can describe, forms the solver's instance $\N$ from the reply, prices it and solves it. Eligibility filters,
deferral and the 12-node cap are given in \cref{sec:planning}; route commitment in \cref{sec:execution}. New rooms or
object kinds seen while a node is in force are acted upon only when its visit ends, which, together with the committed
route, guards against oscillation. A confident sighting of the target preempts the visit, the warm-up or a fallback
and hands control to target closing. If the closing calls STOP, the episode ends; if it releases the candidate, the
spot is remembered and the interrupted visit resumes; a closing that exceeds its action bound ends the episode as a
failure. Outside the closing, the agent also stops beside a target landmark confirmed from three views within
0.75\,m (\cref{sec:execution}). With nothing to plan over, or when the language model returns no valid answer,
\ours falls back to bounded frontier exploration~\citep{yamauchi1997frontier} rather than assuming a uniform prior
(\cref{sec:planning}).

\begin{algorithm}[tp]
\DontPrintSemicolon
\caption{\ours: LLM-guided object search}
\label{alg:main}
\Input{target category $\cstar$; RGB-D observations $(I_t, D_t)$ and pose $x_t$; action budget $T$}
\Output{\Stop (\Found) at a verified target or beside a confirmed target landmark, or \Fail when the budget or the
closing bound is spent}
\State{scene graph $\G$; described node set $\mathcal{U}$; solver instance $\N$; probabilities $p$ and pass verdict
$\omega$; costs $C$; visit order $\pi$; node in force $v$ (none at the start); committed route $\tau$; closing
outcome $o$}
\BlankLine
\Comment{Warm-up}
\WarmUp{} \Comment*[r]{ten $30^{\circ}$ turns in place}
$\G \gets$ \InitSceneGraph{$x_0$} \Comment*[r]{spawn spot recorded as a completed scan}
\While{$t < T$}{
  \If(\Comment*[f]{decision point}){no node is in force}{
    \Comment{Scene graph and oracle}
    $\G \gets$ \ClassifyRooms{$\G$} \Comment*[r]{\cref{alg:classifyrooms}: LLM labels rooms from their objects}
    $\mathcal{U} \gets$ every unfinished room of the storey with its status, the openings not yet peeked and not covered by such a room's frontier walk, the sighted targets\;
    \lIf(\Comment*[f]{nothing to plan over}){no node of $\mathcal{U}$ could be offered to the solver}{\ExploreFrontier{}; \KwContinue}
    $(p, \omega) \gets$ \LLMProbabilities{$\cstar, \G, \mathcal{U}$} \Comment*[r]{\cref{alg:llmprobabilities}: independent $p_i$}
    \lIf(\Comment*[f]{back-off; never a uniform prior}){no valid LLM reply}{\ExploreFrontier{}; \KwContinue}
    \Comment{Solver instance and costs}
    $\N \gets$ \CandidateNodes{$\G, p, \omega$} \Comment*[r]{\cref{alg:candidatenodes}: eligibility, at most 12}
    \lIf(\Comment*[f]{nothing worth a visit}){$\N = \varnothing$}{\ExploreFrontier{}; \KwContinue}
    $C \gets$ \Dijkstra{$\G, x_t, \N$} \Comment*[r]{travel plus service, \cref{eq:arccost}}
    \Comment{Visit order (receding horizon)}
    $\pi \gets$ \RPT{$\N, p, C$};\quad $v \gets \pi_1$ \Comment*[r]{minimises \cref{eq:objective}; only $\pi_1$ runs}
    $\tau \gets$ \AStar{$x_t, e_v$} \Comment*[r]{committed route to the entry point of $v$}
  }
  \Comment{Navigation and local exploration}
  \Repeat{$v$ finished \KwOr visit budget spent \KwOr $v$ unreachable or blocked \KwOr transit timed out \KwOr
  $\cstar$ sighted \KwOr $t \ge T$}{
    \uIf{not yet arrived at $e_v$}{
      \Navigate{$\tau$} \Comment*[r]{\cref{alg:navigate}: one action; $\tau$ re-planned only if the goal moved, or it is blocked, left or stalled}
    }
    \Else{
      \ExploreLocal{$v$} \Comment*[r]{\cref{alg:explorelocal}: one action of the frontier walk or look}
    }
    \Comment{each call ends with $\G \gets$ \UpdateSceneGraph{$\G, I_t, D_t, x_t$} (\cref{alg:updatesg})}
    \lIf(\Comment*[f]{outside the closing}){a target landmark confirmed from three views lies within 0.75\,m and was not rejected}{\KwRet{\Found}}
    \lIf(\Comment*[f]{mid-visit trigger}){a new object kind appears in the room being searched or crossed}{$\G \gets$ \ClassifyRooms{$\G$}}
    \Comment{other new rooms and object kinds wait for the visit's end; no replanning here}
  }
  \uIf(\Comment*[f]{the sighting preempts the visit}){$\cstar$ sighted}{
    $o \gets$ \TargetClosing{$\cstar$} \Comment*[r]{\cref{alg:targetclosing}: verify, approach, inspect}
    \lIf{$o =$ \Stop}{\KwRet{\Found}}
    \lIf(\Comment*[f]{160-action closing bound}){$o =$ \Fail}{\KwRet{\Fail}}
    resume the interrupted visit of $v$ \Comment*[r]{$o$ released the candidate and stored its spot}
  }
  \lElse(\Comment*[f]{repeated failures defer $v$}){release $v$ with its verdict, leaving no node in force}
}
\KwRet{\Fail}\;
\end{algorithm}

% ============================================================================
% 05_perception
% ============================================================================
% ============================================================================
%  Section 5: Online scene graph (UpdateSceneGraph, ClassifyRooms)
% ============================================================================
\section{Online Scene Graph}\label{sec:perception}

\ours represents what it has observed as a scene graph $\G = (\M, \Ob, \D, \R)$: a 2-D occupancy grid $\M$, confirmed
object landmarks $\Ob$, confirmed doorways $\D$, and rooms $\R$ that partition the observed free space. Unlike layered
3-D scene graphs~\citep{armeni20193d,rosinol20203d,hughes2022hydra}, it is organised on the floor plane. It is updated from the
RGB image $I_t$, the depth image $D_t$ and the pose $x_t$ after every action, without a language model
(\cref{alg:updatesg}); room labels come from a separate, event-triggered step, \textsc{ClassifyRooms}
(\cref{alg:classifyrooms}). Objects and doorways reach the room classifier, the search oracle (\cref{sec:planning}) and
the planner only after confirmation across frames, so one spurious frame cannot redirect the search. The exception is
an object of the target category: one sighting at a confidence of at least 0.50 makes it a node for the planner. The pose is the
simulator's ground truth, and enclosed holes in the noise-free simulator depth (glass, mirrors, screens) are filled
from the surrounding valid pixels before use; both are listed with the other privileged inputs and deviations in
\cref{tab:privileges}. All constants are collected in \cref{tab:hyperparams}.

\subsection{Open-vocabulary detection}\label{sec:perception-detection}

We detect objects with YOLO-World~\citep{cheng2024yolo} (X-size v2, with a CLIP~\citep{radford2021learning} ViT-B/32
text encoder) at 640\,px input, with non-maximum suppression at IoU~0.5 and at most 100 boxes per frame. The detector
is called exactly once per frame with a fixed vocabulary. On Gibson it has 30 prompts: eight goal prompts for the six
categories (with the aliases couch/sofa and tv/television), thirteen room-context objects that inform the room
classifier, four door prompts, two stair prompts, and three distractors (toy, rocking horse, stuffed animal)
\TODO{list the vocabularies used on HM3D, MP3D and RoboTHOR in \cref{app:protocols}}. The distractors are never goals;
they exist to absorb false matches, because an open-vocabulary detector assigns a region to the closest prompt it is
given, so a child's ride-on horse could otherwise be read as a confident chair. Within a frame, a box is discarded if
a stronger box of the same class (aliases merged) overlaps it with IoU of at least 0.5, or a stronger box of any class
with IoU of at least 0.9, since one crop then answers two prompts.

Three confidence thresholds match the cost of an error at each consumer. The detector emits every box from 0.05, the
floor of the door channel: a missed doorway costs a room boundary, whereas a spurious one must still pass geometric and
multi-view tests (\cref{sec:perception-doors}). Object boxes enter the landmark map only from 0.30, a moderate floor
because a landmark must in addition be re-observed before anything acts on it, unless it is of the target category and
was seen at 0.50 or more. A target box starts target closing
(\cref{sec:execution}) only from 0.50, because a takeover commits up to 160 actions to one hypothesis.

\subsection{3-D localisation and multi-frame landmark fusion}\label{sec:perception-fusion}

A bounding box overhangs its object, so the range $d^{*}$ of a retained box $b$ is the 30th percentile of the valid
depths (0.5--5.0\,m) in its central half, which favours the near surface over background seen past the object's edges.
Pixels of that central region within $\max(0.12\,\mathrm{m},\, 0.06\,d^{*})$ of $d^{*}$ must form a connected component
covering at least 35\% of its valid pixels; otherwise the box straddles several surfaces and is rejected. The position
$X_b$ is the per-axis median of the component's points back-projected into the world frame, and the footprint radius is
$\rho_b = w_b\, d^{*} / (2 f_x)$ for a box $w_b$ pixels wide and focal length $f_x$, clamped to 0.15--1.5\,m. Door and
stair prompts never become landmarks. A detection is fused only if it lies between 0.30\,m below and 2.2\,m above the
floor and is supported by observed floor near it, and frames taken from a tilted inspection view are not fused.

Landmarks are fused strictly per class: a detection joins the nearest landmark of its class whose centre lies within
0.35\,m of $X_b$, or whose footprint disc overlaps the detection's with IoU of at least 0.25, provided the heights agree
within 0.35\,m; otherwise it starts a new candidate. A landmark counts at most once per frame and is confirmed into
$\Ob$ after observations in two distinct frames. As in the multi-frame aggregation of Stubborn~\citep{luo2022stubborn},
this keeps single-frame false positives, such as a one-frame flicker of a wrong class, out of the graph. Classes never
merge, because objects of different classes often rest on or beside one another, such as a toilet and the bathtub next
to it, and a cross-class merge folds them into one landmark whose label can flip to the other class, losing a target
and the room evidence it carries.

\subsection{Occupancy mapping}\label{sec:perception-occupancy}

$\M$ is a log-odds grid of 0.10\,m cells covering 80\,m$\times$80\,m, built from every eighth depth pixel in each image
direction, with valid depths of 0.5--5.0\,m. Points 0.15--0.93\,m above the floor are obstacle hits; the band ignores
low floor clutter such as rugs and door sills, and anything the 0.88\,m tall agent passes under. On RoboTHOR, whose
sensor returns 0.1--5\,m and whose agent is 0.9\,m tall, \TODO{state the depth-validity range and obstacle band
used}. Free space comes only
from two sources: 2-D rays cast from the camera to obstacle hits, and directly observed floor points (within
$\pm 0.10$\,m of the floor plane), the latter never clearing a cell that is already occupied. Pixels without a valid
return never clear space, since missing or out-of-range depth, as on glass that hole filling did not repair or on walls
beyond 5\,m, is no evidence of free floor. A cell is free at
an occupancy probability of at most 0.45, occupied at 0.65 or more and unknown in between, so that one ray pass frees a
cell while a weak, isolated hit stays unknown. The map itself is not inflated: the planner adds clearance (0.30\,m
preferred, 0.18\,m minimum; \cref{sec:execution}), so the same map serves both standoffs and the room segmentation
below. Frontiers are free cells adjacent to unknown ones~\citep{yamauchi1997frontier}. One grid and one scene graph are kept per floor;
with stair traversal disabled, as on the single-storey Gibson split, the agent stays on its spawn storey and only that
storey's map is used \TODO{state whether stair traversal is enabled on HM3D and MP3D}.

\subsection{Doorway detection and confirmation}\label{sec:perception-doors}

Doorways are read from the detector's raw boxes, before the per-frame de-duplication of
\cref{sec:perception-detection}, and only in level views. A box from the four door prompts (confidence $\geq 0.05$) is
dropped if a confident ($\geq 0.5$) non-door box overlaps it with IoU $\geq 0.6$, since a confidently detected cabinet
does not bound a doorway. Its range is read from the jamb and
lintel bands of the box rather than from its centre, which in an open doorway sees the next room's far wall. Geometric
gates then require a passable opening: a box at least 0.9 times as tall as it is wide, an apparent width of
0.35--1.5\,m, a top at least 1.4\,m above the floor and a bottom within 0.5\,m of it (each height test waived when the
box is cut by the image border). Candidates of all door prompts are fused into one class within 0.6\,m, and a doorway
is confirmed into $\D$ after three observations, at least one of them from a viewpoint 0.2\,m or 20$^\circ$ away from
the first. Frames from one pose are correlated, and the viewpoint condition stops one appearance-driven false
positive, such as a wardrobe front, from confirming itself.

\subsection{Room segmentation}\label{sec:perception-rooms}

Rooms are recomputed after every action from the free cells of $\M$, excluding the footprints of seen staircases and
healed by one morphological closing and one opening. With $\Phi(x)$ the Euclidean clearance of a free cell $x$, the
seeds are the local maxima of $\Phi$ with $\Phi \geq 0.3$\,m, at least 1.0\,m apart, and rooms are the basins of a
watershed~\citep{vincent1991watersheds} of $-\Phi$ flooded from them: wide areas, bounded by the narrow places between
them. Adjacent basins whose dynamics (the clearance lost from the shallower peak down to the saddle) is below 0.2\,m
are merged, which repairs L-shaped or furniture-split rooms, and regions under 0.5\,m$^2$ are dropped. Confirmed
doorways are carved out of the flooding mask and always separate rooms: basins on opposite sides of a cut are never
merged. Because a detected door inherits the error of its depth estimate, the cut is snapped to the narrowest severing
passage within 0.9\,m, or, if none exists, carved as a disc of 0.75\,m radius. The agent's trajectory adds a second
kind of cut, the set $\Theta$ of walked-through thresholds: a clearance dip along the travelled path to at most
0.55\,m, with at least 0.30\,m more clearance on both sides within 1.5\,m of travel, marks a threshold, which is cut
only where a severing passage lies within 0.5\,m. A uniformly narrow corridor never dips and is therefore never cut.
These cuts let a room boundary appear where the detector never confirmed a door.

This construction belongs to the distance-transform family of grid-map room segmentation surveyed by
\citet{bormann2016room}. It differs from separating rooms on a sparsified Voronoi graph cut at
doors~\citep{honerkamp2024language,mohammadi2025more} and from segmenting the bird's-eye view of a completed
reconstruction~\citep{werby2024hierarchical}. On a partial map, walls exist only where they have been seen, so the
medial axis of the explored space is one connected component that door cuts alone cannot separate. In a clearance
watershed a missing wall merely lowers the ridge between two basins, so the partition degrades gracefully while
coverage grows instead of collapsing into one dominant room.

Room identities persist by matching each new region to its predecessors (a region holding a room's birth cell, then
IoU $\geq 0.15$, then at least 60\% containment of the smaller region); a room that disappears in one partition can
reclaim its identifier when a matching region returns, and identifiers are never reused, so the labels and search
history attached to a room survive the re-segmentation after every action. Each room $r$ carries its area $A_r$; its
objects $\Ob_r \subseteq \Ob$, the confirmed landmarks whose cell lies in $r$ or, since furniture stands on occupied
cells, whose nearest cell of $r$ lies within the landmark's footprint radius plus 0.6\,m; its frontier clusters that
the agent can reach with clearance; its search effort (the actions spent inside it); its recency (the time since the
agent last stood in it); and its status (never entered, entered, or finished; \cref{sec:planning}). Together with the label
$\ell_r$, these attributes describe the room to the oracle.

\subsection{LLM room classification}\label{sec:perception-labels}

The objects a room contains are strong cues for its category~\citep{pronobis2012large}. \textsc{ClassifyRooms} gives
Qwen2.5-3B-Instruct~\citep{yang2024qwen25}, served locally and decoded greedily (temperature 0, fixed seed) in JSON
mode, the room's confirmed objects as class--count pairs, and asks for a label from a fixed set, a confidence and a
one-sentence reason (\cref{app:prompts}). The prompt warns that observations are incomplete and may contain false
detections, and asks for \emph{unknown} when the evidence is insufficient or contradictory. The label set has nine
residential types (kitchen, bedroom, bathroom, living room, dining room, office, hallway, storage closet, laundry
room), five institutional ones (lobby, waiting area, patient room, exam room, reception) and \emph{unknown}; a reply
outside the set is read as \emph{unknown} \TODO{the five institutional labels are carried over from an earlier
hospital deployment and are offered on every benchmark; justify them or remove them before the final runs}. A room with no confirmed object, or with only generic ones (cabinet, potted plant, book, vase, clock, cup,
bottle), stays \emph{unknown} without a query. Unknown is an explicit label, not a failure: such rooms remain
candidates, so the oracle and the planner can weigh discovering an unidentified room against searching a recognised
one (\cref{sec:planning}).

With $K_r$ the set of object classes in $\Ob_r$, $\kappa_r$ the returned confidence, $K_{\mathrm{gen}}$ the generic
classes and $\operatorname{sig}(k)$ the room type of a signature object $k$ (bed $\to$ bedroom; toilet, shower, bathtub
$\to$ bathroom; oven, stove, refrigerator $\to$ kitchen), the strength $\beta_r$ of a label $\ell_r$ is
\begin{equation*}
  \beta_r = \textit{firm} \quad\text{iff}\quad
  \kappa_r \geq 0.65 \;\wedge\; \ell_r \neq \textit{unknown} \;\wedge\;
  \left( \left( |K_r| \geq 2 \wedge K_r \setminus K_{\mathrm{gen}} \neq \emptyset \right)
  \vee \exists\, k \in K_r : \operatorname{sig}(k) = \ell_r \right),
\end{equation*}
and \emph{tentative} otherwise. One bed can thus give a firm ``bedroom'', whereas one chair gives only a tentative label,
which reaches the oracle marked with a question mark (``living room?'').

The classifier never runs on ordinary steps. It runs (i) at each decision point, immediately before the candidate
nodes are formed, for every room whose multiset of object classes has changed or that has not been re-queried for 50
actions, reusing a cached verdict for an identical multiset; and (ii) in the middle of a visit, once per new object
kind seen in the room being searched or traversed, so that its label is current when the visit ends. If the model
fails, the labels already held stand, no default labelling is substituted, and the agent explores for 25 actions before
asking again, doubling this back-off after each further failure up to 200 actions (\cref{sec:planning}).

\paragraph{Optional components.} Two further options are implemented but are off by default
\TODO{state whether either option is ablated in \cref{sec:ablations}}: reading a target box's depth on a box-prompted
SAM~\citep{kirillov2023segment} or MobileSAM~\citep{zhang2023faster} mask, so that an occluder in front of the target
cannot shorten its range; and Grounding DINO~\citep{liu2024grounding} proposals, alone or added to YOLO-World's, that
are kept only when BLIP-2~\citep{li2023blip} verifies them.

% ============================================================================
% 06_planning
% ============================================================================
% ============================================================================
%  06 -- Probabilistic Room Search (CandidateNodes, costs, LLMProbabilities,
%        RPT*, re-planning, exploration fallback, multiple storeys)
% ============================================================================
\section{Probabilistic Room Search}
\label{sec:planning}

At every decision point of \cref{alg:main} the agent chooses which part of the partially mapped home to search next, in four separate steps. \textsc{CandidateNodes} (\cref{alg:candidatenodes}) extracts a few places worth visiting from the scene graph $\G=(\M,\Ob,\D,\R)$ of \cref{sec:perception}; shortest paths on the occupancy grid, plus the time each visit takes, price moving between them; a language model, the \emph{oracle}, which sees what each place contains but not how far away it is, estimates how likely searching it is to reveal $\cstar$ (\cref{alg:llmprobabilities}); and RPT$^*$~\citep{lyu2026rpt} orders the places to minimise the expected time-to-find of \cref{eq:objective}. Only the first node of the order is executed before the agent decides again. All costs are measured in simulator actions.

% ----------------------------------------------------------------------------
\subsection{Candidate nodes}
\label{sec:planning-nodes}

The solver's instance $\N$ is drawn from three kinds of node, $\N_{\mathrm{room}}$, $\N_{\mathrm{open}}$ and $\N_{\mathrm{target}}$, plus stair nodes in the multi-storey mode (\cref{sec:planning-fallback}), and passes the eligibility filters below; each node $i$ has an entry point $e_i$, the place where its visit begins.

\emph{Rooms.} $\N_{\mathrm{room}}$ holds the unfinished rooms of the current storey (the finishing rules are given in \cref{sec:execution}); the warm-up rotation counts as a scan of the spawn room. A room's entry point is the frontier of the room whose view would reveal the most unknown floor per metre walked (\cref{sec:execution}); a room with no such frontier left is entered at its doorstep, the room's passable cell nearest the agent along the floor, failing which at its vantage point, the cell of the room with at least 0.4\,m of clearance that lies farthest from anything outside the room, preferring cells the agent can reach, or at the room's centroid if no cell qualifies.

\emph{Openings.} $\N_{\mathrm{open}}$ stands for unseen space that belongs to no room yet: reachable frontiers merged within 1.5\,m into exits at least 0.5\,m wide. We exclude frontiers within 1.2\,m of the agent, which mark the camera's blind zone rather than unexplored space, short frontiers in the shadow of confirmed floor-standing objects of at least 0.25\,m radius unless at least 2\,m$^2$ of unknown space lies within 1.5\,m of the frontier, frontiers with a seen staircase or a drop within 1.5\,m on their unknown side, and exits that a completed peek has already covered. An opening is a \emph{doorway} if a confirmed door lies within 1.0\,m and a \emph{gap} otherwise; its entry point is its threshold. An opening that leaves a room which is itself a node is folded into that room, whose visit walks its frontiers.

\emph{Sighted targets.} $\N_{\mathrm{target}}$ holds landmarks of the target class that are confirmed or were seen once with confidence at least 0.50, have not yet been inspected from close on this storey, and do not lie where the closing routine of \cref{sec:execution} has disproved a candidate from close. A landmark can be confirmed without closing ever having started, for instance when its detections were confident enough to enter the object map (\cref{sec:perception}) but none reached the takeover threshold of the closing routine. The entry point is the reachable cell, nearest to the agent along the floor, on a 1.5\,m standoff ring around the landmark, or on a 2.25\,m ring if no cell of the first is reachable.

\emph{Eligibility.} Each entry point is snapped to the nearest passable cell within 2\,m; a node with no such cell, or one unreachable from the agent, is dropped, as is a node whose scaled probability is $p_i<0.05$ (\cref{sec:planning-oracle}). A node whose transit has failed three times (unreachable, timed out or blocked) is deferred for 180 actions, and a released node cools down for 120 actions, or for 30 after such a failure; if every remaining node is cooling, those whose visit did not already exhaust them are admitted. An instance keeps at most 12 nodes: the one with the lowest cost from the agent, plus the 11 most probable of the rest. The reserved slot matters because the objective weights the first arc most heavily: a cap based on probability alone could discard the cheap nearby node that the full instance would visit first.

\emph{Second pass.} Every unfinished room, and the room whose visit is in force, is described to the oracle together with its status; other finished rooms are left off the node list, though the house summary still names them with their status (these rooms, the openings not yet peeked or folded and the sighted targets form the described set $\mathcal{U}$ of \cref{alg:llmprobabilities}). By default a finished room is never offered to the solver again, and once the storey is covered the exploration fallback (\cref{sec:planning-fallback}) carries the search. An optional second pass, off by default, offers finished rooms again while the storey is covered, with no unfinished room, no opening and no reachable frontier revealing at least 1\,m$^2$ left; the \emph{pass} verdict of the oracle's reply (\cref{sec:planning-oracle}) is requested and recorded but does not open it. A revisit scans from a vantage point at least 1.5\,m away from every earlier scan point, so that it sees the room from a new position.

% ----------------------------------------------------------------------------
\subsection{Travel and service costs}
\label{sec:planning-costs}

The agent (the depot, vertex $0$) and the candidates form a complete graph with arc weights
\begin{equation}
  C_{ij} \;=\; \frac{d_{ij}+\lambda_i+\lambda_j}{u} \;+\; \sigma_j \quad (i\neq j), \qquad C_{ii}=0 .
  \label{eq:arccost}
\end{equation}
Here $d_{ij}$ is the shortest-path length between $e_i$ and $e_j$ (between the agent's position and $e_j$ for $i=0$), computed by Dijkstra's algorithm~\citep{dijkstra1959note} on the 8-connected 0.1\,m grid with unknown cells blocked. A pair is measured at the 0.30\,m preferred clearance where such a route connects it, and at the 0.18\,m body radius otherwise, matching the routes that the A$^*$ planner of \cref{sec:execution} prefers except where it takes the body-radius route because the clearance route would be more than 1.5 times as long. Cells are binary (passable or not) rather than carrying the planner's clearance-shaped cost: the shaped cost inflates lengths near walls and is not measured in metres, whereas binary weights keep $d_{ij}$ a symmetric length in metres. The speed $u=0.25$\,m per action, one forward step, converts length into actions; turns are not charged. $\sigma_j$ is the service time of node $j$: for a room, 4 actions plus 4 per live frontier left in it (a frontier cluster of the room whose view would reveal at least 1\,m$^2$ of unknown floor), at most 60, the bound of a room visit; 4 for an opening or a sighted target (a look along three headings), with $\sigma_0=0$; stair nodes are charged 12 actions. The leaf length $\lambda_i$ is non-zero only for stair nodes (\cref{sec:planning-fallback}). Folding $\sigma_j$ into every arc entering $j$ is exact, because under \cref{eq:objective} the service of $j$ is paid with the same survival probability as the arc that reaches it, and, together with the leaf lengths, it preserves the triangle inequality of $d$, since $C_{ij}+C_{jk}-C_{ik}\ge 2\lambda_j/u+\sigma_j\ge 0$ whenever $d_{ij}+d_{jk}\ge d_{ik}$.

% ----------------------------------------------------------------------------
\subsection{Success probabilities from a language model}
\label{sec:planning-oracle}

\textsc{LLMProbabilities} (\cref{alg:llmprobabilities}) asks Qwen2.5-14B-Instruct~\citep{yang2024qwen25} (4-bit quantised, temperature 0, fixed seed, JSON output) for each node's success probability $p_i$: the probability that searching node $i$ next finds $\cstar$, that is, an instance is there and the agent would see it by going and looking. The prompt (\cref{app:prompts}) names the target and gives a house summary $H$ (room types found and their status, unidentified rooms, openings not yet looked into, whether reachable frontier remains, budget used) and one line $\eta_i$ per node. A room line lists its label (possibly tentative or unknown), area, frontier count, scan status and how long ago it was scanned, search effort, up to eight confirmed object classes, and whether the agent stands in it; an opening line names the room it leaves, whether it is a doorway or a gap, and the objects glimpsed through it.

The prompt contains no distance and instructs the model not to reason from distance, travel time or the room the agent stands in, which it must judge exactly like the others; that flag is the only cue to the agent's position that the model receives. This separation is the central design choice of the planner. The language model contributes commonsense about where objects are kept, while distance, which the occupancy grid measures directly, is left to the solver, which already charges every action of travel through \cref{eq:arccost}; a model that also discounted far rooms would charge travel twice, and the trade-off between likelihood and cost would be made implicitly, inside an uncalibrated score, instead of by the objective. In the same spirit, SayCan~\citep{ichter2023do} multiplies a language model's estimate of a skill's usefulness by a separately learned estimate of its feasibility rather than asking one model for both. Without distances the prompt also changes only when the observed content of the home does (times and the budget are coarsened for the same reason), and a byte-identical prompt reuses the previous reply without a new call.

The reply is reasoning-first~\citep{wei2022chain}: before any number, the model states what kind of home this appears to be, what the house summary implies and what stage the search has reached; it then gives the pass verdict and, per node, a short reason followed by an integer score $r_i$ from 0 to 100. Identifiers not in the prompt are discarded, and a reply that scores none of the listed nodes counts as a failure. Scores, clamped to $[0,100]$, are mapped to $p_i=\min(r_i/100,\,1-10^{-6})$, the ceiling keeping every survival probability positive, and a listed node the reply omits receives $p_i=0.01$, below the eligibility floor. Sighted targets are not shown to the model: the detector has already seen the target category there in at least two frames, or once with confidence at least 0.50, which is stronger evidence than any description the model could weigh, so they receive a fixed $p_i=0.85$ without a call. An opening that the current reply has not scored, for instance one folded into its room when the prompt was built, takes its score from an earlier reply, or $p_i=0.10$ if it was never scored. Fixed rules then scale these values, whatever the model wrote, before the eligibility floor and the solver use them, and $p_i$ denotes the scaled value from here on: a gap keeps the fraction $\min(1,\max(0.25,\,a_i/4\,\mathrm{m}^2))$ of its value, where $a_i$ is the unknown floor the camera would see from its threshold looking through it, while a doorway keeps all of it; a room the agent has stood in and left unfinished keeps half of its value, or a fifth once the agent has spent at least 20 actions in it and at most one live frontier is left.

The scores are not normalised (\cref{sec:problem}); the prompt tells the model not to make them sum to 100, and we do not claim that they are calibrated.

A failed call or an unusable reply triggers one schema-repair prompt. If the repaired reply is unusable too, \ours never substitutes a uniform prior, which would silently turn semantic search into geometric search; the agent instead runs the exploration fallback (\cref{sec:planning-fallback}) for 25 actions before asking again, doubling this back-off after each consecutive failure up to 200 actions. A usable reply resets the back-off.

% ----------------------------------------------------------------------------
\subsection{Ordering the nodes by expected time-to-find}
\label{sec:planning-order}

Each decision point yields an instance of the Hamiltonian path problem with probabilistic terminals~\citep{lyu2026rpt}: the complete graph on $V=\{0\}\cup\N$ with the weights of \cref{eq:arccost}, a terminal probability $p_i$ per node and $p_0=0$ for the agent. The visit order is the open path $\pi=(\pi_0{=}0,\pi_1,\dots,\pi_m)$ through all $m=|\N|$ nodes that minimises $J(\pi)$ of \cref{eq:objective}, so, under the estimated probabilities and the independent-terminal model, $J$ is the expected number of actions until the target is found or the list is exhausted. Unlike choosing the single best-scored frontier at each step~\citep{yu2023l3mvn,zhou2023esc}, the objective values each candidate by its effect on the expected cost of the whole remaining search, so a moderately likely room next door can precede a slightly likelier one across the house.

RPT$^*$ is an exact best-first search over states $(i,g,q,A)$: the current vertex, the accumulated expected cost, the survival probability and the visited set, starting from $(0,0,1,\{0\})$. Expanding to an unvisited vertex $j$ yields $(j,\,g+q\,C_{ij},\,q(1-p_j),\,A\cup\{j\})$. States are ordered by $g+h$, where the admissible heuristic $h$ scales a precomputed lower bound on the expected cost of visiting the remaining vertices by the current survival probability, and a state is pruned when another at the same vertex has no larger cost and a superset of its visited vertices. The first state popped that has visited every vertex is certified optimal for the instance. We re-implemented the solver from its published description, correcting a few points of that description (among them the initial state and the cost of the bound table). The bound table costs $O(|V|^3)$ time, but the search is exponential in the worst case, with up to $|V|\,2^{|V|}$ pairs of vertex and visited set. We therefore cap the instance at 12 nodes ($|V|\le 13$), chosen so that exact solves fit within a budget of 2\,s of computation per decision \TODO{report solve times against node count (\cref{tab:runtime}) and how often the cap binds}. If the budget expires, the solver returns the best complete order it has found, or a nearest-neighbour order if it has completed none, without a certificate.

RPT$^*$ requires metric costs, and the solver checks the triangle inequality before searching. Because \cref{eq:arccost} mixes distances measured at two clearances, metricity is not guaranteed: a pair measured on a longer, wider route can exceed the sum of two pairs measured at the body radius. When the check fails, or the solver raises any other error, the implementation falls back to a single draw of the next node with probability proportional to $p_i$, from a generator seeded per episode, and the resulting choice carries no optimality guarantee (\cref{tab:runtime}).

Execution is receding-horizon, as in the search system that \citet{lyu2026rpt} build on the solver: the agent commits to $\pi_1$, discards the rest of the order when that visit ends, and re-solves with fresh costs and probabilities at the next decision point, so that the per-instance optimum (\cref{sec:problem}) follows the nodes, costs and probabilities as they change.

% ----------------------------------------------------------------------------
\subsection{When to re-plan}
\label{sec:planning-replan}

A decision point occurs when the visit of the node in force ends and when a fallback move completes. Target closing (\cref{sec:execution}) preempts the search on a confident sighting; if it releases its candidate, the interrupted visit resumes rather than triggering a new decision. A visit ends when its node is exhausted (no frontier of the room left to walk, a full rotation completed inside the room, the peek complete or its opening resolved on the way, or the room finished by sight on the way in), when its budget runs out (60 actions for a room visit, 30 for a peek), or when its transit fails (no route within 5 actions, over 120 actions in transit, or blocked for 6). While no node is in force, as during fallback exploration, the oracle is also re-queried when a room without a score appears or a room gains a new kind of object, when a new opening appears (at most once every 6 actions) and, in the second pass, every 25 actions of budget, so that the scanned rooms are re-weighed against the time left.

Four mechanisms prevent the oscillation that re-planning on noisy evidence invites. First, the order is committed until the node is released, and the A$^*$ route to its entry point is itself committed (\cref{sec:execution}), so neither is recomputed at every step. Second, rooms and object kinds discovered while a node is in force, in transit or during its visit, are deferred to the end of the visit, where every node is re-valued anyway, so the agent does not abandon a half-finished frontier walk each time a room is split or gains an object. Third, failed nodes are deferred and released ones cooled, except that a visit which merely spent its budget is not cooled: the budget forces a reassessment, not a departure, and the same node may be chosen again. Fourth, in the multi-storey mode an arrival-return guard keeps the agent from re-taking the staircase by which it has just arrived.

% ----------------------------------------------------------------------------
\subsection{Exploration fallback and multiple storeys}
\label{sec:planning-fallback}

\paragraph{Exploration fallback.} When there is nothing to plan over (no eligible node, the oracle backing off, repeatedly refused transits, or an error in the decision), the agent explores rather than idles; the oracle is not queried while no node exists that its answer could act on. The fallback is a ladder of rungs tried in order. The first rung is the most useful reachable exit frontier~\citep{yamauchi1997frontier}, excluding the blind zone and object shadows as for openings, with utility $U(f)=(n_f+s_f)/[(g_f+0.25)(1+|\Delta\psi_f|/\pi)]$ for a cluster of $s_f$ frontier cells at geodesic distance $g_f$ (in metres) and heading error $\Delta\psi_f$, where $n_f$ counts the unknown cells within camera range that a full turn at the frontier would bring into view; an exit that would reveal less than 1\,m$^2$ is demoted, and the current goal is kept unless another has twice its utility. Then come a floor-change rule (multi-storey mode only), demoted and previously retired frontiers, relocation to the farthest reachable known cell at least 1.5\,m away, a downward-tilted rotation that maps the floor around the agent, taken once per position when at least a quarter of the cells within 1.2\,m are unknown and the benchmark has tilt actions, an escape probe that faces the turn heading with the thinnest occupied band (at most 0.4\,m, with at least 0.5\,m of free or unknown floor beyond it) and takes one forward step, and finally an in-place turn, so that the agent never stands still.

\paragraph{Multiple storeys.} In the multi-storey mode each staircase seen on the current storey becomes a stair node, with its entry point where the flight starts on this storey, scored in the same oracle call (which then also summarises the other storeys) and ordered in the same instance. Its leaf length $\lambda_i$ in \cref{eq:arccost} is the flight's length plus an 8\,m floor-change penalty, charged on every arc into and out of the node, so orders that pass through a staircase and return are expensive; no clock decides a floor change. Taking the stairs walks the connector and saves the storey's state, and a storey entered for the first time begins with a fresh warm-up rotation. The staircase just arrived by is withheld for 40 actions and for as long as the new storey still has an unfinished room or opening. The connector geometry is privileged: it is simulator ground truth, revealed only once a staircase is visible (a perfect, visibility-gated stair detector), together with storey metadata (\cref{tab:privileges}). The mode is off on the single-storey Gibson split, where the same privileged geometry serves only to mask seen stairs as obstacles. \TODO{state whether stair traversal is enabled on HM3D and MP3D.}

% ============================================================================
% 07_execution
% ============================================================================
\section{Execution, Local Exploration and Target Closing}
\label{sec:execution}

This section describes how the agent reaches the node chosen in \cref{sec:planning}, how it searches that node once
there, and how a detection of $\cstar$ becomes a STOP. Every layer emits exactly one discrete action per step and
re-reads the agent's pose before each action (\cref{alg:navigate}, \cref{alg:explorelocal} and
\cref{alg:targetclosing}; constants in \cref{tab:hyperparams}).

\subsection{Path Following}
\label{sec:execution-navigate}

\paragraph{Planning and committed routes.}
Routes are planned with A$^*$~\citep{hart1968formal} on the 8-connected 0.1\,m occupancy grid $\M$, with unknown cells
blocked, an admissible (octile) heuristic and a step cost that grows near obstacles. Cells within 0.30\,m of an
obstacle are excluded whenever a route exists at that clearance; otherwise the required clearance is relaxed in
0.05\,m decrements down to the 0.18\,m body radius, which is never crossed. A route found at a clearance above the
body radius is, however, replaced by the route at the body radius when it is more than 1.5 times as long. A route is
\emph{committed}: it is kept
until the goal moves by more than 0.5\,m, the agent strays more than 0.5\,m from it, the remaining route collides with
newly mapped obstacles, or 30 actions pass without 0.10\,m of progress. Re-planning at every step on a growing map
would flip the route from one side of an obstacle to the other and make the agent oscillate.

\paragraph{Discrete-action converter.}
A converter turns the route into actions (\cref{alg:navigate}). On benchmarks with camera tilt, a commanded pitch is
executed first, one 30$^\circ$ tilt per step until the camera is within half a tilt of it, before any turn or step.
The converter then aims at the point 0.5\,m (two steps) beyond the agent's progress along the path, turns toward it
when the heading error exceeds 15$^\circ$ (half a turn, the closest alignment that 30$^\circ$ turns can guarantee),
and otherwise moves forward. Progress is a high-water mark that only advances, and its projection stays on the current
leg while the agent walks within $0.5\tan 15^\circ \approx 0.134$\,m of it (the lookahead times the tangent of half a
turn), so a later leg passing nearby cannot capture the progress and send the agent backwards. Arrival is declared
within $\max(0.25\,\mathrm{m}, r_{\min})$ of the last waypoint. The reach floor $r_{\min} = s/(2\cos(\theta/2)) =
0.129$\,m, for step $s = 0.25$\,m and turn $\theta = 30^\circ$, is the distance inside which a forward step taken
after aligning to within half a turn is no longer guaranteed to close in; it protects coarser action spaces, while the
0.25\,m tolerance is the binding one on all four benchmarks. When the aim point is the goal itself, a step is taken
only if it strictly reduces the distance and the route otherwise counts as walked, so the agent never circles a point
it cannot hit.

\paragraph{Blocked steps and idle steps.}
A forward step that advances less than 0.05\,m along the heading it was issued with counts as blocked, so that sliding
along a wall is not mistaken for progress. The cell one step ahead is then marked occupied, the route is re-planned,
and until a forward step moves again the converter aims one step (0.25\,m) ahead instead of two. A step on which the
converter has nothing left to do is spent on a turn in place, never on STOP, because a STOP ends the episode.

\subsection{Warm-up and Local Visits}
\label{sec:execution-local}

\paragraph{Warm-up.}
Each episode begins with 10 in-place turns of 30$^\circ$, the fewest after which every bearing has been inside the
79$^\circ$ field of view, so that the first decision sees all
around the spawn point rather than a 79$^\circ$ wedge. A target sighted during the rotation starts target closing at
once; if the closing moves the agent more than 0.5\,m, the rotation restarts so that it is completed from one spot.
The spawn position is then recorded as a completed scan, which finishes the spawn room before any room is chosen.

\paragraph{Room visit.}
A room is searched by walking to its frontiers (\cref{alg:explorelocal}). At every action the frontier clusters
credited to the room are ranked by the unknown floor that a full circle of camera rays from each (0.6 to 5\,m, blocked
by walls and by cells the sight ledger has settled) would reveal, plus the cluster's own cells, over the walking
distance to it and the turn needed to face it; a cluster that would reveal less than 1.0\,m$^2$ is not walked to. The
agent walks to the best cluster on a copy of the map in which other rooms count as unknown (the full map if that
leaves no route), keeps a committed goal while it stays passable, near the room and informative, and moves on to the
next once the camera has resolved it. The visit ends when no cluster is left, or at once when a full rotation
completed inside the room since the visit began, such as a full path glance, has finished the room. The whole visit is
bounded by 60 actions.

\paragraph{When a room is finished.}
A room is \emph{finished} when a completed scan point (the spot of the warm-up or of a completed full path glance)
lies inside it; when a scan saw at least half of its cells within the 5\,m depth range through observed free space and
no live frontier remains in it; when no live frontier remains and the agent has stood in it at least 0.5\,m in from
its edge (anywhere in it, if it is too narrow for that), or the camera has either looked into it (one recorded pose
saw at least half of it with a clear line of sight) or walked through it while it is narrow (widest clearance at most
0.9\,m); or when it is a doorless fragment under 3\,m$^2$ with no live frontier. A live frontier is one the room visit
would walk to, revealing at least 1.0\,m$^2$. A room holding a confirmed object that usually shares a room with the
target (a sink, shower or bathtub, for a toilet) is finished only by a completed rotation inside it, because the
target may stand behind the door jamb. A finished room is judged
again if it grows beyond 1.5 times the size at which it was finished.

\paragraph{Opening peek.}
An opening is searched by walking to its threshold and looking along its direction and 30$^\circ$ to either side, the
three headings swept in one direction starting from the side nearer the current yaw. With the 79$^\circ$ field of
view, adjacent headings overlap by 49$^\circ$ and together cover about 140$^\circ$ beyond the opening for as few as
two turns. When the agent stands more than 0.35\,m short of the threshold and the threshold cell is passable, before
this look or once the look has mapped it, the agent walks on to the threshold (at most 8 actions) and looks from there
along five headings, up to 60$^\circ$ to either side of the way it came through. Completing the look retires the
opening. The walk to the threshold has a bound that scales with
its length; on overrun the agent looks from where it stands if it is within 1.5\,m of the opening and retires the
opening otherwise. The look from the threshold is capped at 30 actions. A sighted but unverified target landmark is
visited in the same way, without the walk on, from a 1.5\,m standoff facing it.

\paragraph{Path glances and the sight ledger.}
In transit, at most one brief in-place glance is scheduled on the route in force, at a point up to 6\,m ahead, when a
90$^\circ$ look to one side or a full rotation there would reveal at least 3\,m$^2$ of unknown floor that the forward
walk would not, at a rate of at least 0.5\,m$^2$ per action spent. Side space that one look reveals is thus cleared
without a later dedicated visit. A second kind of glance responds to a detection cut off by the left or right image
edge (confidence at least 0.4, spanning at least 30\% of the image height): the agent turns toward it by at most three
turns, so that the partly seen object and what stands beside it come into view. For the target's category and the
objects that usually share its room this glance is unconditional; for other categories it requires at least 2\,m$^2$
of unknown floor on that side (3\,m$^2$ for a three-turn glance). The sight ledger records every camera pose and settles two kinds of unknown cell so
that they stop producing frontiers: cells within 2.5\,m that camera rays through observed free space have reached from
two distinct poses without resolving them, and enclosed pockets of unknown space of at most 3\,m$^2$ that are
narrower than 0.6\,m or that the agent has stood on. Without it, such
cells would keep drawing the agent back to places it has already looked at. Settled cells are hidden only from the
frontier logic; the planner keeps the true map.

\subsection{Target Closing}
\label{sec:execution-closing}

\paragraph{Why STOP is gated.}
Detections arrive from up to 5\,m away and include false positives, while a STOP ends the episode. On HM3D, MP3D and
RoboTHOR success requires a STOP close to the target (within 0.1\,m of a goal view point on HM3D and MP3D, with the
target visible within 1\,m on RoboTHOR), so a STOP on a false or distant detection forfeits the episode
where a detour would only cost time; on Gibson STOP is not required, but a premature one still forfeits the remaining
budget. \ours therefore never stops on a single sighting. A box of class $\cstar$ that has confidence at least 0.50
and a depth-supported 3-D point at a plausible height, is not a sliver cut by the image border, and lies outside the
remembered rejections described below starts a three-phase routine VERIFY $\rightarrow$ APPROACH $\rightarrow$ INSPECT
(\cref{alg:targetclosing}) that preempts the search on that very frame. While it runs, the occupancy grid $\M$, the
sight ledger and the landmarks $\Ob$ keep updating, but the rooms $\R$, the oracle and the local visits are suspended.
\TODO{confirm that the HM3D, MP3D and RoboTHOR runs use this closing routine}

\paragraph{VERIFY.}
The candidate is \emph{locked} after two consecutive frames whose 3-D detections fall inside a cylinder around the ray
from the agent to the first sighting, of radius 0.5\,m widened by 0.15\,m per metre of range beyond 2\,m; a frame
in which the candidate shows only as a box of the target's class without usable depth, within 30$^\circ$ of its
bearing (a \emph{glimpse}), does not count toward the two but does not break the run either. Once a
candidate exists, associated re-sightings count from confidence 0.30 and are fused into a precision-weighted 3-D
estimate whose noise grows with range. A candidate with a suspect context, in a room whose firm label excludes the
category or at a height outside the category's band, counts at 0.70 of its confidence and needs four consecutive
frames from at least two viewpoints 0.3\,m apart. Meanwhile the agent faces the box (or the glimpse) and steps toward
it while it is farther than 1.25\,m (the terminal radius plus one step); closer than that, it takes the view that
keeps the box in frame. A candidate first seen within 1.3\,m, without an approach behind it, is not verified from that
single close view, unless that view's confidence is at least 0.92 or a box of the target's class with confidence at
least 0.50 was in the frame on two of the six preceding actions: the agent first backs off by up to three steps to a
spot it knows to be clear, faces the remembered position again and verifies from there (at most 20 actions, once per
object). A candidate not locked within 12
actions, not counting the back-off or up to 24 steps toward a glimpse, is released, and so is one that six
consecutive actions of sweeping about its bearing have not brought back into view.

\paragraph{APPROACH.}
The A$^*$ goal lies on a standoff ring of radius 0.65\,m around the estimate: the 1.0\,m terminal radius minus the
0.25\,m arrival tolerance and one map cell, so that arrival leaves the agent within 0.90\,m of the estimated centre.
Ring bearings are tried from the agent's side outwards, in 30$^\circ$ increments up to 90$^\circ$ on either side and
then the far side. Each goal is projected onto the simulator's navigable surface, a privileged input
(\cref{tab:privileges}), and accepted only if A$^*$ reaches within 0.75\,m of the target; when no bearing has such a
route on the known floor, the ring is tried again with unknown cells passable at eight times the cost of open floor.
The route persists through
occlusion, and a refined estimate keeps the current goal while that goal lies within 1.0\,m of the refined estimate,
so that a small refinement cannot trade a reachable goal for an unreachable one. The lock is released if the last 20
approach actions stayed inside a circle of radius 0.20\,m, or if no safe path exists for six consecutive actions. On
benchmarks with tilt, a lock that has gone two actions without a path first maps the floor around the agent, since the
blind disk under a level camera is the usual reason for a missing path: when at least a quarter of the cells within
1.2\,m are still unknown, the camera tilts down, the agent turns until the camera has seen the full circle (10 turns)
and the camera tilts back (at most 16 actions, cut short once a path exists), and the six pathless actions are counted after this sweep. A release for want
of a path remembers the agent's own spot rather than the target and refuses new candidates from within 1.0\,m of it.

\paragraph{INSPECT.}
The inspection begins when the standoff is reached, or earlier as soon as a fresh measured range is within 1.05\,m; if
a fresh reading at the standoff is still beyond 1.05\,m, the approach resumes toward a closer goal, except that a fresh
box that is off centre but within 30$^\circ$ of the heading, measured less than one step (0.25\,m) beyond 1.05\,m after
an in-range sighting earlier in the same inspection, is closed by a single forward step (at most three per
inspection). Otherwise the agent makes no forward step and turns to centre the fresh box. Its measured range is the smaller of the planar range to the fresh
centroid and the near-edge range, the 10th percentile of the horizontal distances of the box's depth points. STOP is
issued on a fresh box that is centred (it spans the image's centre column, lies within 15$^\circ$ of the heading, or
lies within 30$^\circ$ of it just after a centring turn carried it past the centre) at
a measured range of at most 1.05\,m: the 1.0\,m terminal radius, matched to Gibson's 1\,m success region and measured
to the object's near surface rather than to its centre, plus 0.05\,m of projection tolerance; after an
\emph{extended approach} (at least 3\,m walked and 1\,m of range closed) a fresh box within 1.05\,m suffices even
when it is not centred. A fresh box below the 0.50 threshold is first given at most two further frames, from views
that keep it whole, in which the detector may read it more strongly; STOP then follows on the evidence available.
Missed detections cycle through at most nine views (three yaw offsets times three pitches; three views without tilt),
and the inspection is bounded by 24 actions. Without such a fresh box, STOP is allowed only on recent remembered
evidence, in five cases: (i)~after an extended approach, when the estimate is within reach (1.05\,m, or 1.30\,m once
the standoff route has been walked) and three consecutive frames show no box; (ii)~when the inspection has run out of
views, the estimate is within reach, and either the approach was extended or a fresh box of the target's supporting
surface lies beneath the estimate; (iii)~when the inspection has run out after the target was seen fresh within range
earlier in the same inspection; (iv)~on benchmarks without tilt, when the target has slipped below the frame
(described below); and (v)~on benchmarks with tilt, on the frame after the centring tilt described below, if the
detector has lost the box and the estimate is within 1.05\,m. Otherwise the inspection fails and the candidate is
released. Whichever of these rules decides a STOP, including the centring-tilt rule below, the STOP is issued only
after an overshoot that consults no further frame: one forward step toward the estimate, and a second if the estimate
is then still beyond 0.90\,m, each taken only if afterwards the estimate would still lie at least 0.45\,m ahead and the
near edge measured on the deciding frame at least the 0.18\,m body radius, and only while the estimate lies within
90$^\circ$ of the heading; a blocked step ends the overshoot at once, and it is bounded by five actions, turns
included. Outside the closing, the agent also stops when a target landmark confirmed from three
views (one of them at least 0.2\,m or 20$^\circ$ from the first) lies within 0.75\,m and the closing has not rejected
it.

\paragraph{Release and remember.}
A released candidate's spot is remembered as a rejection disk of radius 1.0\,m, doubled to 2.0\,m when an inspection
disproved it, when it went unverified a second time, or when it was never seen again after a back-off (no disk at all
in this last case if one of its frames had a confidence of at least 0.92). New candidates
inside a disk are refused while they are farther than 2\,m, since a closer view is new evidence, and all candidates
farther than 2\,m are refused during a 15-action cooldown after a release; a confidence of at least 0.80 overrides the
cooldown and the 1.0\,m disk of a single unverified release, but not a 2.0\,m disk. The search then resumes the
interrupted task. There is no per-episode cap on the number of closings: the rejection memory and the cooldown are
what keep the same false positive from being retried. The whole closing is bounded by 160 actions, beyond which the
episode ends as a failure rather than with a STOP.

\paragraph{Pitch with and without tilt.}
On benchmarks with 30$^\circ$ look-up and look-down actions, an image-space servo owns the pitch whenever a fresh box
lies within 1.5\,m; without a fresh box the pitch follows the target's estimated elevation. A box cut by the bottom
margin, or whole but centred in the lowest 25\% of the rows, is tilted down by the rounded number of steps that best
centres it without pushing its other edge out of the frame, and symmetrically upwards; a box cut by both margins is
left alone. On the approach the servo reads the box predicted after the next forward step rather than the current one.
The central band is never narrower than one tilt's rounding band, $f_y\tan 15^\circ$ rows on either side of the
principal point with $f_y$ the vertical focal length in pixels, so the servo cannot hunt between two pitches; for the
640$\times$480, 79$^\circ$ camera this is 104 rows, and the 25\% regions are the binding limit. Each inspection allows
one centring tilt before STOP, after which the agent stops on the next frame, or on the previous frame's evidence if
the detector loses the box. \TODO{report the effect of the servo and the centring tilt on a benchmark with tilt}

Without tilt actions, as in Gibson's published four-action protocol, a level camera 0.88\,m high sees the floor only
from about 1.4\,m out, so low targets leave the image during the last steps of an approach. Three rules replace the
servo there. A verification
step is taken only if the box is predicted to stay in frame afterwards. A frame whose box is cut by the lower image
edge while the target's base is out of view counts as a sighting but is not fused into the estimate: its centroid is
that of the visible upper part, and fusing it would drift the estimate, and the standoff measured from it, toward that
part. Finally, rule (iv) issues STOP without a fresh box when the target's predicted visible part spans less than 15\% of the
image height, its bearing lies inside the field of view, its measured top is below the camera, the estimate is within
1.05\,m, and a fresh box was seen within the last six actions. Our Gibson runs add the look-up and look-down actions to
that protocol, so the first and third rules are inactive on all four benchmarks as we run them.

% ============================================================================
% 08_experiments
% ============================================================================
% ============================================================================
%  08_experiments.tex -- Experimental setup
% ============================================================================
\section{Experimental Setup}\label{sec:experiments}

We evaluate \ours on four datasets and five validation splits that span two simulators, three definitions of
success and goal vocabularies of 6 to 21 categories. This section fixes the protocols, metrics, baselines,
privileged inputs and implementation; \cref{app:protocols} gives the protocol details and \cref{app:hyperparams}
all hyper-parameters.

\subsection{Benchmarks and Protocols}\label{sec:experiments-benchmarks}

\Cref{tab:benchmarks} summarises the five splits. A single evaluation harness implements each benchmark's own
success rule and metrics. Because every evaluator is a reimplementation, the harness also checks itself: after every
action it verifies that the simulator executed the declared motion, the path length it accumulates must agree within
1\,mm with the one the benchmark environment reports, and any metric the simulator computes natively must agree with
the harness's value. On the Habitat benchmarks~\citep{savva2019habitat} the agent is a cylinder of radius 0.18\,m
whose camera, 0.88\,m above the floor, delivers $640\times480$ RGB-D images with a 79$^\circ$ horizontal field of
view and a 0.5--5.0\,m depth range; a step moves 0.25\,m, a turn rotates 30$^\circ$, and an episode ends at STOP or
after 500 actions.

\begin{table}[t]
  \centering
  \footnotesize
  \setlength{\tabcolsep}{4pt}
  \caption{Benchmarks. All five are validation splits with a budget of 500 actions. Actions: S = STOP,
  F = forward 0.25\,m, T = turn left/right 30$^\circ$, U/D = camera tilt up/down 30$^\circ$. Evaluator: the
  reference whose definitions our harness reimplements; HM3D v1/v2 and MP3D are scored on the navigation mesh
  recomputed for the ObjectNav agent.}
  \label{tab:benchmarks}
  \begin{tabular}{@{}p{1.45cm}p{2.25cm}rrrp{3.85cm}p{2.0cm}p{1.75cm}@{}}
    \toprule
    Benchmark & Split & Episodes & Scenes & Cat. & Success rule & Actions & Evaluator \\
    \midrule
    Gibson & SemExp v1.1, tiny val & 1000 & 5 & 6 & FMM distance to the target's ground-truth map dilated by 1\,m
      is 0 at episode end; STOP not required; no visibility test & S, F, T, U/D (U/D added by us); sliding on & SemExp (FMM) \\
    HM3D v1 & 2022 challenge val (HM3D-Sem v0.1) & 2000 & 20 & 6 & STOP and geodesic distance $<0.1$\,m to the
      nearest goal view point & S, F, T, U/D & habitat-lab \\
    HM3D v2 & val (HM3D-Sem v0.2) & 1000 & 36 & 6 & as HM3D v1 & S, F, T, U/D & habitat-lab \\
    MP3D & 2021 challenge val & 2195 & 11 & 21 & as HM3D v1 & S, F, T, U/D & habitat-lab \\
    RoboTHOR & 2021 challenge val & 1800 & 15 & 12 & STOP while the target is visible within 1\,m; no distance
      test & S, F, T, U/D (pitch $\pm30^\circ$); actuation noise & RoboTHOR challenge \\
    \bottomrule
  \end{tabular}
\end{table}

\paragraph{Gibson.} The SemExp split~\citep{chaplot2020object} (v1.1, tiny partition) of Gibson~\citep{xia2018gibson}
has 1,000 validation episodes, 200 in each of five buildings (Collierville, Corozal, Darden, Markleeville and
Wiconisco), with every start and goal on one storey, over six categories (chair, couch, potted plant, bed, toilet,
tv). SemExp's action set has no camera tilt; we add Habitat's look-up and look-down actions, with the camera starting
level and kept between 30$^\circ$ above and 60$^\circ$ below level, a deviation from the SemExp protocol
(\cref{tab:privileges}). Sliding along walls is enabled.
Success follows SemExp's evaluator, which we reimplement: at the end of the episode, the Fast
Marching~\citep{sethian1996fast} distance on a 5\,cm grid from the agent to the target category's ground-truth
semantic map dilated by 1\,m must be zero. STOP is not required, so an episode that exhausts its budget inside this
region also succeeds, and visibility is not tested; the rule as stated by \citet{loo2025open} adds a visibility
condition, which SemExp's released evaluator does not check. \ours does not wait for the budget: it issues STOP when
its terminal rule (\cref{sec:execution}) holds, after at most two further forward steps toward the target, a rule chosen to sit
inside the 1\,m region \TODO{report how many of
our Gibson successes nevertheless end at the budget without STOP}. As in SemExp's code, the reference length is the
start FMM distance plus 1\,m, and path length is measured in the horizontal plane. We run habitat-sim 0.2.4,
whereas SemExp's reference used 0.1.5, and load the published navigation mesh as it is; the mesh governs motion and collisions but not the score, which is computed on the semantic
map \TODO{state whether SemExp's habitat-sim 0.1.5 recomputed the Gibson navigation mesh for the 0.18\,m agent}.

\paragraph{HM3D.} HM3D v1 is the validation split of the Habitat 2022 ObjectNav challenge on HM3D-Semantics
v0.1~\citep{ramakrishnan2021habitat,yadav2023habitat}, with 2,000 episodes in 20 buildings. HM3D v2 has 1,000
validation episodes on HM3D-Semantics v0.2 in 36 buildings and keeps the same agent, as the papers we compare with do,
rather than the 2023 challenge's Stretch robot. Both splits have six categories, tilt actions and no sliding, and
HM3D v1 includes episodes whose goal lies on another floor. Following habitat-lab~\citep{szot2021habitat}, success
requires STOP within a geodesic distance of 0.1\,m of the nearest goal view point, one of the navigable points placed
within about 1\,m of a goal instance; reference length and final distance are geodesic distances to these points,
and path length is measured in 3-D. The shipped navigation meshes are built for a 0.10\,m radius and 1.50\,m height;
habitat-sim silently recomputes them for the 0.18\,m radius and 0.88\,m height that habitat-lab configures for the
ObjectNav agent, so every published distance and score on these splits comes from the recomputed mesh. We evaluate
on it explicitly; with it our start distances reproduce the geodesic distances published with the HM3D episodes
exactly, whereas the shipped mesh gives distances 5.5\% short.

\paragraph{MP3D.} The Habitat 2021 ObjectNav validation split of Matterport3D~\citep{chang2017matterport3d} has 2,195
episodes in 11 buildings with 21 goal categories. Agent, actions, success rule, metrics and navigation mesh are as for
HM3D, scored by the same reimplementation of habitat-lab's ObjectNav measures. Both HM3D and MP3D contain multi-storey
buildings: an agent confined to its start storey cannot solve the 14.6\% of HM3D v1 and 9.6\% of MP3D episodes that
\citet{yokoyama2024vlfm} report as requiring stairs \TODO{state whether stair traversal is enabled on HM3D and MP3D}.

\paragraph{RoboTHOR.} The RoboTHOR ObjectNav 2021 validation split~\citep{deitke2020robothor} in
AI2-THOR~\citep{kolve2017ai2} has 1,800 episodes, 120 in each of 15 apartments, over 12 categories, many of them small
objects such as mugs, apples and spray bottles. The LoCoBot agent (radius 0.175\,m) carries a $640\times480$ RGB-D
camera 0.8688\,m above the floor with a 79$^\circ$ horizontal field of view and a 0.1--5\,m depth range. Rotations
and steps are perturbed by Gaussian actuation noise (standard deviations of 0.5$^\circ$ per rotation and 5\,mm per
step), the camera pitch is limited to $\pm30^\circ$ and starts every episode at 30$^\circ$ down, and a failed action
still consumes a step. Success requires STOP while the target object
is visible within 1\,m; no distance is computed. \SPL uses the reference lengths shipped with the episodes, 6.0\% of
which are zero, so such an episode scores \SPL~0 whenever the agent moves, an upstream artefact we keep. The official
visibility check covers only the first-listed instance of the category; we record both that verdict and an
any-instance verdict \TODO{state which verdict the RoboTHOR results report}. The official runner gives the agent only
the goal category and the RGB and depth images; our runs add the simulator pose (\cref{tab:privileges}).

\subsection{Metrics}\label{sec:experiments-metrics}

For episode $i$ of $n$, let $S_i\in\{0,1\}$ be the benchmark's success verdict (\cref{eq:success}), $\ell_i$ the
benchmark's reference path length, $L_i$ the
length of the path travelled (turns, tilts and STOP add nothing), and $d_{0,i}$, $d_{\mathrm{end},i}$ the distances to
the goal at the start and at the end of the episode. Following~\citet{anderson2018evaluation},
\begin{equation}
  \mathrm{SR} = \frac{1}{n}\sum_{i=1}^{n} S_i,
  \qquad
  \mathrm{SPL} = \frac{1}{n}\sum_{i=1}^{n} S_i\,\frac{\ell_i}{\max(\ell_i,\,L_i)}.
  \label{eq:spl}
\end{equation}
\SoftSPL~\citep{datta2021integrating} credits partial progress and is gated on neither success nor STOP:
\begin{equation}
  \mathrm{SoftSPL} = \frac{1}{n}\sum_{i=1}^{n} \max\!\left(0,\, 1-\frac{d_{\mathrm{end},i}}{d_{0,i}}\right)\,
  \frac{\ell_i}{\max(\ell_i,\,L_i)}.
  \label{eq:softspl}
\end{equation}
habitat-lab writes $d_{0,i}$ for $\ell_i$ in the efficiency ratio. The two coincide on HM3D and MP3D; on Gibson
$\ell_i=d_{0,i}+1$\,m, and our form keeps the ratio identical to that of \SPL. We know of no published \SoftSPL on
Gibson or RoboTHOR, nor of a published \DTG on RoboTHOR, so these serve as supplementary figures. The ratio
$\ell_i/\max(\ell_i,L_i)$ is taken as 1 when $\ell_i=L_i=0$, as in RoboTHOR's evaluator. Distance to goal is averaged
over the episodes whose final position can still reach a goal, and the others are counted separately:
\begin{equation}
  \mathrm{DTG} = \frac{1}{|\mathcal{F}|}\sum_{i\in\mathcal{F}} d_{\mathrm{end},i},
  \qquad \mathcal{F}=\{\,i : d_{\mathrm{end},i}<\infty\,\}.
  \label{eq:dtg}
\end{equation}
The distance is geodesic to the nearest view point on HM3D and MP3D and to the nearest instance on RoboTHOR. On Gibson
it is the FMM distance to the 1\,m success region, and unreachable cells receive a finite sentinel, so every episode
enters the mean; Gibson \DTG is therefore a distance to the success boundary. The Gibson literature mixes this quantity,
often printed as DTS $=\max(d-1\,\mathrm{m},0)$ with $d$ the distance to the goal, with distances to the object itself,
so we compare distances only between rows that share a definition. We give 95\% Wilson intervals for \SR and seeded
percentile-bootstrap intervals for \SPL, and compare variants run on the same episodes (\cref{sec:ablations}) with an
exact McNemar test on success and a sign-flip permutation test on \SPL.

\subsection{Baselines}\label{sec:experiments-baselines}

We compare with published results, grouped by the navigation training they use. \emph{Trained in domain}:
SemExp~\citep{chaplot2020object}, PONI~\citep{ramakrishnan2022poni}, 3D-aware ObjectNav~\citep{zhang20233d},
T-Diff~\citep{yu2024trajectory}, GOAL~\citep{li2025goal}, Habitat-Web~\citep{ramrakhya2022habitat}, PIRLNav~\citep{ramrakhya2023pirlnav},
OVRL-v2~\citep{yadav2023ovrl} and EmbCLIP~\citep{khandelwal2022simple}, all trained on ObjectNav, and
SGM~\citep{zhang2024imagine}, whose map model is trained self-supervised on the benchmark's training scenes; these
indicate what in-domain training buys rather than serve as like-for-like competitors. \emph{Zero-shot, trained on
another task}: ZSON~\citep{majumdar2022zson} and PSL~\citep{sun2024prioritized} (ImageNav),
PixNav~\citep{cai2024bridging} (pixel-goal navigation) and VLFM~\citep{yokoyama2024vlfm} (a PointNav policy trained on
HM3D). \emph{Training-free}: frontier-based exploration (FBE)~\citep{yamauchi1997frontier}, CoW~\citep{gadre2023cows},
GoW, the GLIP-based variant of CoW implemented by ESC~\citep{zhou2023esc}, ESC itself, L3MVN~\citep{yu2023l3mvn}, LGX~\citep{dorbala2024can}, LFG~\citep{shah2023navigation},
SemUtil~\citep{chen2023how}, VoroNav~\citep{wu2024voronav}, OpenFMNav~\citep{kuang2024openfmnav},
InstructNav~\citep{long2024instructnav}, TriHelper~\citep{zhang2024trihelper}, GAMap~\citep{yuan2024gamap},
SG-Nav~\citep{yin2024sg}, UniGoal~\citep{yin2025unigoal}, TopV-Nav~\citep{zhong2024topvnav},
ApexNav~\citep{zhang2025apexnav}, WMNav~\citep{nie2025wmnav}, BeliefMapNav~\citep{zhou2025beliefmapnav},
CogNav~\citep{cao2025cognav}, OSG-Nav~\citep{loo2025open}, the method of \citet{unlu2024reproducibility}, VLTNet~\citep{vltnet2024},
RF-Nav~\citep{yuan2024defa}, PIGEON~\citep{peng2025pigeon} and HRO~\citep{jia2026hro}. L3MVN and TriHelper, although training-free in their navigation, use a
semantic segmenter fine-tuned in domain, and PONI's implementation of FBE a detector fine-tuned on Gibson training
images; we mark these rows. Each benchmark's table (\cref{tab:gibson,tab:hm3d,tab:mp3d_robothor}) lists the verified
results on that split, apart from those excluded for the reasons stated in its notes.

We do not re-run baselines. Numbers are quoted at their printed precision from the original papers or, where a table
note says so, from the authors' code release (SemExp on Gibson) or the official leaderboard (EmbCLIP on RoboTHOR),
with fractions converted exactly to percentages. Where a later paper reprints a different value for the same method
we use the original. Figures that exist only as third-party evaluations are marked; among them are OSG-Nav's
re-implementations of LGX and LFG on Gibson, ApexNav's re-evaluations on HM3D v2, and SG-Nav's runs of L3MVN and
OpenFMNav on MP3D and RoboTHOR and of VLFM on RoboTHOR. Five further caveats apply, detailed in
\cref{app:protocols}. (i)~\emph{Success radius.} ApexNav (with its HM3D v2 re-evaluations) and SG-Nav's released MP3D
configuration use 0.2\,m, UniGoal 1.6\,m and InstructNav's code 0.25\,m, and some papers state none or an unclear one;
a looser radius can only raise \SR and \SPL, so we mark these rows and also re-score our HM3D episodes at 0.2\,m
\TODO{add the 0.2\,m re-scoring for MP3D, or state that it is not reported}. (ii)~\emph{Episode sets.} ZSON's
Gibson result (another episode set) and OSG-Nav's 400-episode HM3D rows are excluded; subsets and unstated episode
counts are marked; for SemExp on MP3D we use the third-party run on the full split that SGM~\citep{zhang2024imagine}
prints. (iii)~\emph{Test splits.} ProcTHOR~\citep{deitke2022procthor} and EmbCLIP test-split figures are excluded;
EmbCLIP's validation leaderboard entry and PIRLNav's validation result are our trained references.
(iv)~\emph{Reproducibility.} ESC's RoboTHOR result (38.1 \SR, 22.2 \SPL) was not reproduced by a later paper that
re-ran it, which obtained 34.5 and 18.2~\citep{unlu2024reproducibility}; we list both. (v)~\emph{Inputs and
actions.} ZSON, PSL, PixNav and OVRL-v2 use RGB without depth, UniGoal and CoW act without look-up and look-down
actions, and WMNav acts through at most 40 polar moves proposed by a vision-language model; we mark these rows
\TODO{check whether PIRLNav and EmbCLIP use RGB without depth and, if so, mark them as well}.

\subsection{Privileged Inputs, Deviations and Development Hygiene}\label{sec:experiments-privileges}

\Cref{tab:privileges} lists every input or assumption in which our evaluation departs from the bare protocol or from
the typical baseline. \ours is training-free and builds a metric map. No goal coordinates, distance fields or evaluator
signals reach the policy; as a safeguard, the harness rejects episode metadata whose fields name goal positions,
shortest paths or geodesic distances.

\begin{table}[tp]
  \centering
  \footnotesize
  \setlength{\tabcolsep}{4pt}
  \caption{Privileged inputs and protocol deviations of our evaluation. ``Not reported'' means that the compared
  papers do not describe using the input.}
  \label{tab:privileges}
  \begin{tabular}{@{}p{2.55cm}p{3.6cm}p{3.35cm}p{4.6cm}@{}}
    \toprule
    Input / assumption & \ours & Typical baselines & Note \\
    \midrule
    Agent pose & simulator ground truth & ground-truth pose (GPS+compass) for most Habitat methods; none for
      RGB-only policies & deviation on RoboTHOR, whose official runner provides no pose; ESC runs there without it \\
    Depth & noise-free simulator depth; enclosed holes filled & noise-free simulator depth; zero-depth pixels dropped
      & deviation; a hole is filled only if valid depth encloses $\geq$50\% of its rim and it covers $\leq$50\% of
      the frame \\
    Object semantics & predicted by an open-vocabulary detector & predicted; some segmenters fine-tuned in domain & no
      ground-truth semantics \\
    Navigation training & none (training-free \cmark) & none to large-scale RL/IL & see \cref{sec:experiments-baselines} \\
    Map & metric occupancy grid & metric maps in modular methods; none in end-to-end policies; OSG-Nav non-metric
      & \\
    Foundation models & open-weight, run locally (3B and 14B LLMs) & often proprietary API models & \\
    Stair geometry & ground-truth connector geometry from the navigation mesh, revealed only when visible, plus
      storey (floor-level) metadata in the multi-storey mode & not reported & on single-storey Gibson used only to
      mask seen staircases as obstacles \\
    Depth range handling & below-minimum readings invalid, at/beyond-maximum read as no return & clamped into the
      valid range (habitat-lab) & changes perception, not scoring \\
    Target standoff goal & projected onto the navigable mesh & not reported & closing routine only; rejected beyond
      0.5\,m horizontally or 0.2\,m vertically \\
    Goal location, distances, evaluator signals & none & none & metadata check in the harness \\
    Action space & Gibson: S, F, T plus look-up/look-down (30$^\circ$ steps, 30$^\circ$ up to 60$^\circ$ down) &
      Gibson (SemExp): S, F, T only & deviation on Gibson; a tilt uses one action of the budget and moves nothing \\
    STOP for success & not required on Gibson; \ours issues STOP & same protocol per benchmark & Gibson (SemExp)
      convention \\
    Simulator & habitat-sim 0.2.4 (Habitat benchmarks) & SemExp reference: 0.1.5 (Gibson) & on Gibson the published
      navigation mesh is loaded as it is \\
    Success radius & official 0.1\,m (HM3D, MP3D) & mostly 0.1\,m; 0.2\,m in ApexNav (HM3D, MP3D) and SG-Nav (MP3D)
      & our HM3D episodes also re-scored at 0.2\,m; other radii in \cref{sec:experiments-baselines} \\
    RoboTHOR visibility & official first-listed-instance verdict and an any-instance verdict, both recorded &
      official verdict & see \cref{sec:experiments-benchmarks} \\
    \bottomrule
  \end{tabular}
\end{table}

Development used Gibson \emph{training} buildings, with start positions generated by SemExp's training criteria, but
the validation split was not left untouched. During development we ran and inspected the first three episodes of
each validation scene and corrected faults that they exposed, the depth-hole filling was designed on the
spawn-view depth frames of one validation episode, and two component choices were compared on validation episodes:
the detector checkpoint size, on the first episode of each validation scene, and the default local explorer, on two
validation episodes besides fifteen training buildings. After a full validation run we read its aggregate and
per-category results, examined nine of its failed episodes and revised the target-closing routine in response; the
revisions comprise the low-target terminal rules for the then tilt-less Gibson protocol, the rule that keeps a reachable
standoff goal, the confidence-gated STOP with its weak-box holds, and the image-space pitch servo with its centring
tilt (\cref{sec:execution}), and the optional mask-based depth, off by default, was also added after that run
\TODO{confirm the complete list of closing changes made after the validation run}. A later campaign used 15 further
validation episodes with the toilet goal, three per scene drawn by a seeded random sample (later runs kept two of each
three): from the per-step traces of 18 runs on them and about ten re-runs of one or two of them, and by the \SR and \SPL of
these runs, we changed rules and values throughout the search, perception, mapping and target-closing stages, chose the steps taken after a decided
STOP between variants run on these episodes, and added the camera tilt on Gibson. Our Gibson results
(\cref{sec:results}) therefore come from a split that informed the search, perception, mapping, the closing routine and the
action space and are not a held-out estimate; we report on this split because it is
the benchmark's evaluation set, as every method in \cref{tab:gibson} does. \TODO{frozen-evaluation statement: the date the configuration was frozen, confirmation that each reported run was
executed once with it and no later change, and whether any HM3D, MP3D or RoboTHOR episode was inspected during
development}. Each run records hashes of the source tree, the episode data, the detector checkpoint and its prompt
features, the language-model digests and the seed. The harness refuses to resume an interrupted run whose
configuration differs; the only changes it can be told to accept concern three named items (the motion-check
tolerances, the source-tree hash and the placement of services on the GPU), and any such change is recorded with the
run.

\subsection{Implementation Details}\label{sec:experiments-implementation}

Detector settings, vocabulary and thresholds are those of \cref{sec:perception} and \cref{tab:hyperparams}; the
optional Grounding DINO, BLIP-2 and SAM modules are off in the main configuration. Rooms are labelled by Qwen2.5-3B-Instruct and node probabilities are estimated by
Qwen2.5-14B-Instruct~\citep{yang2024qwen25}, quantised to 4 bits (about 9\,GB)
\TODO{quantisation of the 3B model and Ollama version}. Both are served locally by Ollama and kept resident, with
temperature 0, seed 17 and JSON output. Oracle calls have an 8192-token context and at most 1536 output tokens, label
calls at most 768 output tokens; they time out after 900\,s and 120\,s respectively. A timed-out request is not
retried, a reply that violates the oracle's output schema gets one repair request, and a failed oracle call triggers
the exploration fallback (\cref{sec:planning}). The Habitat benchmarks run on habitat-sim 0.2.4 and RoboTHOR on
AI2-THOR \TODO{AI2-THOR version and Unity build; the challenge reference is AI2-THOR 2.7.2}.

The main configuration is a workstation whose single NVIDIA RTX 4090 (24\,GB) holds the LLM server, the detector and
the renderer \TODO{workstation CPU model and RAM}. A reduced variant runs on a laptop with an 8\,GB RTX 5070 Laptop GPU
(Intel i9-14900HX, 30\,GB RAM), where the 14B model does not fit beside the other services and Qwen2.5-3B-Instruct
serves both roles; it is reported separately (\cref{tab:llm}). Detection and mapping run at every action, the oracle
once per decision point (an unchanged prompt reuses the previous reply), and RPT$^*$ orders at most 12 nodes per
instance. \TODO{per-step and per-decision latency on both machines, oracle calls per episode, wall-clock time per
episode, and the number of runs per benchmark}.

% ============================================================================
% 09_results
% ============================================================================
\section{Results}\label{sec:results}

We compare \ours with published zero-shot and trained ObjectNav agents on Gibson, HM3D (v1 and v2), MP3D and
RoboTHOR, under the protocols of \cref{sec:experiments} and \cref{tab:benchmarks}. Baseline numbers are quoted as
described in \cref{sec:experiments}. Protocols differ in success radius, episode subset,
observations, actions and perception training, so each table carries a regime column and notes for the deviations we
could verify, and differences of a few points between rows are not evidence on their own. \ours is training-free, but
it receives ground-truth pose and noise-free depth with hole filling, it projects the closing standoff goal onto
the navigable mesh, and on Gibson it adds two camera-tilt actions to the four actions the baselines use
(\cref{tab:privileges}); its rows should be weighed with these inputs in mind.

\paragraph{Gibson.}
\Cref{tab:gibson} reports the SemExp tiny validation split (1{,}000 single-storey episodes in five scenes, six
categories), scored with SemExp's evaluator, under which success means ending within the 1\,m-dilated goal region
by fast-marching distance~\citep{sethian1996fast} and STOP is not required (\cref{sec:experiments}). For
SemExp~\citep{chaplot2020object} we print the figure that the authors' code release states for their released model
on the v1.1 episode file, the file and evaluator we use, rather than the paper's own figure, which was measured on an
earlier episode set. The distance column needs care. Our \DTG, like SemExp's evaluator, is the distance to the 1\,m
success region, which several papers call DTS, whereas SemUtil and L3MVN, and OSG-Nav for its own runs, define the
distance to the goal object; we therefore compare distances only among rows that share our definition. The strongest
zero-shot rows rely on GPT-4V~\citep{yuan2024gamap}, on a PointNav policy trained on HM3D~\citep{yokoyama2024vlfm},
on a segmenter fine-tuned in domain~\citep{zhang2024trihelper} or, for HRO, on a Mask R-CNN detector of unstated
training, whereas \ours uses two open-weight Qwen2.5
models~\citep{yang2024qwen25} served locally and a detector that is not fine-tuned. \ours reaches \TODO{SR},
\TODO{SPL} and \TODO{DTG}; \TODO{state its position relative to the best training-free row, to the training-free
rows that share our distance definition, and to the trained group, claiming a difference only where our 95\%
interval excludes the published figure}. As disclosed in \cref{sec:experiments}, validation episodes were inspected
during development and the target-closing routine was revised after inspecting failures of an earlier validation run;
the later revisions of the search, the perception, the mapping, the closing and the action set were diagnosed on a seeded sample of 15
validation episodes with the toilet target and kept or reverted by their \SR and \SPL on them; \TODO{final frozen-evaluation statement for the reported Gibson
number}.

\begin{table}[t]
\centering
\small
\caption{Gibson ObjectNav, SemExp v1.1 tiny validation split (1{,}000 episodes). \SR and \SPL in \%, \DTG in
metres; higher \SR and \SPL and lower \DTG are better. Baseline numbers are as printed in the cited papers or,
where the note says so, by a third party or the authors' code release; -- means not reported.}
\label{tab:gibson}
\begin{tabular}{@{}llccc@{}}
\toprule
Method & Regime & \SR & \SPL & \DTG \\
\midrule
\multicolumn{5}{@{}l}{\emph{Trained}} \\
SemExp, released v1.1 model~\citep{chaplot2020object} & Trained & 65.7 & 33.9 & 1.474\textsuperscript{a} \\
PONI~\citep{ramakrishnan2022poni} & Trained & 73.6 & 41.0 & 1.25\textsuperscript{a} \\
3D-aware ObjectNav~\citep{zhang20233d} & Trained & 74.5 & 42.1 & 1.16\textsuperscript{a} \\
SGM~\citep{zhang2024imagine} & Self-sup. & 78.0 & 44.0 & 1.11\textsuperscript{a} \\
T-Diff~\citep{yu2024trajectory} & Trained & 79.6 & 44.9 & 1.00\textsuperscript{a} \\
GOAL~\citep{li2025goal} & Trained & 83.5 & 44.2 & 0.83\textsuperscript{a} \\
\midrule
\multicolumn{5}{@{}l}{\emph{Zero-shot}} \\
LGX\textsuperscript{c}~\citep{dorbala2024can} & TF & 31.0 & 5.2 & 4.775\textsuperscript{b} \\
FBE\textsuperscript{c}~\citep{yamauchi1997frontier} & TF$^{\ast}$ & 64.3 & 28.3 & 1.78\textsuperscript{a} \\
LFG\textsuperscript{c}~\citep{shah2023navigation} & TF & 64.5 & 40.6 & 1.812\textsuperscript{b} \\
SemUtil~\citep{chen2023how} & TF & 69.3 & 40.5 & 1.488\textsuperscript{b} \\
OSG-Nav~\citep{loo2025open} & TF & 73.4 & 38.6 & 1.722\textsuperscript{b} \\
L3MVN~\citep{yu2023l3mvn} & TF$^{\ast}$ & 76.1 & 37.7 & 1.101\textsuperscript{b} \\
VLFM value map + FMM\textsuperscript{c,d}~\citep{yokoyama2024vlfm} & TF & 82.8 & 48.5 & -- \\
VLFM~\citep{yokoyama2024vlfm} & Other task & 84.0 & 52.2 & -- \\
HRO\textsuperscript{e}~\citep{jia2026hro} & TF & 84.0 & 46.5 & -- \\
TriHelper~\citep{zhang2024trihelper} & TF$^{\ast}$ & 85.2 & 43.1 & 0.600\textsuperscript{a} \\
GAMap\textsuperscript{d}~\citep{yuan2024gamap} & TF & 85.7 & 55.5 & -- \\
\midrule
\textbf{\ours (ours)} & TF & \tbd & \tbd & \tbd \\
\bottomrule
\end{tabular}
\par\smallskip
\parbox{\linewidth}{\footnotesize Regime: \emph{Trained}, trained on ObjectNav; \emph{Self-sup.}, map model
trained self-supervised on in-domain training scenes; \emph{Other task}, zero-shot with a navigation policy
trained on another task; \emph{TF}, training-free; \emph{TF$^{\ast}$}, training-free navigation with a segmenter
fine-tuned on in-domain images. \textsuperscript{a}Distance to the 1\,m success region (often called DTS), the
definition our evaluator uses. \textsuperscript{b}Distance to the goal object as defined in the paper's text
(SemUtil and L3MVN run SemExp-derived code, so their printed value may be closer to \textsuperscript{a}); not
directly comparable with \textsuperscript{a}. \textsuperscript{c}Third-party implementation: FBE by PONI, LGX and
LFG by OSG-Nav, and VLFM's value map with an FMM planner by GAMap. \textsuperscript{d}GAMap's text states a 0.1\,m
success rule and attributes its Gibson episodes to Al-Halah et al., which conflicts with the SemExp/L3MVN protocol it
says it follows; we print its Table 1 as given. \textsuperscript{e}HRO's detector is a Mask R-CNN of unstated training. SemExp: the released model's figure is stated in the authors' code
release and reprinted by OSG-Nav and VLFM; the SemExp paper reports 54.4/19.9 on an earlier episode set, and PONI's
re-evaluation of the released model gives 71.7/39.6, and L3MVN prints 65.2/33.6. ZSON~\citep{majumdar2022zson} reports 31.3/12.0 on a different Gibson episode set with RGB-only input
and is omitted, as is ZER (11.3 \SR), also on another episode set. Also omitted: OSG-Nav's own FBE run (64.1/28.3), which
nearly duplicates PONI's; RegQLearn (Gireesh et al., CASE 2022), whose printed figure is the best of five runs;
Think-Remember-Navigate and REST, which do not state their Gibson episode set; and L3MVN's feed-forward variant,
which trains object--room heads on HM3D. \ours adds LOOK\_UP and LOOK\_DOWN (30$^\circ$ steps) to the four actions of
the SemExp protocol that the baselines use (\cref{tab:privileges}).}
\end{table}

\paragraph{HM3D.}
\Cref{tab:hm3d} places HM3D v1 (2{,}000 episodes, 20 scenes) beside HM3D v2 (1{,}000 episodes, 36 scenes). The two
are not interchangeable: under one and the same protocol ApexNav reaches 59.6 \SR on v1 and 76.2 on v2, so
comparisons are meaningful only within a column. The official rule requires STOP within 0.1\,m of a goal view point,
whereas ApexNav~\citep{zhang2025apexnav} evaluates itself at 0.2\,m on both versions, as it does its re-evaluations of
L3MVN, VLFM and SG-Nav on v2; a looser radius can only raise \SR and \SPL. We therefore score \ours at 0.1\,m and
re-score the same episodes at 0.2\,m for comparison with those rows; both scores use the navigation mesh that habitat-sim recomputes for
the ObjectNav agent (\cref{sec:experiments}). HM3D v1 contains episodes whose goal lies on another storey, which an
agent confined to its start storey cannot solve; \TODO{state whether stair traversal is enabled on HM3D and report
the share of failures whose goal lies on another storey}. \TODO{describe how \ours compares with the training-free
rows at 0.1\,m and with ApexNav and its re-evaluations at 0.2\,m on each version, and its gap to the trained
references}.

\begin{table}[tp]
\centering
\small
\caption{HM3D ObjectNav validation. HM3D v1: HM3D-Semantics v0.1, 2022 challenge episodes (2{,}000). HM3D v2:
HM3D-Semantics v0.2 episodes (1{,}000) with the same agent and actions as v1. \SR and \SPL in \%; -- means not
reported. The two versions are not comparable with each other.}
\label{tab:hm3d}
\begin{tabular}{@{}llcccc@{}}
\toprule
 & & \multicolumn{2}{c}{HM3D v1} & \multicolumn{2}{c}{HM3D v2} \\
\cmidrule(lr){3-4}\cmidrule(l){5-6}
Method & Regime & \SR & \SPL & \SR & \SPL \\
\midrule
\multicolumn{6}{@{}l}{\emph{Trained}} \\
SemExp\textsuperscript{b}~\citep{chaplot2020object} & Trained & 37.9 & 18.8 & -- & -- \\
SGM~\citep{zhang2024imagine} & Self-sup. & 60.2 & 30.8 & -- & -- \\
OVRL-v2\textsuperscript{e}~\citep{yadav2023ovrl} & Trained & 64.7 & 28.1 & -- & -- \\
PIRLNav~\citep{ramrakhya2023pirlnav} & Trained & 70.4 & 34.1 & -- & -- \\
\midrule
\multicolumn{6}{@{}l}{\emph{Zero-shot}} \\
ZSON\textsuperscript{e}~\citep{majumdar2022zson} & Other task & 25.5 & 12.6 & -- & -- \\
GoW, ESC's implementation~\citep{zhou2023esc} & TF & 33.1 & 18.8 & -- & -- \\
PixNav\textsuperscript{c,e}~\citep{cai2024bridging} & Other task & 37.9 & 20.5 & -- & -- \\
ESC~\citep{zhou2023esc} & TF & 39.2 & 22.3 & -- & -- \\
VoroNav~\citep{wu2024voronav} & TF & 42.0 & 26.0 & -- & -- \\
PSL\textsuperscript{e}~\citep{sun2024prioritized} & Other task & 42.4 & 19.2 & -- & -- \\
L3MVN~\citep{yu2023l3mvn} & TF$^{\ast}$ & 50.4\textsuperscript{c} & 23.1\textsuperscript{c} & 36.3\textsuperscript{a,b} & 15.7\textsuperscript{a,b} \\
VLFM value map + FMM (GAMap)\textsuperscript{b}~\citep{yokoyama2024vlfm} & TF & 50.9 & 23.6 & -- & -- \\
TopV-Nav\textsuperscript{d}~\citep{zhong2024topvnav} & TF & 52.0 & 28.6 & -- & -- \\
VLFM~\citep{yokoyama2024vlfm} & Other task & 52.5 & 30.4 & 63.6\textsuperscript{a,b} & 32.5\textsuperscript{a,b} \\
GAMap~\citep{yuan2024gamap} & TF & 53.1 & 26.0 & -- & -- \\
SG-Nav~\citep{yin2024sg} & TF & 54.0 & 24.9 & 49.6\textsuperscript{a,b} & 25.5\textsuperscript{a,b} \\
UniGoal\textsuperscript{d,e}~\citep{yin2025unigoal} & TF & 54.5 & 25.1 & -- & -- \\
OpenFMNav~\citep{kuang2024openfmnav} & TF & 54.9 & 24.4 & -- & -- \\
TriHelper~\citep{zhang2024trihelper} & TF$^{\ast}$ & 56.5 & 25.3 & -- & -- \\
VLFM* shortest-path planner (ApexNav)\textsuperscript{b}~\citep{yokoyama2024vlfm} & TF & -- & -- & 56.9\textsuperscript{a} & 27.5\textsuperscript{a} \\
InstructNav\textsuperscript{c,d}~\citep{long2024instructnav} & TF & -- & -- & 58.0 & 20.9 \\
WMNav\textsuperscript{d,e}~\citep{nie2025wmnav} & TF & 58.1 & 31.2 & 72.2\textsuperscript{c} & 33.3\textsuperscript{c} \\
ApexNav\textsuperscript{a}~\citep{zhang2025apexnav} & TF & 59.6 & 33.0 & 76.2 & 38.0 \\
BeliefMapNav~\citep{zhou2025beliefmapnav} & TF & 61.4 & 30.6 & -- & -- \\
PIGEON-ZeroShot\textsuperscript{d}~\citep{peng2025pigeon} & TF & 67.1 & 35.8 & 79.6\textsuperscript{c} & 38.1\textsuperscript{c} \\
CogNav\textsuperscript{c}~\citep{cao2025cognav} & TF & 72.5 & 26.2 & -- & -- \\
\midrule
\textbf{\ours (ours)}, 0.1\,m & TF & \tbd & \tbd & \tbd & \tbd \\
\textbf{\ours (ours)}, re-scored at 0.2\,m\textsuperscript{a} & TF & \tbd & \tbd & \tbd & \tbd \\
\bottomrule
\end{tabular}
\par\smallskip
\parbox{\linewidth}{\footnotesize Regime as in \cref{tab:gibson}; here TF$^{\ast}$ marks a segmenter fine-tuned
on HM3D. \textsuperscript{a}0.2\,m success radius (ApexNav's setting) instead of the official 0.1\,m.
\textsuperscript{b}Evaluated by a third party: SemExp by L3MVN; VLFM's value map with an FMM planner in place of its
PointNav policy, on v1, by GAMap; and the v2 rows of L3MVN, VLFM, VLFM with a shortest-path planner (VLFM*) and SG-Nav
by ApexNav. \textsuperscript{c}Subset or unstated episode count: CogNav's own failure count implies about 1{,}008
episodes; WMNav's v2 number comes from a module ablation; PixNav and L3MVN do not state their HM3D v1 episode counts;
PIGEON describes v2 as 1{,}000 episodes over 10 scenes, where other papers count 36; InstructNav's paper states no
HM3D version or episode count, and its released code loads the v2 episodes and samples 500 by default, so we place it
under v2. \textsuperscript{d}Other or unclear success radius: UniGoal 1.6\,m, InstructNav 0.25\,m in its released
code, TopV-Nav `$<1$\,m', WMNav inconsistent (1.0\,m and 0.1\,m), PIGEON not stated. \textsuperscript{e}Non-standard
input or actions: RGB only (ZSON, PixNav, PSL, OVRL-v2), VLM-proposed polar moves (WMNav), or no look-up/down actions
(UniGoal). PIGEON rows are from its second arXiv version, which differs from the first. SG-Nav's v1 row is its GPT-4 variant (LLaMA-7B: 53.9/24.8); ApexNav does not
state which variant it re-evaluated on v2. ESC, TopV-Nav and CogNav rows are from their latest arXiv versions, which
differ from earlier ones. Omitted: results on self-sampled subsets, OSG-Nav~\citep{loo2025open} (400 episodes) and
ImagineNav (about 200); test-standard or unstated-split numbers, ProcTHOR and ProcTHOR-ZS~\citep{deitke2022procthor};
method variants, L3MVN's feed-forward model (trained on HM3D) and PIRLNav's behaviour-cloning-only model; and figures
we could not check against their own papers (those PIGEON quotes for STRIVE, ASCENT, Uni-NaVid and ImagineNav++).}
\end{table}

\paragraph{MP3D.}
The left half of \cref{tab:mp3d_robothor} covers MP3D (2{,}195 episodes, 11 scenes, 21 categories). Its vocabulary
adds categories such as cushions, towels, counters and fireplaces to Gibson's six, which tests whether the room
labels and the oracle's commonsense extend beyond the usual household furniture; \TODO{describe how the detector
prompts and the room-context vocabulary are extended to the 21 MP3D categories}. SemExp's own MP3D figure was
measured on its 2{,}000-episode, six-category subset, so we print a third-party evaluation on the full split.
ApexNav and SG-Nav~\citep{yin2024sg} score MP3D at 0.2\,m, and several training-free rows exist only as SG-Nav's
re-runs. As on HM3D, part of the episodes require stairs (\cref{sec:experiments}), so the storey caveat applies here
as well. \TODO{describe the standing
of \ours on MP3D relative to the strongest training-free row at 0.1\,m, to the 0.2\,m rows, and to the trained
references}.

\paragraph{RoboTHOR.}
The right half of \cref{tab:mp3d_robothor} covers RoboTHOR (1{,}800 episodes, 15 apartments, 12 categories). Most of
its targets are small objects such as mugs, apples and spray bottles, which shifts the difficulty from deciding where
to search towards detecting and closing on a small object, and success requires the target to be visible at STOP;
\TODO{describe the detector prompts for the 12 RoboTHOR categories}. ESC's result~\citep{zhou2023esc} did not
reproduce when re-run by \citet{unlu2024reproducibility}, so we list both; L3MVN, VLFM and OpenFMNav exist
only as SG-Nav's re-runs; and EmbCLIP~\citep{khandelwal2022simple}, trained on RoboTHOR and evaluated on the same
validation episodes, serves as the trained reference, while test-split results are left out. Our RoboTHOR runs use
ground-truth pose, which the official runner does not provide and which ESC, for one, runs without.
\TODO{describe the standing of \ours on RoboTHOR relative to the strongest training-free row and to EmbCLIP}.

\begin{table}[tp]
\centering
\small
\caption{MP3D ObjectNav validation (2021 challenge, 2{,}195 episodes) and RoboTHOR ObjectNav validation (2021
challenge, 1{,}800 episodes). \SR and \SPL in \%; -- means not reported. Validation splits only.}
\label{tab:mp3d_robothor}
\begin{tabular}{@{}llcccc@{}}
\toprule
 & & \multicolumn{2}{c}{MP3D} & \multicolumn{2}{c}{RoboTHOR} \\
\cmidrule(lr){3-4}\cmidrule(l){5-6}
Method & Regime & \SR & \SPL & \SR & \SPL \\
\midrule
\multicolumn{6}{@{}l}{\emph{Trained}} \\
SemExp\textsuperscript{b}~\citep{chaplot2020object} & Trained & 28.3 & 10.9 & -- & -- \\
PONI~\citep{ramakrishnan2022poni} & Trained & 31.8 & 12.1 & -- & -- \\
Habitat-Web~\citep{ramrakhya2022habitat} & Trained & 35.4 & 10.2 & -- & -- \\
SGM~\citep{zhang2024imagine} & Self-sup. & 37.7 & 14.7 & -- & -- \\
EmbCLIP~\citep{khandelwal2022simple} & Trained & -- & -- & 52.2 & 26.0 \\
\midrule
\multicolumn{6}{@{}l}{\emph{Zero-shot}} \\
CoW\textsuperscript{e}~\citep{gadre2023cows} & TF & 7.4 & 3.7 & 26.7 & 16.9 \\
ZSON\textsuperscript{e}~\citep{majumdar2022zson} & Other task & 15.3 & 4.8 & -- & -- \\
ESC~\citep{zhou2023esc} & TF & 28.7 & 14.2 & 38.1\textsuperscript{d} & 22.2\textsuperscript{d} \\
VLFM* shortest-path planner\textsuperscript{b}~\citep{yokoyama2024vlfm} & TF & 32.5\textsuperscript{a} & 15.9\textsuperscript{a} & -- & -- \\
L3MVN\textsuperscript{b}~\citep{yu2023l3mvn} & TF & 34.9 & 14.5 & 41.2 & 22.5 \\
TopV-Nav\textsuperscript{f}~\citep{zhong2024topvnav} & TF & 35.2 & 16.4 & -- & -- \\
VLFM~\citep{yokoyama2024vlfm} & Other task & 36.4 & 17.5 & 42.3\textsuperscript{b} & 23.0\textsuperscript{b} \\
OpenFMNav\textsuperscript{b}~\citep{kuang2024openfmnav} & TF & 37.2 & 15.7 & 44.1 & 23.3 \\
ApexNav~\citep{zhang2025apexnav} & TF & 39.2\textsuperscript{a} & 17.8\textsuperscript{a} & -- & -- \\
SG-Nav, LLaMA-7B~\citep{yin2024sg} & TF & 40.1\textsuperscript{a} & 16.0\textsuperscript{a} & 47.3\textsuperscript{c} & 23.7\textsuperscript{c} \\
SG-Nav, GPT-4~\citep{yin2024sg} & TF & 40.2\textsuperscript{a} & 16.0\textsuperscript{a} & 47.5\textsuperscript{c} & 24.0\textsuperscript{c} \\
UniGoal\textsuperscript{f}~\citep{yin2025unigoal} & TF & 41.0 & 16.4 & 48.0\textsuperscript{c} & 24.2\textsuperscript{c} \\
WMNav\textsuperscript{e,f}~\citep{nie2025wmnav} & TF & 45.4 & 17.2 & -- & -- \\
PIGEON-ZeroShot\textsuperscript{f}~\citep{peng2025pigeon} & TF & 52.6 & 18.5 & -- & -- \\
CogNav~\citep{cao2025cognav} & TF & 46.6\textsuperscript{c} & 16.1\textsuperscript{c} & 54.6 & 24.3 \\
VLTNet~\citep{vltnet2024} & TF & -- & -- & 33.2 & 17.1 \\
ESC, re-run by Unlu et al.~\citep{unlu2024reproducibility} & TF & -- & -- & 34.5 & 18.2 \\
LGX~\citep{dorbala2024can} & TF & -- & -- & 35.0 & 21.9 \\
Unlu et al.~\citep{unlu2024reproducibility} & TF & -- & -- & 35.2 & 18.3 \\
GoW, ESC's implementation~\citep{zhou2023esc} & TF & -- & -- & 36.1 & 21.6 \\
RF-Nav (DEFA + CDM)~\citep{yuan2024defa} & TF & -- & -- & 36.8 & 22.3 \\
\midrule
\textbf{\ours (ours)} & TF & \tbd & \tbd & \tbd & \tbd \\
\bottomrule
\end{tabular}
\par\smallskip
\parbox{\linewidth}{\footnotesize Regime as in \cref{tab:gibson}. \textsuperscript{a}0.2\,m success radius on MP3D
(ApexNav's paper and configuration, SG-Nav's released configuration) instead of 0.1\,m. \textsuperscript{b}Evaluated
by a third party: SemExp on the full MP3D split by 3D-aware ObjectNav~\citep{zhang20233d} (as printed by SGM),
VLFM with a shortest-path planner in place of its PointNav policy (VLFM*) by ApexNav, and L3MVN and OpenFMNav on both
benchmarks and VLFM on RoboTHOR by SG-Nav; these methods report no such result themselves, and SG-Nav does not state
which segmenter its L3MVN re-run used. \textsuperscript{c}Subset or unstated episode count: CogNav describes its MP3D
evaluation as 1.8K episodes and 20 categories; SG-Nav, and UniGoal, which follows its setting, do not state that all
1{,}800 RoboTHOR episodes were run. \textsuperscript{d}Not reproduced; see the ESC re-run row.
\textsuperscript{e}Non-standard input or actions: ZSON uses RGB only and selects its checkpoint on ObjectNav validation
\SR; WMNav uses VLM-proposed polar moves within 40 macro steps; CoW acts without look-up/down actions.
\textsuperscript{f}Other or unclear success radius: UniGoal states 1.6\,m for ObjectNav and otherwise follows SG-Nav's
setting (0.2\,m on MP3D); TopV-Nav states `$<1$\,m'; WMNav's text gives both 1.0\,m and 0.1\,m; PIGEON states none (its second arXiv
version; the first printed 41.6/14.4). CoW is its OWL-ViT B/32 variant (its best, B/16, reaches 27.5/17.2 on RoboTHOR), and
LGX's own implementation of GoW reaches 33.2/20.3 on RoboTHOR. TopV-Nav is its second arXiv version. EmbCLIP's figure
is its entry on the official RoboTHOR validation leaderboard. SemExp's own MP3D result (36.0/14.4) was measured on a
2{,}000-episode, six-category subset and is not printed. RoboTHOR test-split results are omitted. \ours uses
ground-truth pose on RoboTHOR (\cref{tab:privileges}).}
\end{table}

\subsection{Ablation Studies}\label{sec:ablations}

\Cref{tab:ablation} changes one component at a time and keeps every other setting at its default
(\cref{app:hyperparams}). Every variant is to be run on the same episodes as the full method, so that differences
can be tested with the paired tests of \cref{sec:experiments}. \TODO{state the episodes used for the ablations (full
splits or a disclosed subset), or run the ablations on the Gibson training scenes or a disclosed subset, as the
development protocol requires; mark the differences that the paired tests support, and confirm that no ablation result
was used to select the defaults}.

\emph{Search ordering.} Rows (b) and (c) keep the candidate nodes, probabilities and cost matrix but replace
RPT$^{*}$~\citep{lyu2026rpt} with greedy rules that each ignore one factor of \cref{eq:objective}:
greedy-by-probability visits the most probable remaining node whatever it costs to reach, and nearest-first visits
the cheapest remaining node whatever its probability. Comparing them with row (a) measures what an exact expected-time
order buys over heuristics with the same inputs; \TODO{report the effect on \SR and \SPL}. Because the solver falls
back to drawing the next node in proportion to $p_i$ when an instance violates its precondition, row (a) is read
together with the fallback rates of \cref{tab:runtime}.

\emph{Node probabilities.} Uniform probabilities (d) would remove commonsense entirely, so that the order follows
geometry alone; the current implementation refuses a uniform prior, so this variant still has to be built. Row (e) adds hand-coded rules to the oracle: a firmly labelled room whose type cannot hold the target is
removed from the candidates, and unexplored nodes and rooms holding an object that usually accompanies the target
receive a minimum probability; it tests whether such rules complement or contradict the oracle's own commonsense.
Removing room labels from the oracle prompt (f), a variant that also still has to be built, would measure how much
of the oracle's signal comes from room classification rather than from the raw object lists. \TODO{report the effect of each}.

\emph{Candidate nodes.} Rows (g) and (h) remove opening nodes and sighted-target landmark nodes, and row (i) adds the
second pass over finished rooms, which is off by default. \TODO{report the effect, in particular in large multi-room houses where openings dominate the
candidate set}.

\emph{Local exploration.} Rows (j)--(m) replace the frontier-walk room visit with a sweep visit that spends ten actions on frontier
goals inside the room, disable path glances or the sight ledger, or replace the local visit with a bounded coverage
explorer adapted from FALCON~\citep{zhang2025falcon}. \TODO{report the effect, especially on \SPL, since these
components trade actions for coverage}.

\emph{Target closing, perception and protocol.} Row (n), which has no implementation yet, would remove the
verification gates of the closing routine (\cref{sec:execution}): the approach would start without the two-frame lock,
and STOP would be issued on arrival at the standoff without the inspection's freshness, centring, confidence and range
tests; the implemented switches instead disable individual closing components. It would measure how many wrong STOPs the gates
prevent and how many actions they cost. Row (o)
replaces YOLO-World~\citep{cheng2024yolo} with Grounding DINO~\citep{liu2024grounding} proposals verified by
BLIP-2~\citep{li2023blip}, and row (p) adds Grounding DINO proposals to those of YOLO-World and verifies them with
BLIP-2. Row (q) reads the target's depth on a MobileSAM~\citep{zhang2023faster} mask, derived from
SAM~\citep{kirillov2023segment}, instead of on the box. Row (r) removes the depth-hole filling, a disclosed deviation
from the baselines (\cref{tab:privileges}). \TODO{report the effect of each, including the change in the rate of wrong
STOPs}.

\begin{table}[t]
\centering
\small
\caption{Ablations. Each row changes one component of \ours and keeps all others at their defaults. \SR and \SPL
in \%, \DTG in metres. \protect\TODO{state the split of each column group (validation, training scenes or a disclosed
subset)}.}
\label{tab:ablation}
\begin{tabular}{@{}cp{7.4cm}ccccc@{}}
\toprule
 & & \multicolumn{3}{c}{Gibson} & \multicolumn{2}{c}{HM3D v1} \\
\cmidrule(lr){3-5}\cmidrule(l){6-7}
 & Variant & \SR & \SPL & \DTG & \SR & \SPL \\
\midrule
(a) & \ours (full) & \tbd & \tbd & \tbd & \tbd & \tbd \\
\midrule
\multicolumn{7}{@{}l}{\emph{Search ordering}} \\
(b) & Greedy by probability instead of RPT$^{*}$ & \tbd & \tbd & \tbd & \tbd & \tbd \\
(c) & Nearest-first instead of RPT$^{*}$ & \tbd & \tbd & \tbd & \tbd & \tbd \\
\multicolumn{7}{@{}l}{\emph{Node probabilities}} \\
(d) & Uniform probabilities instead of the LLM oracle \TODO{not implemented yet; implement and freeze before running} & \tbd & \tbd & \tbd & \tbd & \tbd \\
(e) & Hand-coded room-type exclusions and probability floors added to the oracle & \tbd & \tbd & \tbd & \tbd & \tbd \\
(f) & No room labels in the oracle prompt \TODO{not implemented yet; implement and freeze before running} & \tbd & \tbd & \tbd & \tbd & \tbd \\
\multicolumn{7}{@{}l}{\emph{Candidate nodes}} \\
(g) & No opening nodes & \tbd & \tbd & \tbd & \tbd & \tbd \\
(h) & No target-landmark nodes & \tbd & \tbd & \tbd & \tbd & \tbd \\
(i) & Second pass over finished rooms added & \tbd & \tbd & \tbd & \tbd & \tbd \\
\multicolumn{7}{@{}l}{\emph{Local exploration}} \\
(j) & Sweep visit (10 local actions) instead of frontier walk & \tbd & \tbd & \tbd & \tbd & \tbd \\
(k) & No path glances & \tbd & \tbd & \tbd & \tbd & \tbd \\
(l) & No sight ledger & \tbd & \tbd & \tbd & \tbd & \tbd \\
(m) & Bounded FALCON coverage explorer instead of the local visit & \tbd & \tbd & \tbd & \tbd & \tbd \\
\multicolumn{7}{@{}l}{\emph{Target closing, perception and protocol}} \\
(n) & Closing without verification gates \TODO{not implemented yet; implement and freeze before running} & \tbd & \tbd & \tbd & \tbd & \tbd \\
(o) & Grounding DINO + BLIP-2 instead of YOLO-World & \tbd & \tbd & \tbd & \tbd & \tbd \\
(p) & YOLO-World + Grounding DINO proposals, BLIP-2 verification & \tbd & \tbd & \tbd & \tbd & \tbd \\
(q) & MobileSAM mask depth instead of box depth & \tbd & \tbd & \tbd & \tbd & \tbd \\
(r) & No depth-hole filling & \tbd & \tbd & \tbd & \tbd & \tbd \\
\bottomrule
\end{tabular}
\end{table}

\subsection{Analysis}\label{sec:results-analysis}

\paragraph{Language-model size.}
\Cref{tab:llm} varies the two language models independently. The default pairs a 3B room labeller with a 14B node
oracle: labelling a room is a short classification from an object list, whereas the oracle reasons over the whole
house summary. The oracle is called only at decision points, and an unchanged prompt reuses the previous reply, so
its latency adds wall-clock time but no actions. \citet{loo2025open} report that the HM3D \SR of
their scene-graph agent rises from 34.9 with LLaMA~2~7B to 69.3 with GPT-3.5 on a 100-episode subset, which
indicates that such agents can be sensitive to the reasoning model. \TODO{report whether \ours shows a similar
sensitivity between the 3B and 14B oracles, and what a 14B labeller adds}.

\begin{table}[t]
\centering
\small
\caption{Choice of language models on Gibson. Both models are Qwen2.5-Instruct~\citep{yang2024qwen25} served
locally (the 14B model in 4-bit quantisation). Calls are averaged per episode; latency is the median per oracle call
on \protect\TODO{the machine of each row; the 3B/3B pair is the configuration that fits the 8\,GB laptop GPU of
\cref{sec:experiments}}.}
\label{tab:llm}
\begin{tabular}{@{}llccccc@{}}
\toprule
Room labeller & Node oracle & \SR & \SPL & Label calls & Oracle calls & Oracle latency (s) \\
\midrule
3B & 14B (default) & \tbd & \tbd & \tbd & \tbd & \tbd \\
3B & 3B & \tbd & \tbd & \tbd & \tbd & \tbd \\
14B & 3B & \tbd & \tbd & \tbd & \tbd & \tbd \\
14B & 14B & \tbd & \tbd & \tbd & \tbd & \tbd \\
\bottomrule
\end{tabular}
\end{table}

\paragraph{Runtime.}
\Cref{tab:runtime} breaks the cost of the policy into components. Detection, mapping and scene-graph maintenance run
on every frame and set the minimum time per action, whereas the two language models, RPT$^{*}$ and route planning run
at decision points or on re-planning and are amortised over the actions of a visit. RPT$^{*}$ is exact, and its cost
is bounded by the cap of 12 nodes per instance and a 2\,s search budget (\cref{sec:planning}); the last block of the
table gives, for each benchmark, the share of decision points whose order was certified optimal, returned on budget
expiry without a certificate, or drawn by the random fallback. \TODO{report the machine, the median and 95th-percentile latency of each component, the calls per
episode, and the share of episode wall time spent waiting for the oracle}.

\begin{table}[t]
\centering
\small
\caption{Per-component latency of \ours on Gibson, measured on \protect\TODO{machine}. Median and 95th percentile per call
in milliseconds; calls per episode averaged over episodes. The last block gives, per benchmark, the outcome of the
visit-ordering step as a share of decision points.}
\label{tab:runtime}
\begin{tabular}{@{}p{2.9cm}p{3.3cm}p{3.9cm}ccc@{}}
\toprule
Component & Model or algorithm & Invoked & Median & p95 & Calls \\
\midrule
Object detection & YOLO-World X-v2 & every frame & \tbd & \tbd & \tbd \\
Mapping and segmentation & occupancy, landmark fusion, doorways, rooms & every frame & \tbd & \tbd & \tbd \\
Room labelling & Qwen2.5-3B-Instruct & decision points, 50-action refresh, new object kind & \tbd & \tbd & \tbd \\
Node oracle & Qwen2.5-14B-Instruct & decision points, unless the prompt is unchanged & \tbd & \tbd & \tbd \\
Cost matrix & Dijkstra on the occupancy grid & decision points & \tbd & \tbd & \tbd \\
Visit ordering & RPT$^{*}$, at most 12 nodes & decision points & \tbd & \tbd & \tbd \\
Route planning & clearance-weighted A$^{*}$ & route planning and re-planning & \tbd & \tbd & \tbd \\
\midrule
Whole policy step & excluding rendering & every action & \tbd & \tbd & \tbd \\
\midrule
\multicolumn{3}{@{}l}{\emph{Visit ordering outcome (\% of decision points)}} & Certified & Budget & Random \\
Gibson & \multicolumn{2}{l}{RPT$^{*}$ certified / budget expired / random fallback} & \tbd & \tbd & \tbd \\
HM3D v1 & \multicolumn{2}{l}{as above} & \tbd & \tbd & \tbd \\
HM3D v2 & \multicolumn{2}{l}{as above} & \tbd & \tbd & \tbd \\
MP3D & \multicolumn{2}{l}{as above} & \tbd & \tbd & \tbd \\
RoboTHOR & \multicolumn{2}{l}{as above} & \tbd & \tbd & \tbd \\
\bottomrule
\end{tabular}
\end{table}

\paragraph{Failure modes.}
The failed episodes we inspected during development fall into the categories of \cref{tab:failures}; the main
mechanisms are the following (\cref{fig:qualitative}). (i)~\emph{False-positive STOPs.} A detector error that persists over consecutive frames at terminal range passes the
closing gates, for example a kitchen counter taken for a bed, bar seating for a sofa, or a ride-on toy for a chair;
multi-frame confirmation bounds such errors but cannot remove those that are consistent across views.
(ii)~\emph{Timeouts in large houses.} On floors with many rooms and openings, peeks and room visits exhaust the budget
before the target room is reached. (iii)~\emph{Navigation stuck.} Occasionally the agent wedges against geometry
the map does not yet show. (iv)~\emph{Closing failures.} A low target such as a toilet can be reached and then released when the lower image
edge cuts its box, as we observed under Gibson's four-action protocol, which offers no camera tilt; the tilt actions
our Gibson runs add (\cref{tab:privileges}) let the pitch servo of \cref{sec:execution} address this case,
\TODO{state whether they remove it}. Repeated release-and-retake cycles on a wrong support surface also
consume budget. (v)~\emph{Annotation gaps.} The agent sometimes stops at a real
instance of the goal category that the episode's goal set does not contain, which the evaluator scores as a
failure. \TODO{quantify each mechanism on each benchmark}. Per-category results and the failure breakdown are given
in \cref{tab:percategory} and \cref{tab:failures} (\cref{app:additional}).

\begin{figure}[t]
\centering
\fbox{\parbox[c][4.5cm][c]{0.95\linewidth}{\centering\TODO{Qualitative episodes: top-down maps with the online
scene graph (rooms coloured by oracle probability, doorways, openings, target landmarks), the RPT$^{*}$ visit order
at successive decision points and the executed trajectory; one success per benchmark and one example of each
failure mode (false-positive STOP, timeout, navigation stuck, closing failure, annotation gap)}}}
\caption{Qualitative examples of \ours. Each panel shows the scene graph and the planned visit order at selected
decision points, and the trajectory the agent executed.}
\label{fig:qualitative}
\end{figure}

% ============================================================================
% 10_conclusion
% ============================================================================
% ============================================================================
%  10_conclusion.tex -- Limitations and Conclusion
% ============================================================================

\section{Limitations}\label{sec:limitations}

\ours is training-free and evaluated only in simulation; the following conditions bound what our results can claim.

\begin{itemize}[leftmargin=*,itemsep=2pt]
\item \textbf{Privileged inputs.} The agent receives ground-truth pose and noise-free simulator depth whose holes are
filled (the filling is a deviation from the baselines). Staircase geometry comes from the navigation mesh through a
simulated perfect detector that reveals a staircase only once it is visible; on single-storey runs it is used solely to
mask seen stairs as obstacles. The closing routine projects its standoff goal onto the navigable mesh. Ground-truth pose
matches the localisation available to the compared methods on the Habitat benchmarks, but on RoboTHOR the official
protocol provides none (\cref{tab:privileges}). No goal coordinates, distance fields or evaluator signals reach the
policy, but every privileged input above would have to be estimated onboard on a physical robot.

\item \textbf{Uncalibrated probabilities.} The $p_i$ are zero-temperature estimates of a 4-bit quantised 14B-parameter
language model~\citep{yang2024qwen25}, not calibrated frequencies, and the values given to sighted target landmarks
(0.85), to nodes the reply omits (0.01) and to exits it has not valued (0.10) are set by hand, as are the factors,
applied before the probabilities are filtered and ordered, that discount rooms entered and left (0.5 or 0.2) and gaps
that would reveal little unknown floor (to no less than 0.25). We treat the $p_i$ as independent terminal probabilities
without normalisation, although the events that two nodes yield the target are neither independent nor exclusive: a
house may hold several instances, and adjacent rooms share evidence. Nodes scored below 0.05 are not offered at that
decision point, so an underestimate can postpone the visit to the right room. The prompt withholds distances but
marks the room the agent stands in; the model is instructed, not guaranteed, to ignore this.

\item \textbf{Scope of optimality.} RPT$^*$~\citep{lyu2026rpt} certifies an optimal order only for the instance it is
given; re-solving in a receding horizon is a heuristic for the partially observed search, not an optimal policy. To keep
the exact search, whose cost grows exponentially with the number of nodes, within its 2\,s budget, the instance is
capped at 12 nodes (the nearest plus the 11 most probable), so the order is optimal only over this subset. Travel costs
ignore turning, and service times are fixed per node type, except that a room's grows with its live frontier count
(\cref{eq:arccost}). When the cost matrix violates the
solver's metric precondition (mixing distances measured at two clearance levels can cause this) or the solver fails for
any other reason, the next node is drawn at random in proportion to $p_i$; when the budget expires, the best complete
order found, or else a nearest-neighbour order, is used without a guarantee (\cref{tab:runtime}).

\item \textbf{STOP depends on the detector.} Every STOP rests on YOLO-World~\citep{cheng2024yolo} detections: normally
a fresh, centred box with confidence at least 0.50 within 1.05\,m (a weaker box after at most two further frames),
otherwise recent remembered evidence or a landmark
confirmed from three views within 0.75\,m. A missed or truncated detection leaves the target unconfirmed. A persistent false positive at terminal range, such as a kitchen
counter read as a bed, yields a wrong STOP, and repeated release and re-acquisition of a wrong candidate consumes
budget. The optional Grounding DINO~\citep{liu2024grounding} and BLIP-2~\citep{li2023blip} verification is off by
default.

\item \textbf{Room segmentation heuristics.} The clearance watershed~\citep{vincent1991watersheds} cut at doorways is
one of many hand-crafted options~\citep{bormann2016room}. Open-plan layouts merge distinct functional areas into one
room; an L-shaped room may be declared finished once half of its cells are seen from a scan and no live frontier
remains, leaving its hidden leg
unsearched; a misplaced doorway can sever a room; and a signature object glimpsed through a doorway can mislabel a
room and mislead the oracle.

\item \textbf{Language-model latency.} Although the oracle runs only at decision points and reuses its reply to an
unchanged prompt, the 14B model, about 9\,GB of quantised weights, dominates wall-clock time and memory (\cref{tab:runtime}; \TODO{per-call oracle latency on the workstation and the
laptop}); \cref{tab:llm} examines a smaller oracle.

\item \textbf{Multi-storey search.} Every Gibson episode lies on one floor and stair traversal is disabled there.
\TODO{if disabled on HM3D/MP3D: the agent cannot solve the 14.6\% (HM3D v1) and 9.6\% (MP3D) of episodes that
require stairs; if enabled: it relies on the privileged stair geometry and storey metadata}.

\item \textbf{Annotation incompleteness.} A STOP beside a real but unannotated instance counts as a failure, so success
rates penalise correct behaviour to a dataset-dependent degree. In multi-storey Gibson buildings the semantic maps used
for scoring annotate only a reference floor, so instances on other floors cannot be credited.
\TODO{quantify wrong STOPs at unannotated instances}

\item \textbf{Validation exposure and protocols.} Development used Gibson training buildings, but some validation
episodes were inspected, the depth-hole filling was designed on frames from one validation episode, two component
choices were compared on validation episodes, the target-closing routine was revised after inspecting failures
from a validation run, and later revisions of the search, the perception, the target closing, the map update and the Gibson action
space were diagnosed and selected on 15 validation episodes with a toilet goal; \TODO{final frozen-evaluation
statement}. We use habitat-sim~\citep{savva2019habitat} 0.2.4 instead of the 0.1.5 of the SemExp
reference~\citep{chaplot2020object}, on Gibson we add camera tilt (LOOK\_UP and LOOK\_DOWN) to the four actions the
baselines use, and Gibson does not require STOP. On Gibson we load the shipped navigation mesh,
and whether the reference recomputed it for the agent's size is unknown; this affects reachability and collisions,
not the score. HM3D and MP3D use the navigation mesh recomputed for our agent, as habitat-lab does. Baseline numbers are
quoted, some with a 0.2\,m rather than 0.1\,m success radius, so cross-paper comparisons are approximate.

\item \textbf{Static scenes, closed vocabularies.} The map and object memory assume a static scene. Goals come from
each benchmark's fixed category set (6 categories on Gibson and HM3D, 21 on MP3D, 12 on RoboTHOR), and the detector
uses a fixed, hand-written prompt vocabulary (30 prompts on Gibson); attribute- or instance-specified goals are not
evaluated.
\end{itemize}

\section{Conclusion}\label{sec:conclusion}

We presented \ours, a training-free ObjectNav agent that casts object search as a sequence of small probabilistic
routing problems. An online scene graph partitions observed free space into rooms that a small language model labels;
a larger model estimates, without being told any distances, the probability that searching each candidate room or
opening finds the target, while sighted but unverified target landmarks receive a fixed probability; and RPT$^*$ orders
the candidates by expected time-to-find, with travel costs measured on the occupancy grid and service times that are
fixed per node type or, for a room, grow with its frontier count. The
agent executes only the first node, replans at visit boundaries, suspends the search for target closing on a
confident sighting and resumes the interrupted visit if the candidate is released, and normally stops only on a fresh,
confident detection at close range, through a verify--approach--inspect routine with a few bounded exceptions backed
by recent evidence. The contribution is this
coupling, not any single component.

Across the five validation splits on Gibson, HM3D (v1 and v2), MP3D and RoboTHOR, \ours reaches \TODO{SR and SPL
per benchmark}, \TODO{compare with the strongest training-free baselines under matched protocols}; the ablations of \cref{sec:ablations} \TODO{summarise which
components matter: RPT$^*$ versus greedy ordering, oracle size, opening nodes, target closing}.

Future work follows from these limitations. The oracle's probabilities could be calibrated against logged outcomes or
fused with observation-driven beliefs~\citep{zhou2025beliefmapnav}, and service costs could be learned from past visit
durations; replacing the cost matrix by its shortest-path closure, which is metric by construction, would remove the
random fallback triggered by triangle-inequality violations. Substituting onboard estimates for the privileged pose,
depth and stair geometry would enable deployment on ground robots,
where modular map-based agents transfer well~\citep{gervet2023navigating}, and on aerial platforms, with onboard
perception replacing the privileged inputs.
Because the detector is open-vocabulary, goals could extend to free-form or instance descriptions~\citep{yin2025unigoal},
given verification robust to a larger confusion set. Finally, a learned stair detector and multi-floor buildings with
complete annotations would close an important gap between our evaluation and real homes.

\clearpage
\begin{thebibliography}{87}
\providecommand{\natexlab}[1]{#1}

\bibitem[{Anderson et~al.(2018)}]{anderson2018evaluation}
Peter Anderson, Angel Chang, Devendra~Singh Chaplot, Alexey Dosovitskiy, Saurabh Gupta, Vladlen Koltun, Jana Kosecka, Jitendra Malik, et~al.
\newblock On evaluation of embodied navigation agents.
\newblock \emph{arXiv preprint arXiv:1807.06757}, 2018.

\bibitem[{Armeni et~al.(2019)}]{armeni20193d}
Iro Armeni, Zhi-Yang He, JunYoung Gwak, Amir~R. Zamir, Martin Fischer, Jitendra Malik, and Silvio Savarese.
\newblock 3D scene graph: A structure for unified semantics, 3D space, and camera.
\newblock In \emph{IEEE/CVF International Conference on Computer Vision (ICCV)}, 2019.

\bibitem[{Aydemir et~al.(2013)}]{aydemir2013active}
Alper Aydemir, Andrzej Pronobis, Moritz G{\"o}belbecker, and Patric Jensfelt.
\newblock Active visual object search in unknown environments using uncertain semantics.
\newblock \emph{IEEE Transactions on Robotics (T-RO)}, 29(4):986--1002, 2013.

\bibitem[{Batra et~al.(2020)}]{batra2020objectnav}
Dhruv Batra, Aaron Gokaslan, Aniruddha Kembhavi, Oleksandr Maksymets, Roozbeh Mottaghi, Manolis Savva, Alexander Toshev, and Erik Wijmans.
\newblock {ObjectNav} revisited: On evaluation of embodied agents navigating to objects.
\newblock \emph{arXiv preprint arXiv:2006.13171}, 2020.

\bibitem[{Blum et~al.(1994)}]{blum1994minimum}
Avrim Blum, Prasad Chalasani, Don Coppersmith, Bill Pulleyblank, Prabhakar Raghavan, and Madhu Sudan.
\newblock The minimum latency problem.
\newblock In \emph{Proceedings of the 26th Annual ACM Symposium on Theory of Computing (STOC)}, 1994.

\bibitem[{Bormann et~al.(2016)}]{bormann2016room}
Richard Bormann, Florian Jordan, Wenzhe Li, Joshua Hampp, and Martin H\"{a}gele.
\newblock Room segmentation: Survey, implementation, and analysis.
\newblock In \emph{IEEE International Conference on Robotics and Automation (ICRA)}, pages 1019--1026, 2016.

\bibitem[{Cai et~al.(2024)}]{cai2024bridging}
Wenzhe Cai, Siyuan Huang, Guangran Cheng, Yuxing Long, Peng Gao, Changyin Sun, and Hao Dong.
\newblock Bridging zero-shot object navigation and foundation models through pixel-guided navigation skill.
\newblock In \emph{IEEE International Conference on Robotics and Automation (ICRA)}, 2024.

\bibitem[{Cao et~al.(2025)}]{cao2025cognav}
Yihan Cao, Jiazhao Zhang, Zhinan Yu, Shuzhen Liu, Zheng Qin, Qin Zou, Bo~Du, and Kai Xu.
\newblock {CogNav}: Cognitive process modeling for object goal navigation with {LLMs}.
\newblock In \emph{Proceedings of the IEEE/CVF International Conference on Computer Vision (ICCV)}, 2025.

\bibitem[{Chang et~al.(2017)}]{chang2017matterport3d}
Angel Chang, Angela Dai, Thomas Funkhouser, Maciej Halber, Matthias Nie{\ss}ner, Manolis Savva, Shuran Song, Andy Zeng, et~al.
\newblock {Matterport3D}: Learning from {RGB-D} data in indoor environments.
\newblock In \emph{International Conference on 3D Vision (3DV)}, 2017.

\bibitem[{Chaplot et~al.(2020)}]{chaplot2020object}
Devendra~Singh Chaplot, Dhiraj Gandhi, Abhinav Gupta, and Ruslan Salakhutdinov.
\newblock Object goal navigation using goal-oriented semantic exploration.
\newblock In \emph{Advances in Neural Information Processing Systems (NeurIPS)}, 2020.

\bibitem[{Chen et~al.(2023)}]{chen2023how}
Junting Chen, Guohao Li, Suryansh Kumar, Bernard Ghanem, and Fisher Yu.
\newblock How to not train your dragon: Training-free embodied object goal navigation with semantic frontiers.
\newblock In \emph{Robotics: Science and Systems (RSS)}, 2023.

\bibitem[{Cheng et~al.(2024)}]{cheng2024yolo}
Tianheng Cheng, Lin Song, Yixiao Ge, Wenyu Liu, Xinggang Wang, and Ying Shan.
\newblock {YOLO-World}: Real-time open-vocabulary object detection.
\newblock In \emph{IEEE/CVF Conference on Computer Vision and Pattern Recognition (CVPR)}, pages 16901--16911, 2024.

\bibitem[{Datta et~al.(2021)}]{datta2021integrating}
Samyak Datta, Oleksandr Maksymets, Judy Hoffman, Stefan Lee, Dhruv Batra, and Devi Parikh.
\newblock Integrating egocentric localization for more realistic point-goal navigation agents.
\newblock In \emph{Proceedings of the 2020 Conference on Robot Learning (CoRL)}, volume 155 of \emph{Proceedings of Machine Learning Research}, pages 313--328, 2021.

\bibitem[{Deitke et~al.(2020)}]{deitke2020robothor}
Matt Deitke, Winson Han, Alvaro Herrasti, Aniruddha Kembhavi, Eric Kolve, Roozbeh Mottaghi, Jordi Salvador, Dustin Schwenk, et~al.
\newblock {RoboTHOR}: An open simulation-to-real embodied {AI} platform.
\newblock In \emph{Proceedings of the IEEE/CVF Conference on Computer Vision and Pattern Recognition (CVPR)}, 2020.

\bibitem[{Deitke et~al.(2022)}]{deitke2022procthor}
Matt Deitke, Eli VanderBilt, Alvaro Herrasti, Luca Weihs, Jordi Salvador, Kiana Ehsani, Winson Han, Eric Kolve, et~al.
\newblock {ProcTHOR}: Large-scale embodied {AI} using procedural generation.
\newblock In \emph{Advances in Neural Information Processing Systems (NeurIPS)}, 2022.

\bibitem[{Dijkstra(1959)}]{dijkstra1959note}
Edsger~W. Dijkstra.
\newblock A note on two problems in connexion with graphs.
\newblock \emph{Numerische Mathematik}, 1(1):269--271, 1959.

\bibitem[{Dorbala et~al.(2024)}]{dorbala2024can}
Vishnu~Sashank Dorbala, James~F. Mullen~Jr., and Dinesh Manocha.
\newblock Can an embodied agent find your ``cat-shaped mug''? {LLM}-based zero-shot object navigation.
\newblock \emph{IEEE Robotics and Automation Letters (RA-L)}, 9(5):4083--4090, 2024.

\bibitem[{Gadre et~al.(2023)}]{gadre2023cows}
Samir~Yitzhak Gadre, Mitchell Wortsman, Gabriel Ilharco, Ludwig Schmidt, and Shuran Song.
\newblock CoWs on pasture: Baselines and benchmarks for language-driven zero-shot object navigation.
\newblock In \emph{Proceedings of the IEEE/CVF Conference on Computer Vision and Pattern Recognition (CVPR)}, 2023.

\bibitem[{Gervet et~al.(2023)}]{gervet2023navigating}
Theophile Gervet, Soumith Chintala, Dhruv Batra, Jitendra Malik, and Devendra~Singh Chaplot.
\newblock Navigating to objects in the real world.
\newblock \emph{Science Robotics}, 8(79):eadf6991, 2023.

\bibitem[{Gu et~al.(2024)}]{gu2024conceptgraphs}
Qiao Gu, Alihusein Kuwajerwala, Sacha Morin, Krishna~Murthy Jatavallabhula, Bipasha Sen, Aditya Agarwal, Corban Rivera, William Paul, et~al.
\newblock ConceptGraphs: Open-vocabulary 3D scene graphs for perception and planning.
\newblock In \emph{IEEE International Conference on Robotics and Automation (ICRA)}, 2024.

\bibitem[{Hart et~al.(1968)}]{hart1968formal}
Peter~E. Hart, Nils~J. Nilsson, and Bertram Raphael.
\newblock A formal basis for the heuristic determination of minimum cost paths.
\newblock \emph{IEEE Transactions on Systems Science and Cybernetics}, 4(2):100--107, 1968.

\bibitem[{Honerkamp et~al.(2024)}]{honerkamp2024language}
Daniel Honerkamp, Martin B{\"u}chner, Fabien Despinoy, Tim Welschehold, and Abhinav Valada.
\newblock Language-grounded dynamic scene graphs for interactive object search with mobile manipulation.
\newblock \emph{IEEE Robotics and Automation Letters (RA-L)}, 9(10):8298--8305, 2024.

\bibitem[{Huang et~al.(2023)}]{huang2023visual}
Chenguang Huang, Oier Mees, Andy Zeng, and Wolfram Burgard.
\newblock Visual language maps for robot navigation.
\newblock In \emph{IEEE International Conference on Robotics and Automation (ICRA)}, 2023.

\bibitem[{Hughes et~al.(2022)}]{hughes2022hydra}
Nathan Hughes, Yun Chang, and Luca Carlone.
\newblock Hydra: A real-time spatial perception system for 3D scene graph construction and optimization.
\newblock In \emph{Robotics: Science and Systems (RSS)}, 2022.

\bibitem[{Ichter et~al.(2023)}]{ichter2023do}
Brian Ichter, Anthony Brohan, Yevgen Chebotar, Chelsea Finn, Karol Hausman, Alexander Herzog, Daniel Ho, Julian Ibarz, et~al.
\newblock Do as {I} can, not as {I} say: Grounding language in robotic affordances.
\newblock In \emph{Conference on Robot Learning (CoRL)}, volume 205 of \emph{Proceedings of Machine Learning Research}, pages 287--318, 2023.

\bibitem[{Jia and Yu(2026)}]{jia2026hro}
Luyuan Jia and Yinfeng Yu.
\newblock {HRO}: Hierarchical room-to-object framework for zero-shot object goal navigation with large language models.
\newblock In \emph{IEEE International Conference on Systems, Man, and Cybernetics (SMC)}, 2026.

\bibitem[{Khandelwal et~al.(2022)}]{khandelwal2022simple}
Apoorv Khandelwal, Luca Weihs, Roozbeh Mottaghi, and Aniruddha Kembhavi.
\newblock Simple but effective: CLIP embeddings for embodied AI.
\newblock In \emph{Proceedings of the IEEE/CVF Conference on Computer Vision and Pattern Recognition (CVPR)}, 2022.

\bibitem[{Kirillov et~al.(2023)}]{kirillov2023segment}
Alexander Kirillov, Eric Mintun, Nikhila Ravi, Hanzi Mao, Chloe Rolland, Laura Gustafson, Tete Xiao, Spencer Whitehead, et~al.
\newblock Segment anything.
\newblock In \emph{IEEE/CVF International Conference on Computer Vision (ICCV)}, pages 3992--4003, 2023.

\bibitem[{Kollar and Roy(2009)}]{kollar2009utilizing}
Thomas Kollar and Nicholas Roy.
\newblock Utilizing object-object and object-scene context when planning to find things.
\newblock In \emph{IEEE International Conference on Robotics and Automation (ICRA)}, 2009.

\bibitem[{Kolve et~al.(2017)}]{kolve2017ai2}
Eric Kolve, Roozbeh Mottaghi, Winson Han, Eli VanderBilt, Luca Weihs, Alvaro Herrasti, Matt Deitke, Kiana Ehsani, et~al.
\newblock {AI2-THOR}: An interactive {3D} environment for visual {AI}.
\newblock \emph{arXiv preprint arXiv:1712.05474}, 2017.

\bibitem[{Kuang et~al.(2024)}]{kuang2024openfmnav}
Yuxuan Kuang, Hai Lin, and Meng Jiang.
\newblock {OpenFMNav}: Towards open-set zero-shot object navigation via vision-language foundation models.
\newblock In \emph{Findings of the Association for Computational Linguistics: NAACL 2024}, 2024.

\bibitem[{Li et~al.(2023)}]{li2023blip}
Junnan Li, Dongxu Li, Silvio Savarese, and Steven Hoi.
\newblock {BLIP-2}: Bootstrapping language-image pre-training with frozen image encoders and large language models.
\newblock In \emph{International Conference on Machine Learning (ICML)}, volume 202 of \emph{Proceedings of Machine Learning Research}, pages 19730--19742, 2023.

\bibitem[{Li et~al.(2025)}]{li2025goal}
Badi Li, Ren-Jie Lu, Yu Zhou, Jingke Meng, and Wei-Shi Zheng.
\newblock Distilling {LLM} prior to flow model for generalizable agent's imagination in object goal navigation.
\newblock In \emph{Advances in Neural Information Processing Systems (NeurIPS)}, 2025.

\bibitem[{Liu et~al.(2024)}]{liu2024grounding}
Shilong Liu, Zhaoyang Zeng, Tianhe Ren, Feng Li, Hao Zhang, Jie Yang, Qing Jiang, Chunyuan Li, et~al.
\newblock Grounding {DINO}: Marrying {DINO} with grounded pre-training for open-set object detection.
\newblock In \emph{European Conference on Computer Vision (ECCV)}, pages 38--55, 2024.

\bibitem[{Liu et~al.(2026)}]{liu2026global}
Ruimeng Liu, Xinhang Xu, Shenghai Yuan, and Lihua Xie.
\newblock Global planning for object navigation via a weighted traveling repairman problem formulation.
\newblock In \emph{IEEE International Conference on Robotics and Automation (ICRA)}, 2026.

\bibitem[{Long et~al.(2024)}]{long2024instructnav}
Yuxing Long, Wenzhe Cai, Hongcheng Wang, Guanqi Zhan, and Hao Dong.
\newblock InstructNav: Zero-shot system for generic instruction navigation in unexplored environment.
\newblock In \emph{Proceedings of the 8th Conference on Robot Learning (CoRL)}, volume 270 of \emph{Proceedings of Machine Learning Research}, 2024.

\bibitem[{Loo et~al.(2025)}]{loo2025open}
Joel Loo, Zhanxin Wu, and David Hsu.
\newblock Open scene graphs for open-world object-goal navigation.
\newblock \emph{The International Journal of Robotics Research (IJRR)}, 2025.

\bibitem[{Luo et~al.(2022)}]{luo2022stubborn}
Haokuan Luo, Albert Yue, Zhang-Wei Hong, and Pulkit Agrawal.
\newblock Stubborn: A strong baseline for indoor object navigation.
\newblock In \emph{IEEE/RSJ International Conference on Intelligent Robots and Systems (IROS)}, 2022.

\bibitem[{Lyu et~al.(2026)}]{lyu2026rpt}
Yunpeng Lyu, Chao Cao, Ji Zhang, Howie Choset, and Zhongqiang Ren.
\newblock RPT*: Global planning with probabilistic terminals for target search in complex environments.
\newblock \emph{arXiv preprint arXiv:2601.12701}, 2026.

\bibitem[{Majumdar et~al.(2022)}]{majumdar2022zson}
Arjun Majumdar, Gunjan Aggarwal, Bhavika Devnani, Judy Hoffman, and Dhruv Batra.
\newblock ZSON: Zero-shot object-goal navigation using multimodal goal embeddings.
\newblock In \emph{Advances in Neural Information Processing Systems (NeurIPS)}, 2022.

\bibitem[{Mohammadi et~al.(2025)}]{mohammadi2025more}
Mohammad Mohammadi, Daniel Honerkamp, Martin B{\"u}chner, Matteo Cassinelli, Tim Welschehold, Fabien Despinoy, Igor Gilitschenski, and Abhinav Valada.
\newblock MORE: Mobile manipulation rearrangement through grounded language reasoning.
\newblock In \emph{IEEE/RSJ International Conference on Intelligent Robots and Systems (IROS)}, 2025.

\bibitem[{Nie et~al.(2025)}]{nie2025wmnav}
Dujun Nie, Xianda Guo, Yiqun Duan, Ruijun Zhang, and Long Chen.
\newblock {WMNav}: Integrating vision-language models into world models for object goal navigation.
\newblock In \emph{IEEE/RSJ International Conference on Intelligent Robots and Systems (IROS)}, 2025.

\bibitem[{Peng et~al.(2026)}]{peng2025pigeon}
Cheng Peng, Zhenzhe Zhang, Xiaobao Wei, Yanhao Zhang, Heng Wang, Pengwei Wang, Zhongyuan Wang, Cheng Chi, et~al.
\newblock {PIGEON}: {VLM}-driven object navigation via points of interest selection.
\newblock \emph{arXiv preprint arXiv:2511.13207v2}, 2026.

\bibitem[{Pronobis and Jensfelt(2012)}]{pronobis2012large}
Andrzej Pronobis and Patric Jensfelt.
\newblock Large-scale semantic mapping and reasoning with heterogeneous modalities.
\newblock In \emph{IEEE International Conference on Robotics and Automation (ICRA)}, 2012.

\bibitem[{Radford et~al.(2021)}]{radford2021learning}
Alec Radford, Jong~Wook Kim, Chris Hallacy, Aditya Ramesh, Gabriel Goh, Sandhini Agarwal, Girish Sastry, Amanda Askell, et~al.
\newblock Learning transferable visual models from natural language supervision.
\newblock In \emph{International Conference on Machine Learning (ICML)}, volume 139 of \emph{Proceedings of Machine Learning Research}, pages 8748--8763, 2021.

\bibitem[{Rajvanshi et~al.(2024)}]{rajvanshi2024saynav}
Abhinav Rajvanshi, Karan Sikka, Xiao Lin, Bhoram Lee, Han-Pang Chiu, and Alvaro Velasquez.
\newblock {SayNav}: Grounding large language models for dynamic planning to navigation in new environments.
\newblock In \emph{Proceedings of the International Conference on Automated Planning and Scheduling (ICAPS)}, 2024.

\bibitem[{Ramakrishnan et~al.(2021)}]{ramakrishnan2021habitat}
Santhosh~Kumar Ramakrishnan, Aaron Gokaslan, Erik Wijmans, Oleksandr Maksymets, Alexander Clegg, John Turner, Eric Undersander, Wojciech Galuba, et~al.
\newblock {Habitat-Matterport 3D Dataset (HM3D)}: 1000 large-scale {3D} environments for embodied {AI}.
\newblock In \emph{Proceedings of the Neural Information Processing Systems Track on Datasets and Benchmarks (NeurIPS Datasets and Benchmarks)}, 2021.

\bibitem[{Ramakrishnan et~al.(2022)}]{ramakrishnan2022poni}
Santhosh~Kumar Ramakrishnan, Devendra~Singh Chaplot, Ziad Al-Halah, Jitendra Malik, and Kristen Grauman.
\newblock {PONI}: Potential functions for {ObjectGoal} navigation with interaction-free learning.
\newblock In \emph{Proceedings of the IEEE/CVF Conference on Computer Vision and Pattern Recognition (CVPR)}, 2022.

\bibitem[{Ramrakhya et~al.(2022)}]{ramrakhya2022habitat}
Ram Ramrakhya, Eric Undersander, Dhruv Batra, and Abhishek Das.
\newblock {Habitat-Web}: Learning embodied object-search strategies from human demonstrations at scale.
\newblock In \emph{Proceedings of the IEEE/CVF Conference on Computer Vision and Pattern Recognition (CVPR)}, 2022.

\bibitem[{Ramrakhya et~al.(2023)}]{ramrakhya2023pirlnav}
Ram Ramrakhya, Dhruv Batra, Erik Wijmans, and Abhishek Das.
\newblock PIRLNav: Pretraining with imitation and RL finetuning for ObjectNav.
\newblock In \emph{Proceedings of the IEEE/CVF Conference on Computer Vision and Pattern Recognition (CVPR)}, 2023.

\bibitem[{Rana et~al.(2023)}]{rana2023sayplan}
Krishan Rana, Jesse Haviland, Sourav Garg, Jad Abou-Chakra, Ian Reid, and Niko Suenderhauf.
\newblock SayPlan: Grounding large language models using 3D scene graphs for scalable robot task planning.
\newblock In \emph{Conference on Robot Learning (CoRL)}, volume 229 of \emph{Proceedings of Machine Learning Research}, 2023.

\bibitem[{Rosinol et~al.(2020)}]{rosinol20203d}
Antoni Rosinol, Arjun Gupta, Marcus Abate, Jingnan Shi, and Luca Carlone.
\newblock 3D dynamic scene graphs: Actionable spatial perception with places, objects, and humans.
\newblock In \emph{Robotics: Science and Systems (RSS)}, 2020.

\bibitem[{Savva et~al.(2019)}]{savva2019habitat}
Manolis Savva, Abhishek Kadian, Oleksandr Maksymets, Yili Zhao, Erik Wijmans, Bhavana Jain, Julian Straub, Jia Liu, et~al.
\newblock Habitat: A platform for embodied {AI} research.
\newblock In \emph{Proceedings of the IEEE/CVF International Conference on Computer Vision (ICCV)}, 2019.

\bibitem[{Sethian(1996)}]{sethian1996fast}
James~A. Sethian.
\newblock A fast marching level set method for monotonically advancing fronts.
\newblock \emph{Proceedings of the National Academy of Sciences (PNAS)}, 93(4):1591--1595, 1996.

\bibitem[{Shah et~al.(2022)}]{shah2022lm}
Dhruv Shah, B{\l}a{\.z}ej Osi{\'n}ski, Brian Ichter, and Sergey Levine.
\newblock {LM-Nav}: Robotic navigation with large pre-trained models of language, vision, and action.
\newblock In \emph{Conference on Robot Learning (CoRL)}, 2022.

\bibitem[{Shah et~al.(2023)}]{shah2023navigation}
Dhruv Shah, Michael~Robert Equi, B{\l}a{\.z}ej Osi{\'n}ski, Fei Xia, Brian Ichter, and Sergey Levine.
\newblock Navigation with large language models: Semantic guesswork as a heuristic for planning.
\newblock In \emph{Conference on Robot Learning (CoRL)}, 2023.

\bibitem[{Stone(1975)}]{stone1975theory}
Lawrence~D. Stone.
\newblock \emph{Theory of Optimal Search}, volume 118 of \emph{Mathematics in Science and Engineering}.
\newblock Academic Press, New York, 1975.

\bibitem[{Sun et~al.(2024)}]{sun2024prioritized}
Xinyu Sun, Lizhao Liu, Hongyan Zhi, Ronghe Qiu, and Junwei Liang.
\newblock Prioritized semantic learning for zero-shot instance navigation.
\newblock In \emph{European Conference on Computer Vision (ECCV)}, 2024.

\bibitem[{Szot et~al.(2021)}]{szot2021habitat}
Andrew Szot, Alexander Clegg, Eric Undersander, Erik Wijmans, Yili Zhao, John Turner, Noah Maestre, Mustafa Mukadam, et~al.
\newblock {Habitat 2.0}: Training home assistants to rearrange their habitat.
\newblock In \emph{Advances in Neural Information Processing Systems (NeurIPS)}, 2021.

\bibitem[{Unlu et~al.(2024)}]{unlu2024reproducibility}
Halil~Utku Unlu, Shuaihang Yuan, Congcong Wen, Hao Huang, Anthony Tzes, and Yi Fang.
\newblock Reliable semantic understanding for real world zero-shot object goal navigation.
\newblock \emph{arXiv preprint arXiv:2410.21926}, 2024.

\bibitem[{Vincent and Soille(1991)}]{vincent1991watersheds}
Luc Vincent and Pierre Soille.
\newblock Watersheds in digital spaces: An efficient algorithm based on immersion simulations.
\newblock \emph{IEEE Transactions on Pattern Analysis and Machine Intelligence (TPAMI)}, 13(6):583--598, 1991.

\bibitem[{Wen et~al.(2024)}]{vltnet2024}
Congcong Wen, Yisiyuan Huang, Hao Huang, Yanjia Huang, Shuaihang Yuan, Yu Hao, Hui Lin, Yu-Shen Liu, et~al.
\newblock Zero-shot object navigation with vision-language models reasoning.
\newblock In \emph{International Conference on Pattern Recognition (ICPR)}, pages 389--404, 2024.

\bibitem[{Wei et~al.(2022)}]{wei2022chain}
Jason Wei, Xuezhi Wang, Dale Schuurmans, Maarten Bosma, Brian Ichter, Fei Xia, Ed~H. Chi, Quoc~V. Le, et~al.
\newblock Chain-of-thought prompting elicits reasoning in large language models.
\newblock In \emph{Advances in Neural Information Processing Systems (NeurIPS)}, 2022.

\bibitem[{Werby et~al.(2024)}]{werby2024hierarchical}
Abdelrhman Werby, Chenguang Huang, Martin B{\"u}chner, Abhinav Valada, and Wolfram Burgard.
\newblock Hierarchical open-vocabulary 3D scene graphs for language-grounded robot navigation.
\newblock In \emph{Robotics: Science and Systems (RSS)}, 2024.

\bibitem[{Wu et~al.(2024)}]{wu2024voronav}
Pengying Wu, Yao Mu, Bingxian Wu, Yi~Hou, Ji~Ma, Shanghang Zhang, and Chang Liu.
\newblock {VoroNav}: Voronoi-based zero-shot object navigation with large language model.
\newblock In \emph{International Conference on Machine Learning (ICML)}, 2024.

\bibitem[{Xia et~al.(2018)}]{xia2018gibson}
Fei Xia, Amir~R. Zamir, Zhiyang He, Alexander Sax, Jitendra Malik, and Silvio Savarese.
\newblock Gibson {Env}: Real-world perception for embodied agents.
\newblock In \emph{Proceedings of the IEEE Conference on Computer Vision and Pattern Recognition (CVPR)}, 2018.

\bibitem[{Yadav et~al.(2023)}]{yadav2023habitat}
Karmesh Yadav, Ram Ramrakhya, Santhosh~Kumar Ramakrishnan, Theo Gervet, John Turner, Aaron Gokaslan, Noah Maestre, Angel~Xuan Chang, et~al.
\newblock {Habitat-Matterport 3D} semantics dataset.
\newblock In \emph{Proceedings of the IEEE/CVF Conference on Computer Vision and Pattern Recognition (CVPR)}, 2023.

\bibitem[{Yadav et~al.(2023)}]{yadav2023ovrl}
Karmesh Yadav, Arjun Majumdar, Ram Ramrakhya, Naoki Yokoyama, Alexei Baevski, Zsolt Kira, Oleksandr Maksymets, and Dhruv Batra.
\newblock OVRL-V2: A simple state-of-art baseline for ImageNav and ObjectNav.
\newblock \emph{arXiv preprint arXiv:2303.07798}, 2023.

\bibitem[{Yamauchi(1997)}]{yamauchi1997frontier}
Brian Yamauchi.
\newblock A frontier-based approach for autonomous exploration.
\newblock In \emph{Proceedings of the IEEE International Symposium on Computational Intelligence in Robotics and Automation (CIRA)}, pages 146--151, 1997.

\bibitem[{Yang et~al.(2024)}]{yang2024qwen25}
An~Yang, Baosong Yang, Beichen Zhang, Binyuan Hui, Bo~Zheng, Bowen Yu, Chengyuan Li, Dayiheng Liu, et~al.
\newblock {Qwen2.5} technical report.
\newblock \emph{arXiv preprint arXiv:2412.15115}, 2024.

\bibitem[{Yin et~al.(2024)}]{yin2024sg}
Hang Yin, Xiuwei Xu, Zhenyu Wu, Jie Zhou, and Jiwen Lu.
\newblock SG-Nav: Online 3D scene graph prompting for LLM-based zero-shot object navigation.
\newblock In \emph{Advances in Neural Information Processing Systems (NeurIPS)}, 2024.

\bibitem[{Yin et~al.(2025)}]{yin2025unigoal}
Hang Yin, Xiuwei Xu, Lingqing Zhao, Ziwei Wang, Jie Zhou, and Jiwen Lu.
\newblock UniGoal: Towards universal zero-shot goal-oriented navigation.
\newblock In \emph{Proceedings of the IEEE/CVF Conference on Computer Vision and Pattern Recognition (CVPR)}, 2025.

\bibitem[{Yokoyama et~al.(2024)}]{yokoyama2024vlfm}
Naoki Yokoyama, Sehoon Ha, Dhruv Batra, Jiuguang Wang, and Bernadette Bucher.
\newblock VLFM: Vision-language frontier maps for zero-shot semantic navigation.
\newblock In \emph{IEEE International Conference on Robotics and Automation (ICRA)}, 2024.

\bibitem[{Yu et~al.(2023)}]{yu2023l3mvn}
Bangguo Yu, Hamidreza Kasaei, and Ming Cao.
\newblock L3MVN: Leveraging large language models for visual target navigation.
\newblock In \emph{IEEE/RSJ International Conference on Intelligent Robots and Systems (IROS)}, 2023.

\bibitem[{Yu et~al.(2024)}]{yu2024trajectory}
Xinyao Yu, Sixian Zhang, Xinhang Song, Xiaorong Qin, and Shuqiang Jiang.
\newblock Trajectory diffusion for {ObjectGoal} navigation.
\newblock In \emph{Advances in Neural Information Processing Systems (NeurIPS)}, 2024.

\bibitem[{Yuan et~al.(2024)}]{yuan2024defa}
Shuaihang Yuan, Halil~Utku Unlu, Hao Huang, Congcong Wen, Anthony Tzes, and Yi Fang.
\newblock Exploring the reliability of foundation model-based frontier selection in zero-shot object goal navigation.
\newblock In \emph{International Conference on Pattern Recognition (ICPR)}, pages 119--134, 2024.

\bibitem[{Yuan et~al.(2024)}]{yuan2024gamap}
Shuaihang Yuan, Hao Huang, Yu~Hao, Congcong Wen, Anthony Tzes, and Yi~Fang.
\newblock {GAMap}: Zero-shot object goal navigation with multi-scale geometric-affordance guidance.
\newblock In \emph{Advances in Neural Information Processing Systems (NeurIPS)}, 2024.

\bibitem[{Zhang et~al.(2023)}]{zhang20233d}
Jiazhao Zhang, Liu Dai, Fanpeng Meng, Qingnan Fan, Xuelin Chen, Kai Xu, and He~Wang.
\newblock 3D-aware object goal navigation via simultaneous exploration and identification.
\newblock In \emph{Proceedings of the IEEE/CVF Conference on Computer Vision and Pattern Recognition (CVPR)}, 2023.

\bibitem[{Zhang et~al.(2023)}]{zhang2023faster}
Chaoning Zhang, Dongshen Han, Yu~Qiao, Jung~Uk Kim, Sung-Ho Bae, Seungkyu Lee, and Choong~Seon Hong.
\newblock Faster segment anything: Towards lightweight {SAM} for mobile applications.
\newblock \emph{arXiv preprint arXiv:2306.14289}, 2023.

\bibitem[{Zhang et~al.(2024)}]{zhang2024imagine}
Sixian Zhang, Xinyao Yu, Xinhang Song, Xiaohan Wang, and Shuqiang Jiang.
\newblock Imagine before go: Self-supervised generative map for object goal navigation.
\newblock In \emph{Proceedings of the IEEE/CVF Conference on Computer Vision and Pattern Recognition (CVPR)}, pages 16414--16425, 2024.

\bibitem[{Zhang et~al.(2024)}]{zhang2024trihelper}
Lingfeng Zhang, Qiang Zhang, Hao Wang, Erjia Xiao, Zixuan Jiang, Honglei Chen, and Renjing Xu.
\newblock {TriHelper}: Zero-shot object navigation with dynamic assistance.
\newblock In \emph{IEEE/RSJ International Conference on Intelligent Robots and Systems (IROS)}, 2024.

\bibitem[{Zhang et~al.(2025)}]{zhang2025apexnav}
Mingjie Zhang, Yuheng Du, Chengkai Wu, Jinni Zhou, Zhenchao Qi, Jun Ma, and Boyu Zhou.
\newblock ApexNav: An adaptive exploration strategy for zero-shot object navigation with target-centric semantic fusion.
\newblock \emph{IEEE Robotics and Automation Letters (RA-L)}, 10(11):11530--11537, 2025.

\bibitem[{Zhang et~al.(2025)}]{zhang2025falcon}
Yichen Zhang, Xinyi Chen, Chen Feng, Boyu Zhou, and Shaojie Shen.
\newblock {FALCON}: Fast autonomous aerial exploration using coverage path guidance.
\newblock \emph{IEEE Transactions on Robotics (T-RO)}, 41:1365--1385, 2025.

\bibitem[{Zhong et~al.(2025)}]{zhong2024topvnav}
Linqing Zhong, Chen Gao, Zihan Ding, Yue Liao, Huimin Ma, Shifeng Zhang, Xu Zhou, and Si Liu.
\newblock {TopV-Nav}: Unlocking the top-view spatial reasoning potential of {MLLM} for zero-shot object navigation.
\newblock \emph{arXiv preprint arXiv:2411.16425v2}, 2025.

\bibitem[{Zhou et~al.(2023)}]{zhou2023esc}
Kaiwen Zhou, Kaizhi Zheng, Connor Pryor, Yilin Shen, Hongxia Jin, Lise Getoor, and Xin~Eric Wang.
\newblock ESC: Exploration with soft commonsense constraints for zero-shot object navigation.
\newblock In \emph{Proceedings of the 40th International Conference on Machine Learning (ICML)}, volume 202 of \emph{Proceedings of Machine Learning Research}, 2023.

\bibitem[{Zhou et~al.(2024)}]{zhou2024navgpt}
Gengze Zhou, Yicong Hong, and Qi Wu.
\newblock {NavGPT}: Explicit reasoning in vision-and-language navigation with large language models.
\newblock In \emph{Proceedings of the AAAI Conference on Artificial Intelligence (AAAI)}, volume~38, pages 7641--7649, 2024.

\bibitem[{Zhou et~al.(2025)}]{zhou2025beliefmapnav}
Zibo Zhou, Yue Hu, Lingkai Zhang, Zonglin Li, and Siheng Chen.
\newblock {BeliefMapNav}: {3D} voxel-based belief map for zero-shot object navigation.
\newblock In \emph{Advances in Neural Information Processing Systems (NeurIPS)}, 2025.

\end{thebibliography}

\clearpage
\appendix

% ============================================================================
% 11_appendix_algorithms
% ============================================================================
% ============================================================================
%  Appendix A: sub-routines of Algorithm 1 and the hyper-parameter table.
% ============================================================================

\section{Sub-routines}\label{app:algorithms}

This appendix states the sub-routines called by \cref{alg:main} in enough detail to re-implement them, with the default values used in all our experiments; \cref{tab:hyperparams} collects these values in one place. Notation follows \cref{sec:problem}: the scene graph is $\G = (\M, \Ob, \D, \R)$, $x_t$ is the pose after $t$ actions, and every time bound is counted in actions, the unit in which the benchmarks charge the budget. \textsc{Navigate} (\cref{alg:navigate}) and \textsc{ExploreLocal} (\cref{alg:explorelocal}) execute at most one action per call (a call that ends a visit leaves the action to the next decision), and \textsc{UpdateSceneGraph} (\cref{alg:updatesg}) runs after every action, so no decision is taken on a stale map; during target closing (\cref{alg:targetclosing}) it updates every layer except the rooms. The two language models are queried only at the points named in \cref{alg:classifyrooms,alg:llmprobabilities}, and no action is emitted while a reply is pending.

\paragraph{Scene graph.}
\Cref{alg:updatesg} maintains the four layers of $\G$. Three choices shape it. The occupancy grid is never inflated: clearance is added by the planner, so one map serves both the preferred 0.30\,m and the minimum 0.18\,m clearance. Objects are fused per class only and confirmed after two distinct frames, so a single spurious box never reaches the room classifier or the oracle. Rooms are basins of a watershed~\citep{vincent1991watersheds} on the clearance field of the observed free space, cut at confirmed doorways and at narrow passages the agent has walked through; we write $\Theta$ for this set of walked-through thresholds, which lets a room boundary appear where the detector never confirmed a door.

\begin{algorithm}[htbp]
\DontPrintSemicolon
\caption{\textsc{UpdateSceneGraph}$(\G, I_t, D_t, x_t)$}
\label{alg:updatesg}
\Input{scene graph $\G = (\M, \Ob, \D, \R)$; RGB image $I_t$; depth image $D_t$; pose $x_t$}
\Output{updated $\G$; called after every action, never calls an LLM}
\BlankLine
\Comment{Occupancy grid $\M$ (\textsc{IntegrateDepth}): 0.10\,m cells, log-odds, not inflated}
\ForEach{depth pixel on every 8th row and column with $0.5 \le D_t \le 5.0$\,m}{
  back-project it with $x_t$ to a world point $X$\;
  \uIf{$X$ lies 0.15--0.93\,m above the floor}{add an obstacle hit at $X$; free the cells on the ray to $X$\;}
  \lElseIf{$X$ lies within 0.10\,m of the floor}{free the cell of $X$ unless it is occupied}
}
cell state: free if $P_{\mathrm{occ}} \le 0.45$, occupied if $P_{\mathrm{occ}} \ge 0.65$, unknown otherwise\;
\Comment{Objects $\Ob$}
$B_0 \gets$ \textsc{YOLO-World} boxes on $I_t$ (30 prompts) with confidence $\ge 0.05$\;
$B \gets$ boxes of $B_0$ with confidence $\ge 0.30$, minus any box overlapping a stronger one of its class (IoU $\ge 0.5$) or of any class (IoU $\ge 0.9$)\;
\ForEach{box $b \in B$ not of a door or stair prompt}{
  $d^{*} \gets$ 30th percentile of the valid depth in the central half of $b$\;
  $S \gets$ largest connected set of central pixels with depth near $d^{*}$ \Comment*[r]{reject $b$ if $S$ holds $< 35\%$}
  $X_b \gets$ per-axis median of $S$ back-projected with $x_t$;\quad $\rho_b \gets$ footprint radius in $[0.15, 1.5]$\,m\;
  \lIf{$X_b$ is outside $-0.30$ to $+2.2$\,m of the floor height \KwOr lacks floor support}{\KwContinue}
  $o \gets$ nearest object of the same class, not yet fed by this frame, with centres $\le 0.35$\,m apart or footprint IoU $\ge 0.25$, and heights within 0.35\,m\;
  \uIf{$o$ exists}{update its centre and radius by a running mean; $n_o \gets n_o + 1$\;}
  \lElse{add a candidate object $o$ with $n_o \gets 1$}
}
only confirmed objects ($n_o \ge 2$ distinct frames) enter $\G$\;
\Comment{Doorways $\D$ and walked-through thresholds $\Theta$}
\ForEach{box of $B_0$ from a door prompt, not overlapped (IoU $\ge 0.6$) by a non-door box of confidence $\ge 0.5$}{
  read its depth on the jambs and lintel; keep it if 0.35--1.5\,m wide, with its top $\ge 1.4$\,m and its bottom $\le 0.5$\,m above the floor\;
}
$\D \gets$ doors seen $\ge 3$ times, once from $\ge 0.2$\,m or $\ge 20^\circ$ away from the first view\;
$\Theta \gets \Theta\, \cup$ clearance dips on the agent's trail ($\le 0.55$\,m, with $\ge 0.30$\,m more on both sides)\;
\Comment{Rooms $\R$ (\textsc{SegmentRooms}); skipped during target closing}
$F \gets$ free cells of $\M$ minus seen staircases, healed by one closing and one opening\;
$\Phi \gets$ Euclidean clearance of $F$\;
cut $F$ at each door of $\D$ (the narrowest passage within 0.9\,m, else a 0.75\,m disc) and at each threshold of $\Theta$\;
flood $-\Phi$ from its local maxima that are $\ge 1.0$\,m apart with $\Phi \ge 0.3$\,m (watershed); give cut cells to the nearest room\;
merge adjacent basins whose saddle lies $< 0.2$\,m below the shallower peak, never across a cut\;
drop rooms smaller than 0.5\,m$^2$; keep room ids by birth anchor, IoU $\ge 0.15$ or containment $\ge 0.6$\;
\ForEach{room $r \in \R$}{
  $\Ob_r \gets$ confirmed objects on $r$ or within their footprint radius $+\,0.6$\,m of it\;
  record area $A_r$, accessible frontier clusters, time searched, last time inside and scan status\;
}
\KwRet{$\G$}
\end{algorithm}

\paragraph{Room labels.}
\Cref{alg:classifyrooms} is the only use of the small model (Qwen2.5-3B-Instruct~\citep{yang2024qwen25}). It is not called on ordinary steps: rooms are relabelled at a decision point when the multiset of their object classes has changed or after a 50-action refresh, and mid-visit only for the room being searched or crossed when a new kind of object appears in it. A label is \emph{firm} only when the model is confident and the evidence is distinctive: two object kinds of which at least one is not generic, or one signature object that agrees with the label (a bed for a bedroom; a toilet, shower or bathtub for a bathroom; an oven, stove or refrigerator for a kitchen; the
Gibson vocabulary has no stove prompt). Rooms holding only generic objects (cabinet, potted plant, book, vase, clock, cup, bottle) stay \emph{unknown} without a call. Tentative labels reach the oracle marked with a question mark.

\begin{algorithm}[htbp]
\DontPrintSemicolon
\caption{\textsc{ClassifyRooms}$(\G)$}
\label{alg:classifyrooms}
\Input{scene graph $\G$ with rooms $\R$ and confirmed objects $\Ob_r$; trigger (decision point or mid-visit clue)}
\Output{$\G$ with a label $\ell_r$, confidence $\kappa_r$ and strength $\beta_r \in \{\text{tentative}, \text{firm}\}$ per room}
\BlankLine
\uIf(\Comment*[f]{just before the oracle is queried}){called at a decision point}{
  $Q \gets \{ r \in \R :$ the multiset of classes in $\Ob_r$ changed, \KwOr $r$ was labelled $\ge 50$ actions ago$\}$\;
}
\Else(\Comment*[f]{mid-visit clue}){
  $Q \gets$ the room being searched or crossed, if a new kind of object appeared in it\;
}
\ForEach{$r \in Q$}{
  \If(\Comment*[f]{no LLM call}){$\Ob_r = \varnothing$ \KwOr every object in $\Ob_r$ is generic}{
    $\ell_r \gets$ \emph{unknown};\quad $\beta_r \gets$ tentative;\quad \KwContinue\;
  }
  \uIf{an identical object signature was judged before}{reuse that verdict $(\ell_r, \kappa_r)$\;}
  \Else{
    $(\ell_r, \kappa_r) \gets$ \LLM{$\Ob_r$} \Comment*[r]{3B model, temperature 0, JSON; ``class $\times$ count''}
    \lIf{$\ell_r$ is not in the fixed label set}{$\ell_r \gets$ \emph{unknown}}
  }
  $\beta_r \gets$ firm \textbf{iff} $\kappa_r \ge 0.65$, $\ell_r \neq$ \emph{unknown}, and ($\ge 2$ object kinds with one non-generic \KwOr a signature object of type $\ell_r$ in $\Ob_r$)\;
}
\Comment{if the call fails, held labels stand and the agent explores during a back-off (\cref{alg:llmprobabilities})}
\KwRet{$\G$}
\end{algorithm}

\paragraph{Candidate nodes.}
\Cref{alg:candidatenodes} builds the solver's instance $\N$ from the three node types of \cref{sec:planning}. \Cref{alg:main} first queries the oracle on the described set $\mathcal{U}$: every unfinished room of the storey together with its status (and the room whose visit is in force), every opening not yet peeked and not folded into such a room, and the sighted targets. \textsc{CandidateNodes} then forms $\N$ from $\mathcal{U}$, and the lines marked $\star$ use the probabilities of the reply, so that every eligibility rule appears in one place. Finished rooms are therefore withheld and never offered again, unless an optional second pass, off by default, re-offers them once the storey is covered; the pass verdict of the reply is recorded but does not open it. Openings are clusters of reachable frontier cells~\citep{yamauchi1997frontier} that lead out of the mapped floor. The cap of 12 nodes keeps the exact solver well inside its time budget, and reserving one slot for the nearest node keeps a cheap option in every instance even when the oracle prefers distant rooms.

\begin{algorithm}[htbp]
\DontPrintSemicolon
\caption{\textsc{CandidateNodes}$(\G, p, \omega)$}
\label{alg:candidatenodes}
\Input{scene graph $\G$; target category $\cstar$; recorded scan points and completed peeks; pass verdict $\omega$ and probabilities $p$ of the oracle (\cref{alg:llmprobabilities})}
\Output{nodes $\N$ offered to RPT$^*$, each with an entry point $e_i$}
\BlankLine
\Comment{Rooms}
\ForEach{room $r \in \R$}{
  $r$ is \emph{finished} if a completed scan point (of the warm-up or a full path glance) lies in it, a scan saw $\ge 50\%$ of its cells within 5\,m with no live frontier left, it was stood in $\ge 0.5$\,m from its edge (anywhere, if it is too narrow), looked into (half of it seen from one pose) or walked through while narrow (widest clearance $\le 0.9$\,m) with no live frontier left, or it is a doorless fragment $< 3$\,m$^2$ with no live frontier; a live frontier is a frontier cluster of $r$, not under the agent or beside a retired goal, whose view would reveal $\ge 1.0$\,m$^2$ of unknown floor; a verdict is judged again once $r$ grows beyond $1.5\times$ its size at the verdict\;
  \lIf(\Comment*[f]{e.g.\ a sink, shower or bathtub, for a toilet}){$\Ob_r$ holds an object that usually shares a room with $\cstar$}{$r$ is finished only by a completed scan point in it}
}
$\N_{\mathrm{room}} \gets \{ r \in \R : r \text{ not finished} \}$\;
\lIf(\Comment*[f]{second pass, off by default}){the second pass is enabled \KwAnd no unfinished room, opening or reachable frontier revealing $\ge 1.0$\,m$^2$ is left}{$\N_{\mathrm{room}} \gets \R$}
\ForEach{$r \in \N_{\mathrm{room}}$}{
  $e_r \gets$ the live frontier of $r$ that reveals the most unknown floor per metre walked; else its doorstep, the passable cell of $r$ nearest $x_t$ along the floor; else, and first on a revisit, the vantage point: the reachable cell of $r$ farthest from everything outside $r$, with clearance $\ge 0.4$\,m and, on a revisit, $\ge 1.5$\,m from earlier scan points; else the centroid of $r$\;
}
\Comment{Openings onto unseen space}
$\N_{\mathrm{open}} \gets$ reachable frontier exits merged within 1.5\,m and $\ge 0.5$\,m wide, except those within 1.2\,m of the agent, in an object's shadow, with a seen staircase or drop within 1.5\,m on their unknown side, within 1.0\,m of a finished peek, or leading out of a room of $\N_{\mathrm{room}}$ not valued below $p_{\min}$ (folded into that room's visit) \Comment*[r]{$\star$}
\lForEach{$i \in \N_{\mathrm{open}}$}{label $i$ \emph{doorway} if a door of $\D$ lies within 1.0\,m, else \emph{gap};\quad $e_i \gets$ its threshold}
\Comment{Sighted but unverified targets}
$\N_{\mathrm{target}} \gets$ objects of class $\cstar$, confirmed or seen once with confidence $\ge 0.50$, not yet inspected, away from candidates disproved from close\;
\lForEach{$i \in \N_{\mathrm{target}}$}{$e_i \gets$ nearest reachable point on a 1.5\,m (else 2.25\,m) ring about $i$}
$\N \gets \N_{\mathrm{room}} \cup \N_{\mathrm{open}} \cup \N_{\mathrm{target}}$ \Comment*[r]{stair nodes only in multi-storey mode}
\Comment{Filters}
drop $i$ if no passable cell lies within 2.0\,m of $e_i$, or if it is unreachable from $x_t$\;
drop $i$ for 180 actions after its 3rd failed transit (unreachable, timed out or blocked)\;
drop $i$ for 120 actions after a visit to it ended (30 after a failed transit), unless that visit only spent its budget \Comment*[r]{cooldown}
\lIf{every remaining node is cooling}{readmit those whose visit did not exhaust them}
\lIf(\Comment*[f]{$\star$}){$i \in \N_{\mathrm{open}}$ was not scored by the reply (e.g.\ it was folded when the prompt was built)}{$p_i \gets$ its last score, else 0.10}
$p_i \gets f_i\, p_i$: $f_i = \min(1, \max(0.25, a_i / 4\,\mathrm{m}^2))$ for a gap whose threshold view reveals $a_i$ of unknown floor, 0.5 for a room stood in, left and unfinished (0.2 after $\ge 20$ actions in it with $\le 1$ live frontier left), else 1 \Comment*[r]{$\star$}
drop $i$ if $p_i < p_{\min} = 0.05$ \Comment*[r]{$\star$}
\lIf(\Comment*[f]{$\star$}){$|\N| > 12$}{keep the nearest node and the 11 most probable}
\KwRet{$\N$}
\end{algorithm}

\paragraph{Probabilities.}
\Cref{alg:llmprobabilities} queries the larger model (Qwen2.5-14B-Instruct~\citep{yang2024qwen25}) once per decision point, with at most one schema-repair retry. The prompt never contains a distance (\cref{sec:planning}); the flag marking the room the agent stands in is shown, with the instruction never to use it as a reason. The reply is ordered so that the model states its reasoning (what kind of home this is, what the house holds so far, how far the search has progressed) before any number, and a short justification before each probability, in the spirit of chain-of-thought prompting~\citep{wei2022chain}. Times and the budget in the prompt are coarsened, so an unchanged map yields a byte-identical prompt; because decoding is greedy, the previous reply is then reused without a call. A failed query never degrades into a uniform prior: the agent explores frontiers until the back-off expires. The prompts are reproduced in \cref{app:prompts}.

\begin{algorithm}[htbp]
\DontPrintSemicolon
\caption{\textsc{LLMProbabilities}$(\cstar, \G, \mathcal{U})$}
\label{alg:llmprobabilities}
\Input{target category $\cstar$; scene graph $\G$; described nodes $\mathcal{U}$: every unfinished room of the storey (and the room in force) with its status, the openings not yet peeked or folded, and the sighted targets $\N_{\mathrm{target}}$; previous prompt $P'$ and reply $y'$}
\Output{probabilities $p = (p_i)_{i \in \mathcal{U}}$ and the pass verdict $\omega$, or a back-off}
\BlankLine
$H \gets$ house summary: room types found, rooms unidentified, openings not looked into, whether frontier is left, actions used (in steps of 25)\;
\ForEach(\Comment*[f]{never a distance}){$i \in \mathcal{U} \setminus \N_{\mathrm{target}}$}{
  $\eta_i \gets$ for a room: label (``?'' if tentative), area, accessible frontier count, status and scan recency, time searched, up to 8 object classes seen, and whether the agent is in it; for an opening: doorway or gap, the room it leaves, objects glimpsed through it\;
}
$P \gets$ prompt$(\cstar, H, \{\eta_i\})$\;
\uIf(\Comment*[f]{unchanged map: no call}){$P = P'$}{$y \gets y'$\;}
\Else{
  $y \gets$ \LLM{$P$} \Comment*[r]{14B model, temperature 0, JSON, reasoning first}
  \lIf{$y$ violates the reply schema \KwOr scores no listed node}{$y \gets$ \LLM{$P$, schema-repair instruction}}
  \If(\Comment*[f]{never a uniform prior}){$y$ is still unusable}{
    back off $\min(25 \cdot 2^{k-1}, 200)$ actions after the $k$-th consecutive failure, running \ExploreFrontier{} meanwhile\;
    \KwRet{no probabilities}\;
  }
}
$\omega \gets$ pass verdict of $y$ (\emph{first} if missing) \Comment*[r]{$y$: home, house, stage, pass, nodes}
\ForEach(\Comment*[f]{independent, not normalised}){$i \in \mathcal{U} \setminus \N_{\mathrm{target}}$}{
  $p_i \gets \min(1 - 10^{-6}, \max(0, r_i/100))$ if $y$ gives $i$ a score $r_i$ (0--100), else $p_i \gets 0.01$\;
}
\lForEach(\Comment*[f]{detector evidence; not shown to the LLM}){$i \in \N_{\mathrm{target}}$}{$p_i \gets 0.85$}
\KwRet{$p, \omega$}
\end{algorithm}

\paragraph{Navigation.}
\Cref{alg:navigate} turns an A$^*$ route~\citep{hart1968formal} into one discrete action. Routes are committed: A$^*$ runs again only when the goal moves, the agent drifts off the route, the route becomes blocked on the current map, or progress stalls, which avoids the route flicker that a fresh plan per action produces on a growing map. The aim point lies 0.5\,m ahead along the route and progress along it never decreases, so a later leg passing close to the agent cannot capture it. A forward step that fails to advance is turned into map evidence (the cell ahead becomes occupied) rather than retried.

\begin{algorithm}[htbp]
\DontPrintSemicolon
\caption{\textsc{Navigate}$(\tau)$}
\label{alg:navigate}
\Input{route $\tau$ to the goal $g$, committed by A$^*$; pose $x_t$; progress $s$ along $\tau$}
\Output{one action executed; $x_t$ and $t$ updated}
\BlankLine
\If(\Comment*[f]{otherwise the route is kept}){$g$ moved $> 0.5$\,m, $x_t$ is $> 0.5$\,m off $\tau$, $\tau$ collides on $\M$, \KwOr $s$ grew $< 0.10$\,m in 30 actions}{
  $\tau \gets$ \AStar{$x_t, g$};\quad $s \gets 0$ \Comment*[r]{8-connected, clearance-weighted}
}
$s \gets \max(s, \text{arc length of the projection of } x_t \text{ on } \tau)$ \Comment*[r]{progress never decreases}
$z \gets$ point of $\tau$ at arc length $s + 0.5$\,m \Comment*[r]{$s + 0.25$\,m after a blocked step}
\uIf(\Comment*[f]{tilt benchmarks only}){a commanded camera pitch is at least half a tilt step away}{$a \gets$ the \textsc{look} toward it\;}
\uElseIf{a path glance is due at this point of $\tau$}{$a \gets$ its next in-place turn\;}
\uElseIf(\Comment*[f]{arrived}){$z$ is the last waypoint \KwAnd $\lVert x_t - g \rVert \le 0.25$\,m}{$a \gets$ a turn to the final heading if one is needed, else none\;}
\uElseIf{the heading error to $z$ exceeds $15^\circ$}{$a \gets$ a $30^\circ$ turn toward $z$\;}
\uElseIf(\Comment*[f]{never circle the goal}){$z = g$ \KwAnd a forward step would not bring $x_t$ closer to $g$}{$a \gets$ as on arrival\;}
\lElse{$a \gets$ \textsc{forward} (0.25\,m)}
\lIf(\Comment*[f]{idle action; never STOP}){$a$ is none}{$a \gets$ a left turn}
execute $a$; observe $(I_{t+1}, D_{t+1}, x_{t+1})$;\quad $t \gets t + 1$\;
\If(\Comment*[f]{blocked step}){$a = $ \textsc{forward} advanced $< 0.05$\,m along its heading}{
  mark the cell one step ahead occupied in $\M$ and discard $\tau$\;
}
\Comment{\cref{alg:main} then runs \textsc{UpdateSceneGraph}; a box of class $\cstar$ that passes the takeover gates of \cref{alg:targetclosing} flags ``$\cstar$ sighted''}
\end{algorithm}

\paragraph{Local visits.}
\Cref{alg:explorelocal} implements the bounded local visits. A room is searched by walking to its live frontiers, best first by the unknown floor a look-around from each would reveal per metre walked, until none is left; the ranking is re-read every action, so a frontier the camera resolves on the way drops out, and the cap keeps a large or hard-to-walk room from absorbing the budget. A visit that hits its cap ends as \emph{budget spent} and is not cooled down, so RPT$^*$ may choose the node again at once. Openings and sighted targets receive a three-heading peek, widened to five headings when the agent has stepped through onto an opening's threshold. New rooms or object kinds seen during a visit are acted upon only when it ends; only a target sighting preempts a visit, and a sighted-target node is therefore resolved by target closing whenever the peek produces a detection that passes the takeover gates.

\begin{algorithm}[htbp]
\DontPrintSemicolon
\caption{\textsc{ExploreLocal}$(v)$}
\label{alg:explorelocal}
\Input{node $v \in \N$ in force, reached by transit; $k$ actions already spent on the visit}
\Output{at most one action executed (as in \textsc{Navigate}); the visit's verdict when it ends}
\BlankLine
\uIf(\Comment*[f]{room: walk its frontiers}){$v$ is a room}{
  $Q \gets$ live frontiers of $v$ (\cref{alg:candidatenodes}), best first by view gain $U_f$, re-read every action\;
  \uIf{$k \ge 60$}{end the visit: \emph{budget spent}\;}
  \uElseIf(\Comment*[f]{e.g.\ a full path glance}){a full rotation inside $v$ was completed since the visit began}{end the visit: \emph{finished}\;}
  \uElseIf{A$^*$ reaches a goal of $Q$}{\Navigate{$\tau_v$} \Comment*[r]{$\tau_v$: committed route to the best goal reached, kept while it stays passable, near $v$ and informative; other rooms count as unknown}}
  \lElse{end the visit: \emph{finished}}
}
\Else(\Comment*[f]{opening, or sighted target from its standoff}){
  \uIf{$k \ge 30$}{end the visit: \emph{budget spent}\;}
  \uElseIf(\Comment*[f]{resolved on the way}){$v$ is an opening absent from the frontier map for 3 actions running while $> 1.5$\,m away}{end the visit: \emph{finished}\;}
  \uElseIf{$x_t$ is not within 0.35\,m of $e_v$}{\Navigate{$\tau_v$} \Comment*[r]{$\tau_v$: committed route to $e_v$; at most $\max(20, \lceil 1.6\,n \rceil + 6)$ actions, plus at most 8 to step through}}
  \uElseIf{one of the headings $\psi_v$, $\psi_v \pm 30^\circ$ (after a step through, the direction walked and $\pm 30^\circ$, $\pm 60^\circ$ about it) is not yet looked at}{turn toward the next, starting from the side nearer the current yaw\;}
  \Else{
    mark an opening peeked (with the cone it swept) or a sighted target inspected; end the visit: \emph{finished}\;
  }
}
\Comment{a box that passes the takeover gates suspends the visit and hands control to \TargetClosing{} (\cref{alg:targetclosing}); after a release the visit resumes, except that a sighted-target node whose candidate is released is retired}
\end{algorithm}

Here $U_f = (n_f + s_f)/((g_f + 0.25)(1 + |\Delta\psi_f|/\pi))$ is the view gain of a frontier cluster of $s_f$ cells at geodesic distance $g_f$ (in metres) and heading error $\Delta\psi_f$, where $n_f$ counts the unknown cells that a full circle of camera rays from it (0.6--5.0\,m, blocked by walls and by settled unknown) would reach. $\psi_v$ is the direction from the opening's free side toward the unknown (for a sighted target, the bearing of the object from the standoff point), and $n = \lceil d/0.25 \rceil$ is the number of forward steps on the geodesic route of length $d$ to the threshold; if the walk overruns its bound, the agent looks from where it stands when within 1.5\,m of the threshold and retires the opening otherwise. When the transit ends more than 0.35\,m short of an opening's passable threshold, the agent steps through onto it (at most 8 actions, once per peek, never after such an overrun); a threshold that is not yet passable is looked at from where the transit ended and stepped through once that look has made it passable.

\paragraph{Target closing.}
\Cref{alg:targetclosing} is entered on the action in which a box of the target class passes the takeover gates of \cref{sec:execution}: confidence $\ge 0.50$, scaled by 0.70 for a box in a suspect context; no sliver cut by the image border; a depth-supported 3-D point at a plausible height; and, beyond 2.0\,m, outside the remembered release disks and the 15-action cooldown, which a confidence of at least 0.80 can override except for a 2.0\,m disk. The routine then owns every action until it stops or releases the candidate, with room segmentation, room labelling, the oracle and RPT$^*$ suspended; each phase is bounded (12 actions to verify, 24 to inspect, 160 for the whole closing, the first and the last not counting up to 24 steps toward a glimpse), and the body of the loop is evaluated once per action. The line marked ``rules (i)--(v)'' stands for the five STOP rules on remembered evidence listed in \cref{sec:execution}, which also gives the rationale of every gate, the release radii and the pitch control with and without camera tilt.

\begin{algorithm}[htbp]
\DontPrintSemicolon
\caption{\textsc{TargetClosing}$(\cstar)$}
\label{alg:targetclosing}
\Input{target category $\cstar$; the box $b_0$ that passed the takeover gates; scene graph $\G$}
\Output{\Stop at the target, a release that returns control to the search, or \Fail at the closing bound}
\State{anchor $X_0$; fused estimate $\hat{X}$; phase $\phi$; consecutive sightings $n$; standoff goal $g$ and its committed route $\tau_g$}
\BlankLine
$X_0 \gets \hat{X} \gets$ 3-D point of $b_0$;\quad $\phi \gets \textsc{verify}$;\quad $n \gets 1$\;
\While(\Comment*[f]{closing bound}){fewer than 160 closing actions are spent}{
  $b \gets$ fresh box of class $\cstar$, confidence $\ge 0.30$, inside the association cylinder of $X_0$, else $\varnothing$\;
  \lIf(\Comment*[f]{precision-weighted}){$b \neq \varnothing$}{fuse its 3-D point into $\hat{X}$;\quad $d \gets$ its measured range}
  \uIf(\Comment*[f]{is the sighting real?}){$\phi = \textsc{verify}$}{
    $\gamma \gets$ a \emph{glimpse}: a box of class $\cstar$ with confidence $\ge 0.30$, not cut by the border, without usable depth, within $30^\circ$ of the bearing of $X_0$; else $\varnothing$\;
    \lIf{$b \neq \varnothing$}{$n \gets n + 1$ if no frame without a sighting or a glimpse came since the previous sighting, else $n \gets 1$}
    \uIf(\Comment*[f]{once per object; not counted}){$X_0$ was first seen within 1.30\,m at confidence $< 0.92$, a box of class $\cstar$ at $\ge 0.50$ was in fewer than 2 of the 6 frames before, \KwAnd no back-off was made}{
      back off up to 3 steps to a verified clear spot with no mapped obstacle on its line to $X_0$, if one exists, and face $X_0$;\quad $n \gets 0$ \Comment*[r]{$\le 20$ actions}
    }
    \uElseIf(\Comment*[f]{lock; suspect box: 4 frames, 2 viewpoints}){$n \ge 2$}{$\phi \gets \textsc{approach}$\;}
    \uElseIf{12 verification actions are spent \KwOr 6 sweep actions in a row brought neither a box nor a glimpse}{\KwRet{\Release{unverified, 1.0\,m}}\;}
    \uElseIf(\Comment*[f]{steps not counted, $\le 24$}){$b = \varnothing$ \KwAnd $\gamma \neq \varnothing$}{face $\gamma$; once within $30^\circ$ of it, step forward if $\lVert x_t - \hat{X} \rVert > 1.25$\,m\;}
    \uElseIf{$b = \varnothing$}{sweep the bearing of $X_0$ by $0, +30^\circ, 0, -30^\circ$\;}
    \uElseIf{$b$ is more than $30^\circ$ off the heading}{turn to face $b$\;}
    \uElseIf(\Comment*[f]{without tilt: only if $b$ stays in frame}){$d > 1.25$\,m}{step forward\;}
    \lElse{take the view (tilt or turn) that keeps $b$ whole}
  }
  \uElseIf(\Comment*[f]{reach a standoff point}){$\phi = \textsc{approach}$}{
    \uIf{($b \neq \varnothing$ \KwAnd $d \le 1.05$\,m) \KwOr (($\lVert x_t - \hat{X} \rVert \le 1.0$\,m \KwOr $g$ is reached) \KwAnd no fresh $d > 1.05$\,m)}{$\phi \gets \textsc{inspect}$\;}
    \Else{
      \If(\Comment*[f]{goal kept within 1.0\,m}){no $g$ \KwOr $\lVert g - \hat{X} \rVert > 1.0$\,m}{
        $g \gets$ first point on the 0.65\,m ring about $\hat{X}$ at bearings $0, \pm30^\circ, \pm60^\circ, \pm90^\circ, 180^\circ$ from the agent's side, projected onto the navigable surface, with an A$^*$ path $\tau_g$ ending $\le 0.75$\,m from $\hat{X}$; if no bearing has one on the known floor, the bearings again with unknown cells passable at 8 times the cost of free floor\;
      }
      \uIf{no $g$ exists for $2 + 6$ pathless actions}{\KwRet{\Release{boxed in; remember the agent's spot}}\;}
      \uElseIf{the last 20 actions stayed within 0.20\,m}{\KwRet{\Release{stalled}}\;}
      \lElse(\Comment*[f]{committed route; A$^*$ owns the heading}){\Navigate{$\tau_g$}}
    }
  }
  \Else(\Comment*[f]{at most 24 actions, at most 3 of them steps}){
    \uIf{$b \neq \varnothing$ \KwAnd $d \le 1.05$\,m \KwAnd ($b$ centred \KwOr extended approach)}{
      \uIf(\Comment*[f]{weak box}){$b$ has confidence $< 0.50$ \KwAnd fewer than 2 holds were made}{hold a view that keeps $b$ whole\;}
      \lElse{\KwRet{\Stop}}
    }
    \lElseIf(\Comment*[f]{resumption}){$b \neq \varnothing$ \KwAnd $d > 1.05$\,m \KwAnd ($b$ centred \KwOr no sighting yet in this inspection)}{drop $g$ if it was reached;\quad $\phi \gets \textsc{approach}$}
    \lElseIf(\Comment*[f]{one step out}){$b \neq \varnothing$ \KwAnd $1.05 < d \le 1.30$\,m \KwAnd $b$ lies within $30^\circ$ of the heading \KwAnd fewer than 3 such steps were made}{step forward}
    \lElseIf{$b \neq \varnothing$}{turn to centre $b$}
    \lElseIf(\Comment*[f]{rules (i)--(v), \cref{sec:execution}}){a rule on remembered evidence holds}{\KwRet{\Stop}}
    \lElseIf{24 actions spent \KwOr all views blank}{\KwRet{\Release{inspection saw nothing, 2.0\,m}}}
    \lElse{turn to the next view: yaw $0, +30^\circ, -30^\circ$ about the stored bearing}
  }
}
\KwRet{\Fail} \Comment*[r]{closing bound reached: the episode ends unsuccessfully}
\BlankLine
\Comment{\Stop: decided on the current frame, issued after the agent steps toward $\hat{X}$ once, and once more if $\hat{X}$ is still $> 0.90$\,m away, consulting no further frame; a step is taken only if afterwards $\hat{X}$ stays $\ge 0.45$\,m and the near edge measured on the deciding frame one body radius ahead, and while $\hat{X}$ lies within $90^\circ$ of the heading; a blocked step or 5 actions end the steps}
\Comment{\textsc{Release}: store $X_0$ with its radius (if boxed in, the agent's spot), start a 15-action cooldown, resume the interrupted search task}
\Comment{tilt benchmarks: the box servo sets the pitch within 1.5\,m; at most one centring tilt per inspection precedes the STOP}
\end{algorithm}

\section{Hyper-parameters}\label{app:hyperparams}

\Cref{tab:hyperparams} lists the parameters of \ours with their default values. Lengths are in metres and times in actions; with the 0.25\,m forward step, the cost model charges 4 actions per metre of travel. All values are fixed for every episode and scene. They were set during development on Gibson training buildings, except that faults exposed by the first three validation episodes of each scene were corrected, the depth-hole filling was designed on frames of one validation episode, two component choices (the detector checkpoint size and the default local explorer) were compared on validation episodes, and the target-closing routine was revised after inspecting failures of a validation run; later, rules and values throughout the search, perception, mapping and target closing, and the camera tilt on Gibson, were changed after inspecting 15 further validation episodes with the toilet goal and kept or reverted by their validation scores (\cref{sec:experiments}). \TODO{confirm that HM3D, MP3D and RoboTHOR use the same values, including the target-closing routine and whether stair traversal is enabled}. The room classifier chooses from the fixed set of 15 labels given in \cref{sec:perception} (nine residential types, five institutional ones and \emph{unknown}). When an instance violates the triangle inequality that RPT$^*$ requires, or the solver fails, the next node is drawn at random in proportion to $p_i$ instead (\cref{sec:planning}).

\begin{table}[p]
\centering
\small
\caption{Hyper-parameters of \ours (part 1 of 4): perception, mapping, doorways and room segmentation. Lengths in metres, times in actions; the scene graph is rebuilt after every action without an LLM call.}
\label{tab:hyperparams}
\begin{tabular}{@{}>{\raggedright\arraybackslash}p{1.9cm}>{\raggedright\arraybackslash}p{3.1cm}>{\raggedright\arraybackslash}p{3.8cm}>{\raggedright\arraybackslash}p{5.6cm}@{}}
\toprule
Group & Parameter & Value & Role \\
\midrule
Perception & detector & YOLO-World X-v2, 640\,px & open vocabulary~\citep{cheng2024yolo}; CLIP ViT-B/32 text encoder~\citep{radford2021learning} \\
 & NMS IoU; max.\ boxes & 0.5; 100 & one detector call per frame \\
 & vocabulary & 30 prompts & goal classes and aliases, context objects, door, stair and distractor prompts \\
 & score thresholds & 0.05 / 0.30 / 0.50 & emitted / enters the object map / starts target closing \\
 & duplicate boxes & IoU $\ge 0.5$ (same class), $\ge 0.9$ (any) & one object read under two prompt names \\
 & box depth & 30th percentile, central half & valid depth 0.5--5.0\,m on the Habitat benchmarks \TODO{state the depth range the box-depth and occupancy filters use on RoboTHOR} \\
 & depth support & component $\ge 35\%$ & object point = per-axis median of it \\
 & footprint radius & 0.15--1.5\,m & from box width and depth, clamped \\
 & object height & $-0.30$ to $+2.2$\,m & above the floor; floor support required \\
 & same object & $\le 0.35$\,m or footprint IoU $\ge 0.25$ & same class only; heights within 0.35\,m \\
 & confirmation & 2 distinct frames & before an object reaches rooms or LLMs \\
 & depth hole filling & holes $\le 50\%$ of the frame & enclosed holes only; a deviation (\cref{tab:privileges}) \\
 & optional modules & off & Grounding DINO + BLIP-2~\citep{liu2024grounding,li2023blip}; SAM, MobileSAM~\citep{kirillov2023segment,zhang2023faster} \\
\midrule
Mapping & grid & 0.10\,m cells; 80\,m $\times$ 80\,m & never inflated \\
 & depth sampling & every 8th pixel & 4{,}800 points per frame \\
 & obstacle; floor band & 0.15--0.93\,m; $\pm 0.10$\,m & cells on rays to obstacle hits are freed \\
 & log-odds steps & hit $+0.85$ (scaled); free $-0.40$ & clamped to $\pm 4$ \\
 & cell states & free $P_{\mathrm{occ}} \le 0.45$; occupied $P_{\mathrm{occ}} \ge 0.65$ & otherwise unknown \\
 & settled unknown & 2 looks within 2.5\,m; pockets $\le 3$\,m$^2$ narrower than 0.6\,m or stood on & occupied for frontier extraction only \\
 & blocked step & cell 0.25\,m ahead set to $+4$; 2nd bump within 0.5\,m & bump evidence; a repeat bump leaves a body-wide bar occupied for the episode \\
 & seen staircases & masked as obstacles & ground-truth geometry once visible (\cref{tab:privileges}) \\
\midrule
Doorways & door prompts & 4; confidence $\ge 0.05$ & depth read on the jambs and lintel \\
 & door gates & width 0.35--1.5\,m; top $\ge 1.4$\,m; bottom $\le 0.5$\,m & apparent size above the floor \\
 & door confirmation & 3 sightings; one $\ge 0.2$\,m or $\ge 20^\circ$ away & association radius 0.6\,m \\
\midrule
Segmen\-tation & seeds & $\ge 1.0$\,m apart; clearance $\ge 0.3$\,m & clearance maxima~\citep{vincent1991watersheds} \\
 & basin merge & saddle $< 0.2$\,m & never across a cut \\
 & minimum room & 0.5\,m$^2$ & smaller regions dropped \\
 & door cut & snap within 0.9\,m, else 0.75\,m disc & at the narrowest passage \\
 & walked threshold & dip $\le 0.55$\,m; rise $\ge 0.30$\,m & within 1.5\,m along the trail; snap 0.5\,m \\
 & room identity & anchor; IoU 0.15; containment 0.6 & memory of 48 rooms, 1000 updates \\
 & object-to-room reach & footprint radius $+\,0.6$\,m & furniture stands on occupied cells \\
\bottomrule
\end{tabular}
\end{table}

\begin{table}[p]
\centering
\small
\caption*{\Cref{tab:hyperparams} (part 2 of 4): room labels, candidate nodes, costs and oracle.}
\begin{tabular}{@{}>{\raggedright\arraybackslash}p{1.9cm}>{\raggedright\arraybackslash}p{3.1cm}>{\raggedright\arraybackslash}p{3.8cm}>{\raggedright\arraybackslash}p{5.6cm}@{}}
\toprule
Group & Parameter & Value & Role \\
\midrule
Room labels & model & Qwen2.5-3B-Instruct & Ollama; temperature 0, seed 17, JSON~\citep{yang2024qwen25} \\
 & output; timeout & $\le 768$ tokens; 120\,s & one call per room re-labelled \\
 & label set & 15 labels & listed in \cref{sec:perception} \\
 & generic objects & 7 classes & cabinet, potted plant, book, vase, clock, cup, bottle \\
 & signature objects & bed; toilet, shower, bathtub; oven, stove, refrigerator & bedroom; bathroom; kitchen (no stove prompt on Gibson) \\
 & firm label & confidence $\ge 0.65$ & plus distinctive evidence (\cref{alg:classifyrooms}) \\
 & refresh & 50 actions & besides decision points and new kinds \\
\midrule
Candidate nodes & vantage clearance & $\ge 0.4$\,m & room entry after the best frontier and the doorstep (first on a revisit); else the centroid \\
 & revisit standoff & 1.5\,m & second-pass vantage from earlier scan points \\
 & finished by sight & $\ge 50\%$ of cells within 5\,m & seen from a completed scan, with no live frontier left \\
 & live frontier & reveals $\ge 1.0$\,m$^2$ & a room stood in $\ge 0.5$\,m from its edge with none left is finished \\
 & doorless fragment; narrow room & $< 3$\,m$^2$; clearance $\le 0.9$\,m & finished once no live frontier is left (the narrow room once walked through) \\
 & regrowth & $1.5\times$ & a finished room is judged again \\
 & opening & width $\ge 0.5$\,m; merge 1.5\,m & reachable frontier exits \\
 & opening exclusions & 1.2 / 0.30 / 1.5 / 1.0\,m & blind radius / shadow margin / stairs or drop on the unknown side / near peeks \\
 & doorway vs.\ gap & door within 1.0\,m & label shown to the oracle \\
 & sighted target & $p = 0.85$; ring 1.5\,m (else 2.25\,m) & no LLM call \\
 & entry snap & 2.0\,m & otherwise the node is dropped \\
 & probability floor $p_{\min}$ & 0.05 & less likely nodes are not offered \\
 & deferral & 3 failures $\Rightarrow$ 180 actions & unreachable, timed out or blocked \\
 & cooldown & 120 actions; 30 & after a visit that did more than spend its budget; after a failed transit \\
 & node cap & 12 & nearest node plus the 11 most probable \\
\midrule
Costs & travel speed $u$ & 0.25\,m per action & turns are not charged \\
 & distance $d_{ij}$ & Dijkstra, 8-connected, 0.10\,m & at 0.30\,m clearance if possible, else 0.18\,m~\citep{dijkstra1959note} \\
 & service time $\sigma_j$ & room $\min(60, 4 + 4F)$ (12 if $F$ is unknown); stairs 12; others 4 & $F$ live frontiers of the room; folded into entering arcs (\cref{eq:arccost}) \\
 & stair leaf $\lambda_i$ & flight length $+\,8$\,m & multi-storey mode only \\
\midrule
Oracle & model & Qwen2.5-14B-Instruct, 4-bit & Ollama; temperature 0, seed 17, JSON~\citep{yang2024qwen25} \\
 & context; output & 8192; $\le 1536$ tokens & timeout 900\,s \\
 & objects per room & $\le 8$ classes & in the node description \\
 & budget in the prompt & steps of 25 actions & coarse, so unchanged maps repeat the prompt \\
 & score mapping & 0--100 $\to [0, 1 - 10^{-6}]$ & independent, not normalised \\
 & omitted node & 0.01 & a listed node the reply does not score \\
 & unscored opening & last score, else 0.10 & an opening not listed when the prompt was built \\
 & value factors & gap 0.25--1 (full at 4\,m$^2$); entered room 0.5, or 0.2 & fixed rules after the reply; 0.2 after $\ge 20$ actions in the room with $\le 1$ live frontier left \\
 & schema repair & 1 retry & then back-off \\
 & back-off & 25 to 200 actions & doubling; the fallback explores meanwhile \\
 & re-query & 6 actions; 25 actions & new opening; budget step in the second pass \\
\bottomrule
\end{tabular}
\end{table}

\begin{table}[p]
\centering
\small
\caption*{\Cref{tab:hyperparams} (part 3 of 4): RPT$^*$, execution and local visits.}
\begin{tabular}{@{}>{\raggedright\arraybackslash}p{1.9cm}>{\raggedright\arraybackslash}p{3.1cm}>{\raggedright\arraybackslash}p{3.8cm}>{\raggedright\arraybackslash}p{5.6cm}@{}}
\toprule
Group & Parameter & Value & Role \\
\midrule
RPT$^*$ & search & exact best-first & optimal for the instance~\citep{lyu2026rpt} \\
 & time budget & 2.0\,s per solve & then the best order found, uncertified \\
 & precondition & triangle inequality & relative tolerance $10^{-9}$ \\
 & fallback & probability-weighted draw & of the next node if the solve fails; seeded per episode \\
 & horizon & first node only & re-solved at the next decision point \\
\midrule
Execution & warm-up & 10 turns of $30^\circ$ & every bearing seen by the $79^\circ$ camera; restarted if moved $> 0.5$\,m; the spawn spot then counts as scanned \\
 & A$^*$ & 8-connected, octile & unknown cells blocked~\citep{hart1968formal}, except in a last-resort standoff route of target closing \\
 & A$^*$ clearance & 0.30\,m, down to 0.18\,m & relaxed in 0.05\,m steps if no path; 0.18\,m if the clearance route is $> 1.5\times$ longer \\
 & clearance cost & cell cost 1 to 4 & over a 0.8\,m band outside the 0.30\,m clearance \\
 & route commitment & 0.5\,m; 0.5\,m; 0.10\,m per 30 & goal shift; off-route; progress \\
 & follower & lookahead 0.5\,m; $15^\circ$ & heading dead band; monotone progress \\
 & arrival & 0.25\,m & reach floor 0.129\,m never binds \\
 & blocked step & advance $< 0.05$\,m & along the heading; then aim 0.25\,m ahead \\
 & idle action & left turn & never a STOP \\
 & transit release & 5 / 120 / 6 actions & no route / timeout / blocked \\
 & room visit & 60 & whole visit, from arrival; walks the live frontiers by view gain \\
 & opening peek & $0, \pm 30^\circ$; 0.35\,m; 30 & headings; arrival; whole visit \\
 & step through & 8 actions; $0, \pm 30^\circ, \pm 60^\circ$ & onto a passable threshold; headings about the direction walked \\
 & peek resolved & absent 3 actions, $> 1.5$\,m away & the opening's peek ends without a look \\
 & threshold walk & $\max(20, \lceil 1.6\,n \rceil + 6)$ & $n$ forward steps on the route \\
 & path glance & $\ge 3$\,m$^2$; $\ge 0.5$\,m$^2$ per action & unseen floor a side look must reveal; at most one per route, within 6\,m ahead \\
 & frontier utility & $(n_f + s_f) / ((g_f + 0.25)(1 + |\Delta\psi_f|/\pi))$ & room visit and exploration fallback: $n_f$ unknown cells a full circle of camera rays (0.6--5.0\,m) from it reaches, $s_f$ cluster size in cells, $g_f$ geodesic distance in m, $\Delta\psi_f$ heading error; a new fallback goal must score $2\times$ the current one \\
\bottomrule
\end{tabular}
\end{table}

\begin{table}[p]
\centering
\small
\caption*{\Cref{tab:hyperparams} (part 4 of 4): target closing (\cref{alg:targetclosing}).}
\begin{tabular}{@{}>{\raggedright\arraybackslash}p{1.9cm}>{\raggedright\arraybackslash}p{3.1cm}>{\raggedright\arraybackslash}p{3.8cm}>{\raggedright\arraybackslash}p{5.6cm}@{}}
\toprule
Group & Parameter & Value & Role \\
\midrule
Target closing & takeover & $\ge 0.50$; suspect $\times 0.70$ & no border sliver (8\,px margin); valid 3-D point \\
 & takeover vetoes & release disks; 15 actions & beyond 2.0\,m; override at $\ge 0.80$ \\
 & re-sighting & confidence $\ge 0.30$ & inside the association cylinder \\
 & association & 0.50\,m $+$ 0.15\,m per m & beyond 2\,m; depth $+\max(0.60\,\text{m}, 2\rho)$ \\
 & lock & 2 consecutive frames & suspect: 4 frames, 0.30\,m baseline; glimpse frames keep the run consecutive \\
 & glimpse & no usable depth; within $30^\circ$ of the bearing & faced and stepped toward beyond 1.25\,m; $\le 24$ steps, not counted \\
 & verification budget & 12 actions; 6 sweep actions & else released as unverified; sweep: in a row with neither box nor glimpse \\
 & verification step & range $> 1.25$\,m & without tilt: only if the box stays in frame \\
 & close first sighting & within 1.30\,m & back off $\le 3$ steps, $\le 20$ actions, once, to a spot with a clear line to the target; skipped at confidence $\ge 0.92$ or after boxes $\ge 0.50$ in 2 or more of the 6 previous frames \\
 & target estimate & s.d.\ $0.03 + 0.03\,d$\,m & precision-weighted fusion \\
 & standoff ring & 0.65\,m; 8 bearings & on the navigable surface (\cref{tab:privileges}) \\
 & standoff path & ends $\le 0.75$\,m from target & accepted standoff; else through unknown cells at 8 times the free-floor cost \\
 & goal refinement & 0.15\,m; kept within 1.0\,m & when the estimate moves \\
 & terminal range & $1.00 + 0.05$\,m & to the near edge (10th percentile) \\
 & stop confidence & 0.50; $\le 2$ extra frames & for a weaker box \\
 & centred & spans the centre, or $\le 15^\circ$ & yaw error of the box; $< 30^\circ$ when the turn would undo the previous centring turn \\
 & inspection & 24 actions; yaw $0, \pm 30^\circ$ & three pitches each with tilt \\
 & inspection step & range 1.05--1.30\,m; within $30^\circ$; $\le 3$ & one forward step instead of turning \\
 & STOP overshoot & 1 step; a 2nd beyond 0.90\,m; $\le 5$ actions & toward the estimate, leaving it $\ge 0.45$\,m ahead; no frame consulted \\
 & pitch servo & 1.50\,m; regions 0.25 & tilt only; dead band $\ge$ one tilt \\
 & below-frame rule & $< 15\%$ of rows; 6 actions & no tilt; estimate within 1.05\,m \\
 & extended approach & 3.0\,m walked; 1.0\,m closed & then 3 frames without a box \\
 & stall & 20 actions within 0.20\,m & release \\
 & no safe path & $2 + 6$ actions & release; boxed-in spot 1.0\,m \\
 & release radius & 1.0 / 2.0\,m & unverified / disproved; none if a candidate never seen again after a back-off had a frame at $\ge 0.92$ \\
 & closing bound & 160 actions & then the episode ends; steps toward a glimpse not counted \\
\bottomrule
\end{tabular}
\end{table}

\clearpage

% ============================================================================
% 12_appendix_protocols
% ============================================================================
% ============================================================================
%  12_appendix_protocols.tex -- Appendix: LLM prompts, benchmark protocols,
%  additional results
% ============================================================================

\section{LLM Prompts}\label{app:prompts}

\ours queries two instruction-tuned language models of the Qwen2.5 family~\citep{yang2024qwen25}. A small model, Qwen2.5-3B-Instruct, names each room from the objects confirmed in it (\cref{sec:perception}); a larger model, Qwen2.5-14B-Instruct (4-bit quantised, about 9\,GB), estimates the per-node search-success probabilities that RPT$^*$ consumes (\cref{sec:planning}). Both run locally behind an Ollama server, are decoded greedily (temperature~0) with a fixed seed (17) in the server's JSON output mode, and are therefore reproducible up to the numerical nondeterminism of the inference backend. The identity (digest) of each model is recorded at the start of a run and checked before every call, so that a model cannot change in the middle of an evaluation. The room classifier may produce at most 768 output tokens; the oracle runs with an 8192-token context and at most 1536 output tokens.

The prompts below are reproduced verbatim from the implementation \TODO{verify the prompt text against the final code used for the reported runs}. Only the line breaks were changed for typesetting: a line indented by four spaces continues the line above it, and the line break together with those four spaces stands for one space of the original (a fifth leading space, where present, belongs to the text); the model receives the joined line. Words in braces such as \texttt{\{target\}} are fields filled at query time and are described in the text; every other character, including the braces of the JSON schemas, is sent literally.

\subsection{Room classification}\label{app:prompts-room}

The room classifier is queried for a room at a decision point when the room's object evidence (the multiset of confirmed object classes in it) has changed, or when the room has not been re-asked for 50 actions; it is also queried once in the middle of a visit when a new kind of object appears in the room being searched or traversed (\cref{sec:perception}). When the evidence has changed, a verdict already obtained for an identical evidence signature (or, in the middle of a visit, for the same set of object kinds) is reused without a call; only the periodic refresh deliberately asks again about evidence the model has already judged. Rooms that contain no confirmed object, or only generic ones (cabinet, potted plant, book, vase, clock, cup, bottle), are labelled \texttt{unknown} without querying the model. The number of calls for a room is therefore bounded by the number of distinct evidence signatures it passes through, plus at most one refresh every 50 actions.

\paragraph{Input.} The system prompt is fixed except for the label set, which is inserted as a comma-separated list. The user prompt carries a single field, \texttt{\{obj\_list\}}: the confirmed object landmarks inside the room, one line per class in alphabetical order, written as \texttt{- <class> x<count>} (for example \texttt{- bed x1}). Detector aliases are merged before counting, so a couch is reported as \texttt{sofa} and a tv as \texttt{television}. The template also provides the placeholder \texttt{(no objects observed yet)} for an empty list, but it is never sent, because empty rooms do not reach the model.

\paragraph{Reply.} The model returns a JSON object with three keys: \texttt{label} (one entry of the label set), \texttt{confidence} (a number in $[0,1]$) and \texttt{reasoning} (one short sentence). The label is lower-cased and spaces and hyphens are mapped to underscores, so that \emph{Living room} reads as \texttt{living\_room}; any label outside the set becomes \texttt{unknown}. The confidence is clipped to $[0,1]$, a non-finite value is read as~0, and the reasoning is truncated to 200 characters. A label is \emph{firm} only if it is not \texttt{unknown}, its confidence is at least 0.65, and the evidence is either at least two object kinds, at least one of them non-generic, or one signature object that agrees with the label (a bed for a bedroom; a toilet, shower or bathtub for a bathroom; an oven, stove or refrigerator for a kitchen). Otherwise the label is \emph{tentative}, and the oracle sees it with a trailing question mark (\cref{app:prompts-oracle}). An exception from the model server is not cached; the agent keeps the labels it already holds and backs off as described for the oracle below.

\paragraph{Label vocabulary.} The label set offered to the model has 15 entries: \texttt{kitchen}, \texttt{bedroom}, \texttt{bathroom}, \texttt{living\_room}, \texttt{dining\_room}, \texttt{office}, \texttt{hallway}, \texttt{storage\_closet}, \texttt{laundry\_room}, \texttt{lobby}, \texttt{waiting\_area}, \texttt{patient\_room}, \texttt{exam\_room}, \texttt{reception} and \texttt{unknown}. Five of them (\texttt{lobby}, \texttt{waiting\_area}, \texttt{patient\_room}, \texttt{exam\_room} and \texttt{reception}) are non-residential labels carried over from an earlier deployment of the classifier in a hospital setting, and they are offered on the residential benchmarks as well. \TODO{prune non-residential labels} \TODO{or, if they are kept for the reported runs, report how often the model chose one of them on each benchmark}

\noindent System prompt (label set filled in):
\begin{small}
\begin{verbatim}
You are a scene-understanding assistant that
classifies indoor rooms from the objects observed inside them.

You will be given a list of objects observed in one room. Based ONLY on
those objects and common sense about where they occur, output one room label.
Observations are incomplete and may contain false detections. Prefer distinctive
room-function evidence over generic furniture. Use unknown when evidence is
insufficient or contradictory. Do not invent unobserved objects. Reassess the
current evidence rather than assuming an earlier room label was correct.

Choose the label from this set (do not invent new ones):
kitchen, bedroom, bathroom, living_room, dining_room, office, hallway, storage_closet,
    laundry_room, lobby, waiting_area, patient_room, exam_room, reception, unknown

Reply with a JSON object of the form:
{"label": "<one label from the set>",
  "confidence": <float between 0 and 1>,
  "reasoning": "<one short sentence>"}
\end{verbatim}
\end{small}

\noindent User prompt:
\begin{small}
\begin{verbatim}
Room observed objects:
{obj_list}

Classify this room.
\end{verbatim}
\end{small}

\subsection{Search-node probability oracle}\label{app:prompts-oracle}

The oracle is queried at every decision point with the goal category and the full candidate set, except the sighted-but-unverified target landmarks, whose probability is fixed at 0.85 without a model call (\cref{sec:planning}). One call answers the whole question: for each node, the probability that going there next and looking finds the target. The system prompt states the rules of judgement in prose and gives one worked example; the user prompt states the facts of the current episode.

\paragraph{Input.} The user prompt has four fields. \texttt{\{target\}} is the goal category name. \texttt{\{house\}} is the house summary $H$, one line of facts in the prompt's own vocabulary: the room types found so far, each with a question mark if its label is tentative and with its status (scanned, entered or never entered); the number of rooms still unidentified and how many of them were scanned, or ``every known room identified''; the number of openings not yet looked into that are offered as nodes; whether any unexplored frontier is still reachable; and the budget used, rounded down to a multiple of 25 actions (for example ``actions used about 150 of 500''). \texttt{\{n\}} is the number of nodes, and \texttt{\{nodes\}} holds one line $\eta_i$ per node.
A room line gives an integer identifier and the fields \texttt{type=} (the room label, \texttt{unknown} allowed, with a trailing question mark if tentative), \texttt{size=} (the floor area rounded to square metres), \texttt{frontier=} (the number of reachable frontier clusters inside the room), \texttt{status=} (one of \texttt{never\_entered}, \texttt{entered} or \texttt{scanned(<how>, <ago>)}, where \emph{ago} is the time since the scan and \emph{how} is one of ``full rotation'', ``seen from another room's scan'', ``looked into, no frontier left'', ``fragment'' or ``walked, every part of it seen, no frontier left'', mirroring the room-finished rules of \cref{sec:planning}), \texttt{searched=} (the time spent inside the room), \texttt{here=yes} for the room the agent stands in, and \texttt{seen:} (up to eight confirmed object classes in alphabetical order, or ``nothing yet'').
An opening line gives its identifier, whether it is a doorway or a gap, the room it leaves together with that room's label, and the object classes glimpsed through it, that is, confirmed landmarks other than doors that lie 0.3 to 3.0\,m beyond the threshold and within $35^\circ$ of the direction of the opening.
Times are expressed at a nominal one second per action and coarsened to 10\,s steps below one minute and to whole minutes above. Together with the rounded budget, this makes the prompt byte-identical whenever the map has not changed in a way the model can see, and the previous reply is then reused without a call.
\emph{No distance or travel cost appears anywhere in the prompt}: travel is charged by the planner through the cost matrix (\cref{eq:arccost}), and the system prompt explicitly forbids using distance, travel time or the \texttt{here=yes} flag as a reason.

\paragraph{Reply.} The model must return one JSON object whose keys come in a fixed order, so that the reasoning is written before any number~\citep{wei2022chain}: \texttt{home} (the room types where the target normally lives, at most 12 words), \texttt{house} (what the rooms found imply about the rooms not yet identified, at most 25 words), \texttt{stage} (the budget and search stage, at most 20 words), \texttt{pass} (\texttt{first} or \texttt{second}), and \texttt{nodes}, a list of objects \texttt{\{"id", "why", "p"\}} in which a reason of at most 12 words precedes an integer score \texttt{p} in $[0,100]$. The \texttt{pass} verdict is recorded but does not affect the search: finished rooms do not re-enter the candidate set (\cref{sec:planning}).
The parser also accepts the score under the keys \texttt{probability} (multiplied by 100 if at most 1), \texttt{percent} or \texttt{score}; it clamps scores to $[0,100]$, drops identifiers that were not offered, assigns 1\% to an offered node that the reply omits, and rejects a reply that scores no offered node, whereas a reply in which every score is zero is valid. The probability of node $i$ is $p_i = \min(1-10^{-6},\, r_i/100)$ for its score $r_i$; the scores are not normalised (\cref{sec:problem}). Before the solver uses them, fixed rules scale the probabilities of gaps and of rooms the agent has entered and left (\cref{sec:planning-oracle}).

\paragraph{Failures.} A transient transport error is retried once; a timeout is not. An unusable reply, or a call that still fails, triggers one schema-repair request, in which the text shown last below is appended, after a line break, to the same user prompt. If the repaired reply is also unusable, no uniform prior is substituted: the agent runs the exploration fallback for 25 actions before asking again, doubling the wait after each consecutive failure up to 200 actions, and a successful call resets the wait.

\paragraph{Multi-storey supplement.} When a staircase node is offered or another storey is known, the supplement shown below is appended to the system prompt, and the user prompt gains two storey lines: \texttt{\{storey\}} summarises the current storey (the room types found, the time searched there, the number of rooms with frontier left, the number of storeys known to the building and the current storey's rank among them), and \texttt{\{others\}} holds one such summary per other known storey. Stair nodes are described by their direction, whether the destination storey was visited (with its summary) and whether the agent has just arrived by them, and the reply gains the key \texttt{home\_here} (\texttt{found}, \texttt{missing} or \texttt{elsewhere}). Multi-storey search is disabled on the evaluated single-storey Gibson split, so the supplement and the storey lines are never sent there. \TODO{state whether the multi-storey mode, and hence this supplement, is enabled on HM3D and MP3D}

\noindent System prompt:
\begin{small}
\begin{verbatim}
You are the reasoning module of a robot searching ONE home for an instance of ONE target
    object category, within a fixed budget of actions. The robot has partly mapped the
    home into NODES. For EACH node you estimate the probability, in percent, that going
    there NEXT finds the target: the target is there and the robot would see it by going
    and looking. You reason like a person who knows how homes are laid out and where
    things are kept.

NODE LINES.
ROOM: a region the robot segmented. type= was inferred from the objects seen in it
    (unknown = NOT IDENTIFIED YET, never "empty"; a type ending in "?" rests on a single
    kind of object and may be wrong -- a sink alone reads as a kitchen until a toilet
    shows). size= floor area. frontier= unexplored openings or boundaries still
    accessible inside it (0 is NOT proof that nothing is there). status= one of:
    never_entered (only seen from outside, or not at all); entered (the robot stood in
    it and its camera has seen most of what is mapped of it, but it never completed a
    look-around: what is left of it lies behind its frontier= boundaries);
    scanned(<how>, <ago>) (the robot completed a full look-around in it, or saw most of
    it from one, <ago> ago -- what was visible then was seen ONCE by a detector that
    misses objects that are small, far, half-hidden, behind a door or under furniture).
    searched= time spent inside. seen= objects confirmed in it. here=yes: the robot
    stands in it now.
OPENING: a doorway or gap at the edge of the mapped floor leading to space NOT seen yet
    -- a room the robot has not entered. Going there is a quick look from the threshold,
    cheaper than a room visit. Value it by the room likely BEHIND it: glimpsed objects
    name that room (a toilet -> a bathroom; a bed -> a bedroom); nothing glimpsed -> an
    unidentified room, judged by what the home still lacks. An opening whose glimpsed
    objects name the home type of the target (a sink or a shower for a toilet; a sofa or
    a television for a couch) is the home room seen from its door: the best node there
    is, 60 or more.
The HOUSE line is the big picture: every room type found so far (scanned ones included),
    how many rooms are still unidentified, how many openings are not yet looked into,
    whether any unexplored frontier is still reachable, and the budget -- actions used
    of the total. Use ALL of it.

HOW TO JUDGE. Work in this order and write steps 1-3 into the reply before any number.
STEP 1 "home": the room types where this object normally lives, most likely first, and
    where it is rare. Typical homes: a toilet in a bathroom -- a separate one, or an
    en-suite reached through a bedroom; a bed in a bedroom; a sofa or couch in a living
    room (sometimes a family room or an office); a television in a living room, then a
    bedroom, rarely a kitchen, never a bathroom; a chair in a dining room, kitchen,
    office, bedroom or living room; a potted plant in a living room, hallway, kitchen or
    balcony; a refrigerator, oven or stove in a kitchen; a bathtub or shower in a
    bathroom. Small objects that travel (cup, book, bottle) have a home too, only a
    weaker one.
STEP 2 "house": read the home from the rooms found. A typical home has ONE kitchen, ONE
    living room, ONE dining area, one to four bedrooms, one to three bathrooms (one
    often en-suite off the main bedroom), a hallway, perhaps an office, a laundry or a
    garage. If a type that occurs once has been FOUND, an unidentified room is unlikely
    to be another of it, and the odds of the target's home being among the unidentified
    rooms drop. If the home type is still MISSING, the unidentified rooms and the
    openings are where it must be: value them in proportion to how well their size fits
    the home type (bathrooms ARE small, 3-8 m2; bedrooms 9-20 m2; living rooms 15 m2 and
    up) and to what else the home still lacks -- never write a room off for its size. A
    home-type object seen in a room of ANOTHER type (a sink in a "living room", a bed in
    a "kitchen") means the map merged two rooms: value that room as the home type.
STEP 3 "stage": weigh the budget. EARLY -- much budget left and unexplored places
    remain: the target is most likely somewhere not yet looked at; unexplored rooms and
    openings first, scanned rooms low. LATE -- most of the home seen, few or no
    unexplored places or frontier left, the home type found but the target not seen: the
    detector may have MISSED the target in a scanned room of a fitting type (a toilet
    behind a half-open door, a television on a wall the scan grazed, a chair under a
    table), and the chance that an unseen room holds it falls with every place looked
    at; raise the scanned rooms of the home type -- a second look from a different spot
    -- and lower the rest. Write "pass": "first" while an unexplored place in this home
    is still worth looking at before any scanned room is re-checked, else "second". The
    first pass over the home comes before the second, unless a scanned home-type room is
    the only candidate left.
STEP 4 the numbers, one per node, independently, integers 0-100:
 - a scanned room of a type the target does not live in: 0-2.
 - a scanned room of the home type with the target not seen: 3-10 early in the budget,
    15-40 late in it, less if it was scanned moments ago.
 - an entered but unscanned room: the robot has walked it and seen most of it without
    seeing the target; value only what its frontier= boundaries still hide -- at most
    HALF of what an unexplored room of the same type and size would get, and under 10
    with frontier=0. One searched= 20s or longer with frontier= 0 or 1 is as good as
    scanned: 0-5, whatever its type. Never write "known to contain" the target of a room
    the target was not seen in.
 - a never-entered room of unknown type, or an opening with nothing glimpsed: never
    below 10 while the home type is missing; by size fit and by what the home lacks once
    it is found.
 - a never-entered room named as the home type by objects seen through its door, or an
    opening whose glimpse names the home type: 60 or more.
 - a never-entered room of the home type with frontier left, target not seen: 30 or
    more; the same room entered: 10-20, more only if its frontier leads to a large
    unseen part of it.
 - prefer the places the robot has NOT been near: an opening or room away from
    everything entered and scanned leads to a part of the home not seen at all; a gap at
    the edge of a scanned room usually shows a strip of it.
 - a hallway or corridor rarely holds the target, but one with frontier LEADS to unseen
    rooms: value it by what it may lead to.
 - searched long and recently -> low, not zero; searched briefly or long ago -> stays
    promising.
 - SPREAD, do not copy: the unexplored places compete for the same missing rooms. A
    missing bathroom is ONE room, so five unlooked doorways share its chance -- about
    10-15 each, the best-fitting one more, never 35 on every one of them; with many
    openings and one or two missing home-type rooms, most openings get under 15. Several
    openings often lead to the same hallway. The sum over the unexplored places should
    be about the chance that the target is somewhere not yet looked at, not more.
ALWAYS: NEVER use distance, travel time or here=yes as a reason -- the planner charges
    travel, and the room the robot stands in is judged exactly like the others. Several
    rooms may each hold an instance: the numbers are independent search-success
    probabilities, NOT shares of one distribution; do NOT make them sum to 100. Score
    EVERY node id given, and only those ids. Per node write "why" first (at most 12
    words, lower case, naming the room's likely type and what decided it), then "p".
Reply with ONLY this JSON, keys in this order:
{"home":"<step 1, max 12 words>","house":"<step 2, max 25 words>","stage":"<step 3, max
    20 words>",
 "pass":"first|second","nodes":[{"id":<int>,"why":"<why>","p":<int>}, ...]}
EXAMPLE of the FORMAT and the style of reasoning -- a different home every time, so
    never copy its numbers or its words. TARGET toilet; HOUSE: rooms found: kitchen
    (scanned), living_room (scanned), bedroom (never entered), hallway (entered); 1 room
    unidentified; 1 opening not yet looked into; actions used about 150 of 500:
id=2  ROOM  type=kitchen  size=12m2  frontier=0  status=scanned(full rotation, 2min
    ago)  searched=30s  seen: oven, refrigerator, sink
id=3  ROOM  type=living_room  size=24m2  frontier=1  status=scanned(full rotation, 1min
    ago)  searched=40s  seen: sofa, television
id=5  ROOM  type=bedroom  size=14m2  frontier=1  status=never_entered  searched=0s
     seen: bed
id=6  ROOM  type=unknown  size=5m2  frontier=1  status=never_entered  searched=0s  seen:
    nothing yet
id=7  ROOM  type=hallway  size=6m2  frontier=1  status=entered  searched=10s  here=yes
     seen: nothing yet
id=200001  OPENING  doorway off room 7 (type=hallway)  to space NOT seen yet  glimpsed
    through it: nothing yet
{"home":"bathroom; often an en-suite off a bedroom",
 "house":"kitchen and living room found and scanned; no bathroom yet; two unidentified
    rooms, one bathroom-sized; one unlooked doorway",
 "stage":"early, 150 of 500 used, unexplored places remain: first pass",
 "pass":"first",
 "nodes":[{"id":2,"why":"kitchen, scanned, toilets are not in kitchens","p":1},
{"id":3,"why":"living room, scanned, no toilet there","p":1},
{"id":5,"why":"bedroom never entered, an en-suite may open off it","p":25},
{"id":6,"why":"small unidentified room, bathroom sized, bathroom missing","p":55},
{"id":7,"why":"hallway, leads to the unlooked doorway","p":5},
{"id":200001,"why":"unseen room off the hallway, bathroom still missing","p":35}]}
Had the HOUSE line read "every known room identified; 0 openings not yet looked into; no
    unexplored frontier left; actions used about 420 of 500" with a bathroom found and
    scanned 4 minutes ago, the pass would be "second": the scanned bathroom 35, the
    scanned bedroom 8 (an en-suite door the scan may have missed), and the scanned
    kitchen and living room 1 each.
\end{verbatim}
\end{small}

\noindent Multi-storey supplement (appended to the system prompt only when a staircase node is offered or another storey is known):
\begin{small}
\begin{verbatim}
THIS SEARCH MAY CHANGE STOREYS. STAIRS nodes are staircases to another storey (up or
    down). Taking one means leaving this storey and searching the other one. The line
    says whether that storey was visited before and, if so, what was found and how much
    was searched there, and whether the robot arrived on this storey by these stairs a
    moment ago (arrived_by=yes). The THIS STOREY line gives this storey's rank in the
    building (the lowest of N known storeys, the highest, or how many lie below and
    above). Read the storey from the rooms found on it first -- a storey with a kitchen,
    dining or living room is the ground floor whatever its rank -- and from the rank
    where none has: of two storeys the lowest is the ground floor and the highest the
    bedroom floor; of three, the lowest is usually a basement, the middle the ground
    floor and the highest the bedroom floor. Where rooms live in a house: kitchen,
    dining and living room downstairs; bedrooms upstairs; the bathroom beside the
    bedrooms, so upstairs too (a ground floor has at most a small toilet room); an
    office on either. Add to the reply, before "nodes", the key "home_here": "found" (a
    room of the home type is on this storey), "missing" (none found yet, but the home
    type belongs on this storey) or "elsewhere" (the home type lives on another storey).
    STAIRS to an unvisited storey where the home type belongs can be high without
    suppressing the other nodes; stairs to a visited storey where the home type was
    FOUND and barely searched can be highly promising; arrived_by=yes: do not send the
    robot straight back unless it left that storey almost unsearched. An unexplored
    place on this storey is worth about 10 when home_here is "elsewhere" -- a look into
    it comes after the storey the target lives on. Changing storeys costs many actions,
    but that is NOT your concern: the planner adds travel and climbing cost itself.
\end{verbatim}
\end{small}

\noindent User prompt (\texttt{\{storeys\}} is empty in single-storey mode; otherwise it holds the two storey lines shown next, followed by a line break):
\begin{small}
\begin{verbatim}
TARGET: {target}

HOUSE: {house}
{storeys}
NODES ({n}):
{nodes}

Score all {n} nodes independently. JSON only.
\end{verbatim}
\end{small}

\noindent Storey lines (multi-storey mode only):
\begin{small}
\begin{verbatim}
THIS STOREY: {storey}
OTHER STOREYS: {others}
\end{verbatim}
\end{small}

\noindent Schema-repair text (appended after a line break to the user prompt, once, after an unusable reply). The break after \texttt{"first|second",} is the only line break in this appendix that does not stand for a space; the original has none there:
\begin{small}
\begin{verbatim}
Your previous answer did not satisfy the schema. Return exactly
    {"home":"...","house":"...","stage":"...","pass":"first|second",
    "nodes":[{"id":0,"why":"brief reason","p":50}, ...]} with one entry per actual node
    id from the input (not the example id), independent integer p values in 0-100 (do
    not normalise across nodes), and nothing outside the JSON object.
\end{verbatim}
\end{small}

% ----------------------------------------------------------------------------
\section{Benchmark Protocols}\label{app:protocols}

This section records how each benchmark is set up and scored in our evaluation, at the level of detail needed to reproduce the numbers or to judge their comparability with published results. \Cref{tab:benchmarks} summarises the five splits, and \cref{tab:privileges} lists the inputs our agent receives beyond each benchmark's nominal observation.

\subsection{Shared evaluation harness}\label{app:protocols-harness}

All five splits are scored by a single evaluation harness that wraps each simulator behind a common episode interface. The policy receives the RGB-D frame, the agent pose and camera pitch, and the goal category name, together with the privileged inputs disclosed in \cref{tab:privileges}. Geodesic distances, shortest paths, goal positions and view points stay inside the evaluator; as an additional safeguard, the harness refuses at reset any episode metadata whose keys would carry such quantities. Matching between the detector's class names and the benchmark's goal categories is exact set membership after normalisation, never fuzzy matching, because on a benchmark a false accept becomes a false STOP.

\paragraph{Metrics.} Metrics follow \cref{eq:spl,eq:softspl,eq:dtg}; we add the edge cases, which follow one convention throughout: the ratio $\ell/\max(\ell,L)$ is taken as 1 when $\ell = L = 0$ (habitat-lab would divide by zero, and the RoboTHOR evaluators also take 1) and is 0 when $\ell = 0 < L$; the soft progress is 1 when $d_0 = d_{\mathrm{end}} = 0$, 0 when $d_0 = 0 < d_{\mathrm{end}}$, and 0 when $d_{\mathrm{end}} = \infty$. Run-level figures are means over all episodes, failures included.

\paragraph{Path length.} Unless stated otherwise, $L_i$ is the sum of the 3-D chord lengths between successive base positions, starting from the reset position; turns, camera tilts and STOP add nothing, and a blocked forward step adds only the distance actually moved.

\paragraph{Protocol checks.} Before an episode is scored, the harness refuses a success without STOP on every benchmark that requires STOP, requires its own path length to agree with the simulator's within $10^{-3}$\,m, and requires every metric that the simulator also computes natively to agree with its own within $10^{-4}$. After every action it checks that the motion was the one commanded: a turn translates the base by less than 2\,cm and changes the heading by the nominal angle within $3^\circ$; a camera tilt changes the pitch by the nominal angle within $0.5^\circ$, or not at all at a pitch limit; a forward step moves at most the step length plus 3\,cm in the plane, within $10^\circ$ of the heading, and climbs no more than the larger of 0.2\,m and the planar step; STOP moves nothing.

\paragraph{Depth on the Habitat benchmarks.} Readings closer than the minimum range are marked invalid and readings at or beyond the maximum range are marked as returning nothing, whereas habitat-lab clamps both into the valid interval. This changes what the policy perceives, not how it is scored. The filling of enclosed depth holes is a further deviation, disclosed in \cref{tab:privileges}.

\subsection{Detector vocabulary}\label{app:protocols-vocab}

The detector is queried once per frame with a fixed list of text prompts. On Gibson the list has 30 prompts, in this order: the eight goal prompts \emph{chair, couch, sofa, potted plant, bed, toilet, tv, television}; the room-context objects \emph{dining table, oven, sink, refrigerator, book, clock, vase, cup, bottle}; the four door prompts \emph{door, doorway, open doorway, door frame}; the further context objects \emph{cabinet, desk, shower, bathtub}; the two stair prompts \emph{stairs, staircase}; and the three distractors \emph{toy, rocking horse, stuffed animal}. A goal category is accepted under its own prompt and its aliases only (couch under \emph{couch} or \emph{sofa}, tv under \emph{tv} or \emph{television}). Stair boxes never become object landmarks. A distractor class is never a goal and names no room; it exists only to give a confusable object a better-matching prompt than a goal category. \TODO{list the detector vocabularies (goal prompts with their aliases, context, door, stair and distractor prompts) used on HM3D, MP3D and RoboTHOR}

\subsection{Gibson (SemExp v1.1, tiny validation split)}\label{app:protocols-gibson}

\paragraph{Episodes and scenes.} We use the ObjectNav episodes released with SemExp~\citep{chaplot2020object} (version 1.1) on the Gibson scenes~\citep{xia2018gibson}: the validation partition of the tiny split, 1{,}000 episodes, 200 in each of five buildings (Collierville, Corozal, Darden, Markleeville and Wiconisco). Every published start and goal lies on one storey. The six goal categories are chair, couch, potted plant, bed, toilet and tv, with 177, 229, 193, 123, 208 and 70 validation episodes respectively.

\paragraph{Sensors, agent and actions.} The agent carries a $640 \times 480$ RGB-D camera with a $79^\circ$ horizontal field of view at a height of 0.88\,m and a depth range of 0.5--5.0\,m; the agent radius is 0.18\,m and its height equals the camera height. There are six actions, STOP, MOVE\_FORWARD (0.25\,m), TURN\_LEFT/TURN\_RIGHT ($30^\circ$) and LOOK\_UP/LOOK\_DOWN ($30^\circ$, with the camera starting level and kept between $30^\circ$ above and $60^\circ$ below level), and sliding along obstacles is enabled, unlike on HM3D and MP3D. The two tilt actions are our deviation from the SemExp action set, which has only the first four; a tilt uses one action of the budget and moves nothing, so the scoring is unchanged. An episode ends on STOP or after 500 actions.

\paragraph{Success.} Following the SemExp code, an episode succeeds if the geodesic distance from the agent's final position to the goal region is zero, where the goal region is the target category's ground-truth semantic map dilated by 1\,m. STOP is \emph{not} required, so an episode can succeed at the step limit, and there is no visibility test. Our evaluator is an independent reimplementation of SemExp's: on a 0.05\,m grid it dilates the traversable map by 0.10\,m, dilates the goal cells by 1\,m without intersecting them with the traversable area, and computes geodesic distances with the fast marching method~\citep{sethian1996fast}; masked or unreachable cells receive a finite sentinel equal to the largest valid distance plus one cell. Positions are converted to grid cells by truncation, as in SemExp.

\paragraph{SPL and DTG conventions.} The reference length is the fast-marching distance from the start to the goal region plus the 1\,m success radius, $\ell = d_0 + 1$\,m, as in SemExp. The travelled length $L$ accumulates the \emph{horizontal} displacement only, starting from $10^{-5}$\,m; the native \SPL, $\min(S\,\ell/L, 1)$, equals the harness formula. \DTG is the final fast-marching distance to the dilated success region, that is, a distance to the success boundary rather than to the object, and it includes the sentinel value: a final position on a masked cell or outside the floor map reads the sentinel, and that value enters the mean. At least one published start lies on a masked cell, where the reference length itself carries the upstream sentinel; we keep such episodes, as SemExp does, and record them. \TODO{report the number of sentinel starts and of sentinel end positions in our runs}

\paragraph{Simulator.} We run habitat-sim 0.2.4~\citep{savva2019habitat}, whereas the SemExp reference used 0.1.5. The published Gibson navigation mesh is loaded unchanged (\cref{app:protocols-navmesh}). Some published starts lie up to 0.28\,m below the navigation mesh, and one lies 0.033\,m outside it horizontally; starts are never snapped, and the reset check tolerates these offsets. SemExp issues a dummy turn to the simulator when the agent calls STOP; we omit it, which affects no score.

\subsection{HM3D v1 and v2}\label{app:protocols-hm3d}

\paragraph{Episodes and scenes.} HM3D v1 is the validation split of the 2022 Habitat ObjectNav challenge over HM3D-Semantics v0.1~\citep{ramakrishnan2021habitat,yadav2023habitat}: 2{,}000 episodes in 20 scenes, with six goal categories (chair, bed, plant, toilet, tv\_monitor, sofa) that are not balanced (only 84 episodes target a plant). HM3D v2 uses the v2 episodes over HM3D-Semantics v0.2: 1{,}000 validation episodes in 36 scenes, with the same six categories and 135 to 195 episodes each. As in the papers we compare with, the v2 episodes are run with the v1 agent described below; the 2023 challenge itself moved to a different robot with velocity control. The 20 v1 validation scenes are a subset of the 36 v2 scenes, but HM3D-Semantics v0.2 removed debris from four validation scenes, one of which is in the v1 set. The published start rotations are stored in $(x,y,z,w)$ order and are converted to the simulator's $(w,x,y,z)$ order.

\paragraph{Sensors, agent and actions.} The camera is a $640 \times 480$ RGB-D sensor with a $79^\circ$ horizontal field of view at 0.88\,m and a depth range of 0.5--5.0\,m; the agent is 0.88\,m tall with a radius of 0.18\,m, and sliding is disabled. There are six actions: STOP, MOVE\_FORWARD (0.25\,m), TURN\_LEFT/TURN\_RIGHT ($30^\circ$) and LOOK\_UP/LOOK\_DOWN ($30^\circ$). Episodes end on STOP or after 500 actions, STOP included. HM3D v1 contains episodes whose goal lies on another floor; according to VLFM~\citep{yokoyama2024vlfm}, about 14.6\% of HM3D episodes require climbing stairs. \TODO{state whether stair traversal is enabled on HM3D, and how many of our failures need a floor change}

\paragraph{Success, SPL and DTG.} Success requires STOP \emph{and} a geodesic distance below 0.1\,m from the agent to the nearest goal view point; the dataset generator placed the view points within about 1\,m of each goal object. Distances follow habitat-lab's distance-to-goal measure computed to view points: one multi-goal shortest-path query over every view point of every goal instance, recomputed only when the agent has moved. The reference length $\ell$ is this distance at reset, not the geodesic distance stored with the episode, and $L$ is the sum of 3-D chords from a zero initial value. The stored geodesic distance is itself measured to the nearest view point, which is why it is shorter than the straight-line distance to the nearest goal object in 731 of the 2{,}000 v1 episodes. We do not run habitat-lab itself: its distance-to-goal, success, \SPL and \SoftSPL measures are reimplemented on habitat-sim 0.2.4, with habitat-lab 0.2.1 as the reference implementation. An episode with a non-finite reference length cannot be scored (habitat-lab's \SPL would be undefined) and is refused before the run. \TODO{state whether any HM3D validation episode is affected}

\paragraph{Second success radius.} Because some baselines are scored with a 0.2\,m success radius (\cref{app:protocols-caveats}), we additionally re-score the same logged episodes at 0.2\,m. The radius changes only the verdict on the agent's own STOP position, so both scores come from a single run and cannot drift apart. \TODO{state whether the MP3D episodes are also re-scored at 0.2\,m}

\subsection{MP3D}\label{app:protocols-mp3d}

\paragraph{Episodes and scenes.} We use the validation split of the Habitat 2021 ObjectNav challenge on Matterport3D~\citep{chang2017matterport3d}: 2{,}195 episodes in 11 scenes with 21 goal categories (chair, table, picture, cabinet, cushion, sofa, bed, chest\_of\_drawers, plant, sink, toilet, stool, towel, tv\_monitor, shower, bathtub, counter, fireplace, gym\_equipment, seating, clothes).

\paragraph{Sensors, agent, actions and success.} Camera, depth range, agent dimensions, sliding, the six actions and the 500-action limit are as on HM3D. Success requires STOP and a geodesic distance below 0.1\,m to the nearest goal view point. The reference length is the distance to goal at reset ($\ell = d_0$), and $L$ is the sum of 3-D chords from the reset position. The evaluator is an independent reimplementation of habitat-lab's navigation measures (habitat-lab 0.2.2 and 0.2.4 task configurations); the simulator's semantic sensor is never instantiated, so no ground-truth semantics can reach the agent. Some MP3D episodes require changing floors (about 9.6\% according to VLFM~\citep{yokoyama2024vlfm}). \TODO{state whether stair traversal is enabled on MP3D}

\subsection{RoboTHOR}\label{app:protocols-robothor}

\paragraph{Episodes and scenes.} We use the validation split of the 2021 RoboTHOR ObjectNav challenge~\citep{deitke2020robothor} in AI2-THOR~\citep{kolve2017ai2}: 1{,}800 episodes, 120 in each of 15 apartments (three apartment layouts in five configurations each), with 12 goal categories of 150 episodes each (AlarmClock, Apple, BaseballBat, BasketBall, Bowl, GarbageCan, HousePlant, Laptop, Mug, SprayBottle, Television, Vase). The challenge web page states a different episode count; the released files contain 1{,}800.

\paragraph{Sensors, agent and actions.} The camera produces $640 \times 480$ frames; AI2-THOR's field-of-view parameter (63.45$^\circ$) is vertical and corresponds to a $79^\circ$ horizontal field of view. The camera sits 0.8688\,m above the floor (the 0.9\,m agent origin minus the LoCoBot's 0.0312\,m camera offset, verified in the simulator). The policy uses depth in 0.1--5.0\,m; the renderer clips at 0.1--20\,m, and its far value is read as no return. The agent is the LoCoBot capsule (radius 0.175\,m, height 0.9\,m) with the challenge's stochastic actuation: Gaussian rotation noise with a standard deviation of $0.5^\circ$ and Gaussian step noise with a mean of 0.001\,m and a standard deviation of 0.005\,m. The six actions are MoveAhead (0.25\,m), RotateLeft/RotateRight ($30^\circ$), LookUp/LookDown ($30^\circ$, with the camera horizon clamped to $\pm 30^\circ$) and Stop; every episode starts with the camera fully lowered ($+30^\circ$). Episodes end on Stop or after 500 actions, and a failed action still consumes a step.

\paragraph{Success.} An episode succeeds if the agent calls Stop while a goal instance is \emph{visible}, where AI2-THOR's visibility test includes a 1\,m visibility distance; no geodesic distance enters the success test. The official evaluator checks only the first-listed instance of the goal category; our adapter accepts any visible instance and records both verdicts. \TODO{state which verdict the reported RoboTHOR numbers use, and how often the two differ}

\paragraph{SPL and DTG conventions.} The reference length $\ell$ is the shortest-path length shipped with each episode, precomputed by the challenge organisers with AI2-THOR's shortest-path query; it is not recomputed. In 108 of the 1{,}800 episodes (6.0\%) the shipped length is zero. Under the official arithmetic, which we keep, such an episode scores \SPL 0 whenever the agent moves at all ($\ell = 0 < L$), even if it succeeds, and 1 only for a successful Stop without moving; an agent that moves in every episode can therefore reach a mean \SPL of at most 94\%. The travelled length $L$ is the sum of 3-D chords from the spawn position, as in AI2-THOR's metrics utilities. \DTG is the geodesic distance to the nearest goal instance, computed with AI2-THOR's shortest-path query and the progressively relaxed endpoint tolerances used by AllenAct; unreachable end positions are counted separately. No published RoboTHOR \DTG or \SoftSPL exists. The evaluator is transcribed from the official challenge code and AI2-THOR's metrics module, with the Unity build pinned and AI2-THOR 2.7.2 as the reference client.

\paragraph{Deviations.} The official challenge runner supplies only the goal, RGB and depth, with no pose; our runs use the simulator's ground-truth pose (\cref{tab:privileges}). Depth is optional in the official runner, and we use it. Success is additionally recorded under the any-instance visibility rule described above.

\subsection{Navigation-mesh recompute}\label{app:protocols-navmesh}

The HM3D and MP3D scene meshes ship with navigation meshes built for an agent of radius 0.10\,m and height 1.50\,m. habitat-sim silently recomputes the navigation mesh when the agent's dimensions differ from the stored ones, and habitat-lab copies the ObjectNav agent's dimensions (0.18\,m, 0.88\,m) into the simulator configuration, so every published HM3D and MP3D distance, collision test, reference length, \SR and \SPL comes from the recomputed mesh. We therefore evaluate HM3D v1, HM3D v2 and MP3D on the mesh recomputed for the 0.18\,m, 0.88\,m agent. With it, our HM3D reference lengths reproduce the geodesic distances stored with the episodes exactly, whereas the shipped mesh gives distances 5.5\% too short; on one HM3D scene the ObjectNav agent settings reduce the navigable area by 15\%.

The recompute matters. On the Gibson building Collierville, where we measured it, the recomputed mesh has a navigable area of 44.145\,m$^2$ against 58.006\,m$^2$ for the shipped mesh, 16 of 25 sampled geodesic distances differ by more than 5\,cm (by up to 0.6\,m), and 14 of 39 sampled start--goal pairs are connected only on the shipped mesh. An evaluator that loaded the shipped mesh without the recompute would therefore move \SR and \SPL. On Gibson we load the published navigation mesh unchanged, as stated in \cref{app:protocols-gibson}. The Gibson score does not depend on the mesh, because success and \DTG are computed on the ground-truth semantic floor map, but motion and collisions during the episode do. \TODO{state whether SemExp's habitat-sim 0.1.5 recomputed the Gibson navigation mesh for the 0.18\,m agent}

\subsection{Comparability caveats}\label{app:protocols-caveats}

The published numbers we compare with were not all obtained under the same protocol. We list every difference we found that can move a number, so that each comparison in \cref{sec:results} can be read with its caveats. Unless stated otherwise, baseline figures below are SR/SPL in percent, as printed in the cited source.

\begin{enumerate}[leftmargin=*,itemsep=2pt]
\item \textbf{Success radius on HM3D and MP3D.} The official rule is STOP within 0.1\,m of a goal view point. ApexNav~\citep{zhang2025apexnav} states a 0.2\,m success distance, its released configurations use 0.2\,m on HM3D v1, HM3D v2 and MP3D, and its re-evaluations of other methods on HM3D v2 inherit that setting. The released SG-Nav~\citep{yin2024sg} MP3D entry point loads a configuration with a 0.2\,m success distance, although the task configuration bundled with its code says 0.1\,m. UniGoal~\citep{yin2025unigoal} uses 1.6\,m on HM3D, the released code of InstructNav~\citep{long2024instructnav} 0.25\,m, and OSG-Nav~\citep{loo2025open} 1.0\,m together with a visibility condition. A looser radius can only raise \SR and \SPL. We mark such rows and report our own HM3D runs at both 0.1\,m and 0.2\,m (\cref{app:protocols-hm3d}).
\item \textbf{Sensor settings.} ApexNav opens the depth range to 0.0--5.0\,m. SG-Nav's paper states a 0.90\,m camera and a 1.5--10\,m perception range, while its configurations mount the camera at 0.88\,m and set the maximum depth to 10\,m on HM3D. We keep the standard 0.88\,m camera and 0.5--5.0\,m depth range on the Habitat benchmarks.
\item \textbf{Episode subsets.} OSG-Nav's HM3D results use 400 episodes that the authors sampled from the HM3D v0.1 validation split (20 per scene), with a 100-episode subset for its language-model ablation. The appendix failure count of CogNav~\citep{cao2025cognav} implies about 1{,}008 HM3D episodes rather than 2{,}000, and it describes its MP3D set as 1.8K episodes with 20 categories rather than 2{,}195 episodes with 21. InstructNav's code samples a subset of episodes (500 by default). Such rows are not comparable to full-split numbers and are marked or kept apart.
\item \textbf{Test versus validation.} The fine-tuned ProcTHOR~\citep{deitke2022procthor} HM3D figure that is commonly reprinted beside validation numbers (54.4/31.8) comes from the Habitat 2022 leaderboard, that is, the test-standard split. On RoboTHOR, the ProcTHOR (65.2/28.8 after fine-tuning) and ProcTHOR-ZS (55.0/23.7) figures are \emph{test-split} numbers (2{,}040 episodes), yet they appear next to validation numbers in the tables of ESC~\citep{zhou2023esc} and SG-Nav. The EmbCLIP~\citep{khandelwal2022simple} figure of 47.0/20.0 is likewise its test-split result; its validation result on the challenge leaderboard, on the same 1{,}800 episodes we use, is 52.2/26.0. PIRLNav~\citep{ramrakhya2023pirlnav} reports 70.4/34.1 on HM3D validation and 65.0/33.0 on test-standard. We never place a test-split number in a validation table without a mark.
\item \textbf{Quoted versus measured, and drift in reprinted numbers.} ApexNav's statement that it re-evaluated baselines covers HM3D v2 only; its HM3D v1 baseline rows are quotations. SG-Nav's tables print figures that differ from the methods' own papers: on HM3D v1, L3MVN~\citep{yu2023l3mvn} at 48.7/23.0 against its own 50.4/23.1, OpenFMNav~\citep{kuang2024openfmnav} at 52.5/24.1 against 54.9/24.4, and VLFM~\citep{yokoyama2024vlfm} at 52.4/30.3 against 52.5/30.4; on MP3D, VLFM at 36.2/15.9 against 36.4/17.5. Later papers inherit these figures. The MP3D figures of L3MVN and OpenFMNav and the RoboTHOR figures of L3MVN, VLFM and OpenFMNav exist only in SG-Nav's table. ESC's HM3D figures differ across its arXiv versions (44.0/25.2, 38.5/22.0 and, in the conference version we quote, 39.2/22.3). SemExp's MP3D figure (36.0/14.4) was measured on its own 2{,}000-episode, six-category MP3D set rather than the 2{,}195-episode validation split, yet it is commonly reprinted as a validation number; a third-party run on the validation split, as reprinted by SGM~\citep{zhang2024imagine}, obtained 28.3/10.9. Wherever a method's own paper reports the benchmark, we quote that paper.
\item \textbf{Reproducibility.} ESC's RoboTHOR result (38.1/22.2) was not reproduced: a later paper that re-ran ESC with its official implementation obtained 34.5/18.2~\citep{unlu2024reproducibility}. SG-Nav's released code covers MP3D only, so its HM3D and RoboTHOR numbers cannot be reproduced from the release; the 0.1\,m HM3D radius sometimes attributed to it comes from a configuration file that its released entry point does not load.
\item \textbf{Gibson distance and success definitions.} Gibson papers mix two distance metrics under similar names: some report the distance to the goal object, others the distance to the success boundary, $\max(d - 1\,\text{m}, 0)$ (often abbreviated DTS), sometimes in a column labelled DTG. Our \DTG is the fast-marching distance to the 1\,m-dilated goal region, as in the SemExp code, and is therefore of the second kind; distances should be compared only between rows that share a definition. OSG-Nav's written success rule adds that the goal must be directly observable, which the SemExp code does not check; we follow the code. SemExp's paper reports an \SR of 0.544 and an \SPL of 0.199 on its own 1{,}000-episode Gibson evaluation, which predates the v1.1 episode file, whereas the documentation of its released model reports 0.657 and 0.339 on v1.1 (both printed as fractions), so the two are distinct evaluations and must not be interchanged. The Gibson ObjectNav result of ZSON~\citep{majumdar2022zson} uses a different episode set and sensor configuration (RGB only at $128 \times 128$, a $90^\circ$ field of view, a 1.25\,m camera, no depth and no GPS) and is not comparable with numbers on the SemExp split. We use habitat-sim 0.2.4 rather than the reference 0.1.5, and no Gibson paper publishes episode hashes or an evaluator revision, so no row is bit-reproducible. The current public release of PONI~\citep{ramakrishnan2022poni} uses a modified, multi-goal evaluator and episode format; our loader rejects such episodes.
\item \textbf{Navigation-mesh recompute.} The recompute of \cref{app:protocols-navmesh} happens inside habitat-sim, so it affects every Habitat-based HM3D and MP3D result in the same way and is not a difference between, for example, ApexNav and SG-Nav. It does separate any evaluation that loads the shipped meshes without the recompute.
\item \textbf{Training regime.} We call a method training-free only if it trains no navigation policy of any kind. Some training-free methods still use perception trained in-domain: L3MVN and TriHelper~\citep{zhang2024trihelper} use a semantic segmentation network fine-tuned in-domain, and PONI's frontier-exploration baseline uses a detector fine-tuned on Gibson training images. Several zero-shot methods train a navigation module on another task: VLFM (PointNav), PixNav~\citep{cai2024bridging} (pixel-goal navigation), and ZSON and PSL~\citep{sun2024prioritized} (ImageNav); ZSON also selected its checkpoint on ObjectNav validation success. SGM~\citep{zhang2024imagine}, which some tables list as zero-shot, trains a self-supervised map-generation model for navigation on in-domain training scenes and is therefore not training-free. The ``training-free'' and ``non-metric'' columns in OSG-Nav's Gibson table are method attributes, not metrics; \ours is training-free but builds a metric map.
\item \textbf{Inputs and action spaces.} Most Habitat baselines use RGB-D with GPS and compass, which matches our use of ground-truth pose; ZSON, PSL and PixNav use RGB only, UniGoal and CoW act without look-up and look-down actions, and WMNav~\citep{nie2025wmnav} acts through at most 40 polar moves proposed by a vision-language model rather than 0.25\,m and $30^\circ$ primitives. On Gibson we add look-up and look-down actions to the four actions of the SemExp protocol under which the Gibson baselines are evaluated (\cref{app:protocols-gibson}). On RoboTHOR our ground-truth pose exceeds the official protocol, which provides none (ESC runs without it). Our further privileges, noise-free depth with hole filling, visibility-gated ground-truth stair geometry and the projection of the closing standoff onto the navigation mesh, are listed in \cref{tab:privileges}.
\end{enumerate}

% ----------------------------------------------------------------------------
\section{Additional Results}\label{app:additional}

\subsection{Per-category results}\label{app:additional-category}

\Cref{tab:percategory} breaks \SR and \SPL down by goal category on every benchmark. The categories differ in physical size, in where they are typically placed, and in how easily the target can be kept whole in the camera frame at close range, and their frequencies differ between splits (only 70 of the 1{,}000 Gibson episodes target a tv, and only 84 of the 2{,}000 HM3D v1 episodes a plant), so the benchmark means weight them very differently. Fifteen of the 208 Gibson toilet episodes were used to tune the later revisions of the method (\cref{sec:experiments-privileges}), so the Gibson toilet row is the entry most exposed to that tuning. \TODO{discuss per-category trends once the runs are in, in particular for low targets such as toilets on Gibson}

\begin{table}[tp]
\centering
\small
\setlength{\tabcolsep}{4pt}
\caption{Per-category results of \ours (\SR and \SPL in \%). ``Eps'' is the number of validation episodes of the category. (a)~The six-category benchmarks; category names are Gibson's, with the HM3D name in parentheses where it differs. (b)~MP3D. (c)~RoboTHOR.}
\label{tab:percategory}
\textit{(a) Gibson, HM3D v1 and HM3D v2}\\[2pt]
\begin{tabular}{l rrr rrr rrr}
\toprule
& \multicolumn{3}{c}{Gibson} & \multicolumn{3}{c}{HM3D v1} & \multicolumn{3}{c}{HM3D v2} \\
\cmidrule(lr){2-4}\cmidrule(lr){5-7}\cmidrule(lr){8-10}
Category & Eps & \SR & \SPL & Eps & \SR & \SPL & Eps & \SR & \SPL \\
\midrule
chair & 177 & \tbd & \tbd & \tbd & \tbd & \tbd & \tbd & \tbd & \tbd \\
couch (sofa) & 229 & \tbd & \tbd & \tbd & \tbd & \tbd & \tbd & \tbd & \tbd \\
potted plant (plant) & 193 & \tbd & \tbd & 84 & \tbd & \tbd & \tbd & \tbd & \tbd \\
bed & 123 & \tbd & \tbd & \tbd & \tbd & \tbd & \tbd & \tbd & \tbd \\
toilet & 208 & \tbd & \tbd & \tbd & \tbd & \tbd & \tbd & \tbd & \tbd \\
tv (tv\_monitor) & 70 & \tbd & \tbd & \tbd & \tbd & \tbd & \tbd & \tbd & \tbd \\
\midrule
All & 1000 & \tbd & \tbd & 2000 & \tbd & \tbd & 1000 & \tbd & \tbd \\
\bottomrule
\end{tabular}

\medskip
\textit{(b) MP3D}\\[2pt]
\begin{tabular}{l rrr l rrr}
\toprule
Category & Eps & \SR & \SPL & Category & Eps & \SR & \SPL \\
\midrule
chair & \tbd & \tbd & \tbd & stool & \tbd & \tbd & \tbd \\
table & \tbd & \tbd & \tbd & towel & \tbd & \tbd & \tbd \\
picture & \tbd & \tbd & \tbd & tv\_monitor & \tbd & \tbd & \tbd \\
cabinet & \tbd & \tbd & \tbd & shower & \tbd & \tbd & \tbd \\
cushion & \tbd & \tbd & \tbd & bathtub & \tbd & \tbd & \tbd \\
sofa & \tbd & \tbd & \tbd & counter & \tbd & \tbd & \tbd \\
bed & \tbd & \tbd & \tbd & fireplace & \tbd & \tbd & \tbd \\
chest\_of\_drawers & \tbd & \tbd & \tbd & gym\_equipment & \tbd & \tbd & \tbd \\
plant & \tbd & \tbd & \tbd & seating & \tbd & \tbd & \tbd \\
sink & \tbd & \tbd & \tbd & clothes & \tbd & \tbd & \tbd \\
toilet & \tbd & \tbd & \tbd & \textit{All} & 2195 & \tbd & \tbd \\
\bottomrule
\end{tabular}

\medskip
\textit{(c) RoboTHOR}\\[2pt]
\begin{tabular}{l rrr l rrr}
\toprule
Category & Eps & \SR & \SPL & Category & Eps & \SR & \SPL \\
\midrule
AlarmClock & 150 & \tbd & \tbd & HousePlant & 150 & \tbd & \tbd \\
Apple & 150 & \tbd & \tbd & Laptop & 150 & \tbd & \tbd \\
BaseballBat & 150 & \tbd & \tbd & Mug & 150 & \tbd & \tbd \\
BasketBall & 150 & \tbd & \tbd & SprayBottle & 150 & \tbd & \tbd \\
Bowl & 150 & \tbd & \tbd & Television & 150 & \tbd & \tbd \\
GarbageCan & 150 & \tbd & \tbd & Vase & 150 & \tbd & \tbd \\
\midrule
\multicolumn{4}{l}{} & \textit{All} & 1800 & \tbd & \tbd \\
\bottomrule
\end{tabular}
\end{table}

\subsection{Failure analysis}\label{app:additional-failures}

\Cref{tab:failures} assigns every failed episode to one of six mutually exclusive categories and reports each category's share of the failures on each benchmark. \TODO{state the assignment procedure (rules over the episode logs, manual review of a sample, or both), its precedence order when several categories apply, and the number of failed episodes per benchmark} The categories, and the mechanisms behind them that we observed during development, are as follows.

\begin{table}[tp]
\centering
\small
\caption{Failure taxonomy of \ours: share of the failed episodes (\%) in each category, per benchmark. Categories are mutually exclusive, and each column sums to 100\%.}
\label{tab:failures}
\begin{tabular}{l ccccc}
\toprule
Failure category & Gibson & HM3D v1 & HM3D v2 & MP3D & RoboTHOR \\
\midrule
False-positive STOP & \tbd & \tbd & \tbd & \tbd & \tbd \\
Missed target / timeout & \tbd & \tbd & \tbd & \tbd & \tbd \\
Closing failure & \tbd & \tbd & \tbd & \tbd & \tbd \\
Navigation stuck & \tbd & \tbd & \tbd & \tbd & \tbd \\
Annotation gap & \tbd & \tbd & \tbd & \tbd & \tbd \\
Other & \tbd & \tbd & \tbd & \tbd & \tbd \\
\midrule
Failed episodes (count) & \tbd & \tbd & \tbd & \tbd & \tbd \\
\bottomrule
\end{tabular}
\end{table}

\paragraph{False-positive STOP.} The agent ends the episode at an object that is not an instance of the goal category (on Gibson, which does not require STOP, the agent calls STOP outside the success region at such an object). These are detector confusions that survive every check of the closing routine: the two consistent frames of verification, the 0.50 confidence threshold and the near-edge range rule. The confusions we observed are between objects of similar shape or function seen at close range, such as a kitchen counter read as a bed, bar seating read as a sofa, and a child's ride-on toy horse read as a chair; the distractor prompts of the detector vocabulary (\cref{app:protocols-vocab}) exist to absorb confusions of the last kind. Because STOP is issued on detector evidence alone, the closing checks reduce these errors but cannot remove them. A STOP at a genuine but unannotated instance is counted as an annotation gap, not here.

\paragraph{Missed target / timeout.} The action budget runs out without the target ever being verified (on Gibson, without the agent ending inside the success region). Typical mechanisms are a target present in a room that was already scanned but missed by the detector because it was small, distant, partly occluded or only grazed by the scan; large floors on which the budget is spent visiting many rooms and peeking through many openings before the right one; and an oracle ranking that sends the agent elsewhere first, for instance because a room was misnamed from a single object glimpsed through a door. The optional second pass over finished rooms (\cref{sec:planning}), which targets the first mechanism late in the episode, is off by default.

\paragraph{Closing failure.} The target was detected and the closing routine took over, but the episode did not end in success. Without camera tilt, as on Gibson before we added the tilt actions (\cref{app:protocols-gibson}), a mechanism we observed repeatedly is a low target such as a toilet whose box is cut by the lower edge of the frame at close range, so that the candidate is released even though the agent stands next to it; the tilt-free rules of \cref{sec:execution} target this case. Other mechanisms are repeated release-and-retake cycles on a wrong support surface that consume the budget, before the revision described in \cref{sec:experiments}, a reachable standoff goal discarded for an unreachable one when the target's estimated position was refined, and a distant misdetection that is locked before the real target. On the benchmarks with a strict success test, this category also contains STOPs issued next to a true instance but more than 0.1\,m from every view point (HM3D, MP3D) or with no goal instance visible (RoboTHOR).

\paragraph{Navigation stuck.} The agent stops making progress for reasons of geometry rather than search: it is wedged against geometry that the 2-D occupancy map does not show, collision slides keep resetting the blocked-step detector, a split-level floor falls outside the single height band so that neither floor nor obstacle is observed, or the obstacle mask placed around a seen staircase traps the agent. Leaving the start storey by walking down a staircase that had not been seen also falls here.

\paragraph{Annotation gap.} The agent stops at, or ends next to, a real instance of the goal category that the benchmark does not count: the instance is missing from the semantic annotation, is annotated under another category, lies on a part of the scene the annotation does not cover, or (on HM3D and MP3D) has no view point from which the agent's position would count. Under the official RoboTHOR rule, which checks only the first-listed instance, a STOP at another visible instance of the category is a further case. This category can only be assigned by inspecting the episode. \TODO{describe how annotation gaps were confirmed}

\paragraph{Other.} Every remaining failure: episodes ended by an internal consistency check of the agent or the harness, failures of the detector or language-model services that leave the agent in the exploration fallback until the budget is exhausted, and simulator irregularities such as published starts that lie off the navigation mesh. \TODO{list the dominant causes in this category once the runs are in}

\end{document}
